\documentclass[titlepage,12pt,a4paper,twoside,openright,svgnames]{report}%
%\addtolength\oddsidemargin {10cm}\addtolength\evensidemargin {-10cm}

%PAQUETES PABLO
\usepackage{dirtree}
\usepackage{movie15}
\usepackage{indentfirst}
\usepackage{enumitem}
\usepackage{afterpage}
\usepackage{soul}
\usepackage{siunitx}
\usepackage{xcolor}
\usepackage{xspace}
\usepackage{mhchem}
\usepackage{amsmath}
\usepackage{wrapfig}
\usepackage{dirtree}
\usepackage{changepage}
\usepackage{geometry}



%\geometry{bindingoffset=3mm}
\def\videoautorefname{Video}%


\usepackage{tabulary}
\usepackage{tabularx}
\usepackage{booktabs}
\newcommand{\heading}[1]{\multicolumn{1}{c}{#1}}
\usepackage{caption}

%%%%%%%%%code listings


\definecolor{codegreen}{rgb}{0,0.6,0}
\definecolor{codegray}{rgb}{0.5,0.5,0.5}
\definecolor{codepurple}{rgb}{0.58,0,0.82}
\definecolor{backcolour}{rgb}{0.95,0.95,0.92}

%%Code listing style named "mystyle"
%\lstdefinestyle{mystyle}{
  %%backgroundcolor=\color{backcolour},
  %commentstyle=\color{codegreen},
  %keywordstyle=\color{magenta},
  %numberstyle=\tiny\color{codegray},
  %stringstyle=\color{codepurple},
  %basicstyle=\ttfamily\footnotesize,
  %breakatwhitespace=false,
  %breaklines=true,
  %captionpos=b,
  %keepspaces=true,
  %numbers=left,
  %numbersep=5pt,
  %showspaces=false,
  %showstringspaces=false,
  %showtabs=false,
  %tabsize=2
  %stepnumber=5
  %numberfirstline=false
%}


%"mystyle" code listing set
%\lstset{style=mystyle}


%%%%%%%%%%%%%%

%\usepackage{estilos/classicthesis} % Custom classicthesis


\def\chapterautorefname{Chapter}%
\def\figureautorefname{Figure}%
\def\sectionautorefname{Section}%
\def\subsectionautorefname{Subsection}%
\def\tableautorefname{Table}%
\def\equationautorefname{Equation}%

%para que no se corten parrafos a mitad
%\widowpenalties 1 10000
%\raggedbottom

%Video floats
\usepackage{float}
\newfloat{videoFloat}{thp}{lop}[chapter]
\floatname{videoFloat}{Video}

%Pablo custom code colors
\usepackage{color}
\definecolor{atomString}{RGB}{221,17,68}
\definecolor{atomComment}{RGB}{153,153,136}
\definecolor{black}{RGB}{0,0,0}
\definecolor{mygreen}{RGB}{28,172,0} % color values Red, Green, Blue
\definecolor{mylilas}{RGB}{170,55,241}
\usepackage{upquote}
\usepackage{listings}
% CSS
\lstdefinelanguage{CSS}{
  keywords={color,background-image:,margin,padding,font,weight,display,position,top,left,right,bottom,list,style,border,size,white,space,min,width, transition:, transform:, transition-property, transition-duration, transition-timing-function},
  sensitive=true,
  morecomment=[l]{//},
  morecomment=[s]{/*}{*/},
  morestring=[b]',
  morestring=[b]",
  alsoletter={:},
  alsodigit={-}
}

% JavaScript
\lstdefinelanguage{JavaScript}{
  morekeywords={typeof, delete,new, true, false, catch, function, return, null, catch, switch, var, if, in, while, do, else, case, break},
  morecomment=[s]{/*}{*/},
  morecomment=[l]//,
  morestring=[b]",
  morestring=[b]'
}

\lstdefinelanguage{HTML5}{
  language=html,
  sensitive=true,
  alsoletter={<>=-},
  morecomment=[s]{<!-}{-->},
  tag=[s],
  otherkeywords={
  % General
  >,
  % Standard tags
	<!DOCTYPE,
  </html, <html, <head, <title, </title, <style, </style, <link, </head, <meta, />,
	% body
	</body, <body,
	% Divs
	</div, <div, </div>,
	% Paragraphs
	</p, <p, </p>,
	% scripts
	</script, <script,
  % More tags...
  <canvas, /canvas>, <svg, <rect, <animateTransform, </rect>, </svg>, <video, <source, <iframe, </iframe>, </video>, <image, </image>, <header, </header, <article, </article
  },
  ndkeywords={
  % General
  =,
  % HTML attributes
  charset=, src=, id=, width=, height=, style=, type=, rel=, href=,
  % SVG attributes
  fill=, attributeName=, begin=, dur=, from=, to=, poster=, controls=, x=, y=, repeatCount=, xlink:href=,
  % properties
  margin:, padding:, background-image:, border:, top:, left:, position:, width:, height:, margin-top:, margin-bottom:, font-size:, line-height:,
	% CSS3 properties
  transform:, -moz-transform:, -webkit-transform:,
  animation:, -webkit-animation:,
  transition:,  transition-duration:, transition-property:, transition-timing-function:,
  }
}

\lstdefinestyle{htmlcssjs} {%
  % General design
	backgroundcolor=\color{white},
  basicstyle={\footnotesize\ttfamily},
  frame=b,
  % line-numbers
  xleftmargin={0.75cm},
  numbers=left,
  stepnumber=1,
  firstnumber=1,
  numberfirstline=true,
  % Code design
  identifierstyle=\color{black},
  keywordstyle=\color{blue}\bfseries,
  ndkeywordstyle=\color{black}\bfseries,
  stringstyle=\color{red}\ttfamily,
  commentstyle=\color{atomComment}\ttfamily,
  % Code
  language=HTML5,
  alsolanguage=JavaScript,
  alsodigit={.:;},
  tabsize=1,
  showtabs=false,
  showspaces=false,
  showstringspaces=false,
  extendedchars=true,
  breaklines=true,
	 captionpos=b,                    % sets the caption-position to bottom
literate=%
*{0}{{{\color{red}0}}}1
{1}{{{\color{red}1}}}1
{2}{{{\color{red}2}}}1
{3}{{{\color{red}3}}}1
{4}{{{\color{red}4}}}1
{5}{{{\color{red}5}}}1
{6}{{{\color{red}6}}}1
{7}{{{\color{red}7}}}1
{8}{{{\color{red}8}}}1
{9}{{{\color{red}9}}}1,
 % German umlauts
  literate=%
  {Ö}{{\"O}}1
  {Ä}{{\"A}}1
  {Ü}{{\"U}}1
  {ß}{{\ss}}1
  {ü}{{\"u}}1
  {ä}{{\"a}}1
  {ö}{{\"o}}1
}
%
\lstdefinestyle{py} {%
language=python,
backgroundcolor=\color{white},
literate=%
*{0}{{{\color{red}0}}}1
{1}{{{\color{red}1}}}1
{2}{{{\color{red}2}}}1
{3}{{{\color{red}3}}}1
{4}{{{\color{red}4}}}1
{5}{{{\color{red}5}}}1
{6}{{{\color{red}6}}}1
{7}{{{\color{red}7}}}1
{8}{{{\color{red}8}}}1
{9}{{{\color{red}9}}}1,
basicstyle=\footnotesize\ttfamily, % Standardschrift
numbers=left,               % Ort der Zeilennummern
%numberstyle=\tiny,          % Stil der Zeilennummern
%stepnumber=2,               % Abstand zwischen den Zeilennummern
numbersep=5pt,              % Abstand der Nummern zum Text
tabsize=4,                  % Groesse von Tabs
extendedchars=true,         %
breaklines=true,            % Zeilen werden Umgebrochen
keywordstyle=\color{blue}\bfseries,
frame=b,
commentstyle=\color{atomComment}\itshape,
stringstyle=\color{red}\ttfamily, % Farbe der String
showspaces=false,           % Leerzeichen anzeigen ?
showtabs=false,             % Tabs anzeigen ?
xleftmargin=17pt,
framexleftmargin=17pt,
framexrightmargin=5pt,
framexbottommargin=4pt,
%backgroundcolor=\color{lightgray},
showstringspaces=false,      % Leerzeichen in Strings anzeigen ?
}


\usepackage{estilos/pfc}
\usepackage{estilos/funciones}
\PassOptionsToPackage{spanish, english}{babel}
\usepackage[spanish, english]{babel}
\usepackage{lscape}
\usepackage{graphicx}
\usepackage{epstopdf}
\usepackage{rotating}
\usepackage{hyperref}
\usepackage[noabbrev]{cleveref}      % reference object types automatically

\usepackage{fixltx2e}
\usepackage{longtable}

\usepackage{listings}
\usepackage{color}
\usepackage{eurosym}
\usepackage{pdfpages}

\usepackage{algorithm}      % http://ctan.org/pkg/algorithms
\usepackage{algpseudocode}  % http://ctan.org/pkg/algorithmicx
\usepackage{multirow}
\usepackage{appendix}
\usepackage{subfig}         % subfiguras

%\usepackage[3D]{movie15}
%\usepackage{media9}
\usepackage{hyperref}

\usepackage{url}
\usepackage{enumerate}
\usepackage{inputenc}
%\usepackage{multimedia}


%\decimalpoint
\makeindex %Para hacer el índice alfabético

%============ ELADIO GUTIERREZ CARRASCO  =====================
% -*- Demosle unas leccioncillas a LaTeX sobre como colocar figuras -*-
\setcounter{topnumber}{3}     % Max. numero de figs. on top
\setcounter{bottomnumber}{3}  % Max. numero de figs. abajo
\setcounter{totalnumber}{10}   % Max. numero de figs. por pagina
\renewcommand{\topfraction}{1} % Max. fraccion de pagina ocupada por figs.
\renewcommand{\bottomfraction}{1}
\renewcommand{\textfraction}{0}  % Min. fraccion de pagina ocupada por texto
\renewcommand{\floatpagefraction}{0.5} % Max. espacio de pagina solo con figs.
\textfloatsep=0.5mm



%SQUEEZE:
% http://www.eng.cam.ac.uk/help/tpl/textprocessing/squeeze.html

%Page Layout
        %\columnsep: gap between columns
        %\topmargin: gap above header
        %\topskip: between header and text
        %\textheight: height of main text
        %\textwidth: width of text
        %\oddsidemargin: odd page left margin
        %\evensidemargin : even page left margin

%Paragraphs
%\parindent=1cm %: indentation of paragraphs
\parskip=3mm %: gap between paragraphs

%OJOOOOOOOOOOOO MARGENES PARA GUILLOTINAR AL IMPRIMIR, QUITAR PARA VERSIÓN DIGITAL. SOLO HAY QUE PONER EL EVEN EN PRINCICIO
%\addtolength\evensidemargin{3mm}
\addtolength\oddsidemargin{-1.8mm}
%%%%%%%%%%%%%%% ESTO TB PARA IMPRESIÓN,  FIJA EL TAMAÑO EN A4 +3 MM EN CADA LADO
 %\geometry{
		%papersize={216mm,303mm},
 %}
%%%%%%%%%%%%%%%%%%%%%%%%%%%%%%%%%%%

%\oddsidemargin=10mm

%Floats (tables and figures) Parece que no funciona cuando se usa [H]

\floatsep=1mm %: is the space between adjacent floats that are moved to the top or bottom of the text page.
\textfloatsep=1mm %: is the space between the main text and floats at the top or bottom of the page.
\intextsep=5mm %: space left on top and bottom of an in-text float.
%\dbltextfloatsep is \textfloatsep for 2 column output.
%\dblfloatsep is \floatsep for 2 column output.
\abovecaptionskip=0mm %: space above caption
\belowcaptionskip=-2mm %: space below caption

%Maths
\abovedisplayskip=2mm %: space before maths
\belowdisplayskip=2mm %: space after maths
        %\arraycolsep: gap between columns of an array

%Lists
\topsep=1mm %: space between first item and preceding paragraph.
\partopsep=0mm %: extra space added to \topsep when environment starts a new paragraph.
\itemsep=-1.5mm %-2mm %: space between successive items
%\parskip=0mm % 	Space between paragraphs outside of a list, and part of the space between a non-list paragraph and a list item.
%\topsep 	Extra space added to \parskip before the first and after the last item.
%\parsep=0mm %	Paragraph separation within a single item.
%\itemsep 	Extra inter-item spacing added to \parsep.
\partopsep=0mm % 	This is added to the top and/or bottom of the list if and only if there's a blank line above or below the first or last item. Leave this alone unless blank lines become a problem.

%
%By default, LaTeX doesn't like to fill more than 0.7 of a text page with tables and graphics, nor does it like too many figures per page. This behaviour can be changed by placing lines like the following before \begin{document}
%
%\renewcommand\floatpagefraction{.9}
%\renewcommand\topfraction{.9}
%\renewcommand\bottomfraction{.9}
%\renewcommand\textfraction{.1}
%\setcounter{totalnumber}{50}
%\setcounter{topnumber}{50}
%\setcounter{bottomnumber}{50}


%PARA Ahorrar espacio
%http://ftp.eq.uc.pt/software/TeX/macros/latex/contrib/savetrees/savetrees.pdf
%\usepackage{savetrees}

%%% NO HYPENATION %%%
\tolerance=1
\emergencystretch=\maxdimen
\hyphenpenalty=10000
\hbadness=10000
%%% NO HYPENATION %%%

\captionsetup{belowskip=10pt,aboveskip=2pt} %%Caption space between figure and below text

\includeonly{subdocs/portada}

%%Aboveskip space between content and caption
%%Belowskip space between content and surrounding text
\usepackage{glossaries}
\usepackage[acronym]{glossaries}
\makeglossaries

\begin{document}
%  may have noticed the absence of underline. This functionality has to be
% added with the ulem package. Stick \usepackage{ulem} in your preamble.
% By default, this overrides the \emph command with the underline rather than
% the italic style. It is unlikely that you wish this to be the desired effect,
% so it is better to stop ulem taking over \emph and simply call the underline
% command as and when it is needed.
%
% To disable ulem, add \normalem straight after the document environment begins.
% To underline, use \uline{...}.
% To add a wavy underline, use \uwave{...}.
% And for a strike-out \sout{...}.
\normalem
% Para ver si en Paragraph empieza en linea nueva
% \setlength{\parindent}{20pt}
% \parskip = \baselineskip
\selectlanguage{english}

\hypersetup{linkcolor=black, citecolor=black}        % color of links to bibliograp}
\renewcommand{\contentsname}{Index}
\renewcommand{\partname}{Part}
\renewcommand{\chaptername}{Chapter}
\renewcommand{\appendixname}{Appendix}
\renewcommand{\bibname}{References}
\renewcommand{\figurename}{Figure}
\renewcommand{\tablename}{Table}
\renewcommand{\listfigurename}{List of Figures}
\renewcommand{\acronymname}{Acronyms}
\renewcommand{\glossaryname}{Glossary}


% Esto funciona normalmente pero en el caso de tener cargado el parquete CAPTION
% hay que hacerlo de otra manera
\renewcommand{\tablename}{Table}
\renewcommand{\listtablename}{List of Tables}


\renewcommand*\lstlistingname{Code}
\renewcommand*\lstlistlistingname{Code Index}

\newcommand\blankpage{%
    \null
    \thispagestyle{empty}%
    \addtocounter{page}{-1}%
    \newpage}


%\glossarystyle{long4col}
\pagenumbering{alph}

%\part (book y report) -1
%\part (article) 0
%\chapter 0
%\section 1
%\subsection 2
%\subsubsection 3
%\paragraph 4
%\subparagraph 5


% Para definir el espacio antes y después de cada fórmula que aparentemente es excesivo
% \abovedisplayskip=12pt plus 3pt minus 9pt
% \belowdisplayskip=12pt plus 3pt minus 9pt
% \abovedisplayskip=3pt
% \belowdisplayskip=3pt
% \abovedisplayshortskip=0pt plus 3pt
% \belowdisplayshortskip=7pt plus 3pt minus 4pt


%def\TikZ{Ti\emph{k}Z}
%lstset{frameround=fttt}
%lstloadlanguages{TeX}


% Portada

\newlength{\originalVOffset}
\newlength{\originalHOffset}
\setlength{\originalVOffset}{\voffset}
\setlength{\originalHOffset}{\hoffset}

\setlength{\voffset}{0cm}
\setlength{\hoffset}{0cm}

\includepdf[pages=-,link=true,landscape]{./Portada/portada_tfm_pablo.pdf}

\setlength{\voffset}{\originalVOffset}
\setlength{\hoffset}{\originalHOffset}

\addcontentsline*{toc}{chapter}{\numberline{}{}}

\pagenumbering{roman} \setcounter{page}{1}

%--
\thispagestyle{empty}
%\includepdf[pages=-,link=true,landscape,linkname=PortadaPastas]{StarTrackerPortada.pdf}

%\cleardoublepage NO METER
%\thispagestyle{empty} NO METER
%\clearpage{\pagestyle{empty}\cleardoublepage} NO METER
%--
%\newpage

\thispagestyle{empty}
\vspace*{3cm}
%\begin{flushright}
%
%%\textbf{\large ``Design of a multidisciplinary 1U CubeSat Simulation Platform''}
%\vfill
%\parindent=0pt
%\begin{small}
%\begin{flushleft}
%Credits for the cover: \textbf{NASA}.\\
%Printed in Granada, September 2019.
%\\~\\
%All rights reserved.
%\end{flushleft}
%
%\end{small}
%
%\afterpage{\blankpage}
%
%\par\end{flushright}{\large \par}

\newpage


%\thispagestyle{empty}
%\vspace*{3cm}
%\begin{flushright}
%
%%\textbf{\large ``Design of a multidisciplinary 1U CubeSat Simulation Platform''}
%
%\textbf{\large ``Design of a multidisciplinary}\\
%\textbf{\large  1U CubeSat Simulation Platform''}
%
%\afterpage{\blankpage}
%
%\par\end{flushright}{\large \par}
%\clearpage

%\thispagestyle{empty}

\begin{center}
\textbf{\huge \includegraphics[scale=1.45]{FigurasTFM/logo_ugr.pdf}}
\par\end{center}{\huge \par}

\begin{center}
\vspace*{1cm}
\par\end{center}

\begin{center}
\textbf{\large GRADO EN}\\
\textbf{\large INGENIERÍA INFORMÁTICA}
\par\end{center}{\large \par}

\begin{center}
\textbf{\large Trabajo de fin de grado}
\par\end{center}{\large \par}

\begin{center}

\par\end{center}

\begin{center}
\textbf{\emph{\LARGE {}``Limitador de sonido}}\\
\textbf{\emph{\LARGE {} para locales de música''}}

\par\end{center}{\LARGE \par}

\begin{center}
\vspace*{3cm}
\par\end{center}

\begin{center}
{\large CURSO ACADÉMICO: 2020 - 2021}
\par\end{center}{\large \par}

\begin{center}
{\large Alejandro Ruiz Becerra}
\par\end{center}{\large \par}

\newpage
\thispagestyle{empty}

~

\newpage
\thispagestyle{empty}

\begin{center}
\includegraphics[scale=1.45]{FigurasTFM/logo_ugr.pdf}
\par\end{center}

\begin{center}
GRADO EN INGENIERÍA INFORMÁTICA\par\end{center}

\begin{center}
\vspace*{0.1cm}
\par\end{center}

\begin{center}
\textbf{\emph{\LARGE {}``Limitador de sonido}}\\
\textbf{\emph{\LARGE {} para locales de música''}}
\par\end{center}{\Large \par}

\begin{center}
\vspace*{0.3cm}
\par\end{center}

\begin{center}
REALIZADO POR:
\par\end{center}

\begin{center}
\textbf{Alejandro Ruiz Becerra}
\par\end{center}

\begin{center}
DIRIGIDO POR:
\par\end{center}

\begin{center}
\textbf{Andrés María Roldán Aranda}
\par\end{center}

\begin{center}
DEPARTAMENTO:
\par\end{center}

\begin{center}
\textbf{Electrónica y Tecnología de Computadores}
\par\end{center}

\begin{center}
\vfill
\par\end{center}

\vspace*{1.5cm}

\newpage
\thispagestyle{empty}
\noindent
\blankpage

%Begin ----  Para que funcione bien el TOC en PDF
\clearpage
\thispagestyle{empty}
\phantomsection
\addcontentsline{toc}{chapter}{Autorización Lectura}

\noindent D. Andrés María Roldán Aranda, Profesor del departamento
de Electrónica y Tecnología de los Computadores de la Universidad
de Granada, como director del Trabajo Fin de Grado de D. Alejandro Ruiz Becerra,

\vspace*{1cm}

Informa:

\begin{doublespace}
Que el presente trabajo, titulado:
\end{doublespace}

\begin{doublespace}
\begin{center}
\textbf{\emph{\large {}``Limitador de sonido para locales de música''}}
\par\end{center}{\large \par}
\end{doublespace}

\noindent ha sido realizado y redactado por el mencionado alumno bajo
mi dirección, y con esta fecha autorizo a su presentación.

\vspace*{1cm}

\begin{center}
Granada, a 21 de Julio de 2021
\par\end{center}

\bigskip
\bigskip
\begin{center}
\includegraphics[scale=0.2]{FigurasTFM/firmaAndres.png}
\end{center}

\begin{center}
\begin{doublespace}
Fdo. Andrés María Roldán Aranda
\end{doublespace}
\end{center}

\newpage
\thispagestyle{empty}
\noindent

\newpage
\phantomsection
\noindent
\blankpage

\addcontentsline{toc}{chapter}{Autorización Depósito Biblioteca}
\bigskip

\noindent Los abajo firmantes autorizan a que la presente copia de
Trabajo Fin de Grado se ubique en la Biblioteca del Centro y/o
departamento para ser libremente consultada por las personas que lo
deseen.

\vspace*{1cm}

\begin{center}
Granada, a 21 de Julio de 2021
\par\end{center}

\bigskip
\bigskip

\begin{center}
\hspace{0cm}\includegraphics[scale=0.9]{FigurasTFM/firmaJC.png}\hspace{3cm} \includegraphics[scale=0.2]{FigurasTFM/firmaAndres.png}
\end{center}

\begin{doublespace}
\begin{center}
\hspace{0cm}Fdo. Alejandro Ruiz Becerra \hspace{3cm} Fdo. Andrés María Roldán Aranda
\end{center}
\end{doublespace}
~

%Begin ----  Para que funcione bien el TOC en PDF
\clearpage
\phantomsection
\noindent
\thispagestyle{empty}


\addcontentsline{toc}{chapter}{Resumen}
\vspace{-1.48cm}
\begin{center}
%\begin{adjustwidth}{-9pt}{0pt}
    \textbf{\Large Limitador de sonido para locales de música}
%\end{adjustwidth}
\par\end{center}{\Large \par}

\begin{center}
    \textbf{\large Alejandro Ruiz Becerra}
    \par\end{center}{\large \par}

\vspace{0.75cm}


\begin{doublespace}
    \noindent \textbf{PALABRAS CLAVE:}
\end{doublespace}


\begin{singlespace}
    \noindent GranaSAT, Acústica y audio, Ingeniería Acústica, Ingeniería Inversa, Control de ruidos, Ecualización, Electrónica.

%    \glsname{cubesat}, \glsname{altium}, Diseño aeroespacial, \glsname{solid}, \glsname{ground}, \acrshort{EDA}, Electrónica, \acrshort{OBC}, \acrshort{ADCS}, \acrshort{EPS}, \acrshort{OBDH}, Diseño de \acrshort{PCB}, \acrshort{matlab}.

\end{singlespace}

\begin{doublespace}
    \noindent \textbf{RESUMEN:}
\end{doublespace}

\begin{singlespace}

    \noindent El objetivo del presente proyecto es diseñar e implementar el software necesario para la construcción de un limitador de sonido para locales de ocio, de forma que se cumplan las especificaciones legales y ordenanzas exigidas por las instituciones en éste ámbito, siendo beneficiara del presente trabajo la empresa \textbf{Heimdal Sound Control}.

    \noindent El proyecto puede dividirse en tres grandes bloques: ingeniería inversa, diseño e implementación. Durante la primera fase se estudian y analizan limitadores de sonido de la competencia, ya presentes en el mercado; para luego diseñar el sistema en base a los requisitos extraídos del proceso de ingeniería inversa, y finalmente desarrollar y probar el software del producto.

    \noindent Este Trabajo de Fin de Grado se sitúa en el ámbito de un proyecto mayor, ambicioso y de largo recorrido, y se apoya en el trabajo realizado por otros alumnos pertenecientes a diversas competencias. Por tanto, el presente trabajo no debe verse como un todo, sino como un gran engranaje dentro de una máquina mayor, el cuál permite que el conjunto de componentes interaccionen entre ellos.

    \noindent La complejidad y el ámbito multidisciplinar de este Trabajo de Fin de Grado permite cubrir, no sólo algunas de las diferentes especialidades del Grado en Ingeniería Informática, sino también adquirir conocimientos y habilidades transversales o específicos de otros campos de la Ingeniería, como la \textbf{Electrónica} y la \textbf{Acústica}.

    \noindent El resultado de todo lo expuesto culmina con un equipo real de limitación de sonido completo y funcional, que cumple con los requisitos definidos en las etapas iniciales del proyecto, y con el cual se cierra la etapa universitaria de Grado.

%    \noindent El objetivo principal del presente proyecto es desarrollar una Plataforma de Simulación multidisciplinar de \textbf{CubeSats}. Estará compuesta de tres bloques diferenciados, en torno a los cuales pivotará el proyecto: una plataforma de simulación mecánica, un software de gestión de \glsname{ground} y un prototipo de \glsname{cubesat}, que constituirá la base del futuro \textbf{GranaSAT-I}.
%
%    Este Trabajo Fin de Máster se aborda desde una ambiciosa doble perspectiva: por un lado, el desarrollo de una Plataforma de Simulación de amplia utilidad en el ámbito académico, como medio para el acercamiento del alumnado de múltiples titulaciones al mundo aeroespacial y en concreto a los CubeSats, en el contexto de auge actual, fomentado por instituciones como la \textbf{Agencia Espacial Europea }(\acrshort{ESA}); en segundo lugar, en el ámbito de investigación, proveyendo de un medio para la implementación de nuevos algoritmos de comunicación, de control orbital y, en general, para el desarrollo y testeo de tecnologías y técnicas novedosas, de manera previa a su lanzamiento.
%
%    El desarrollo e implementación de este proyecto se lleva a cabo siguiendo metodologías de \textbf{Ingeniería de Sistemas} contrastadas y asentadas en la industria espacial, dotándolo de realismo y acercando al alumno a técnicas profesionales de amplio reconocimiento en el mercado de trabajo. Asimismo, la complejidad y ámbito multidisciplinar de este Trabajo~Fin~de~Máster le permite cubrir, no sólo las diferentes especialidades del Máster~de~Ingeniería~de~\textbf{Telecomunicación}, sino también adquirir conocimientos y habilidades transversales o específicos de otros campos de la Ingeniería, como la \textbf{Mecánica} o la \textbf{Aeroespacial}. Así, además de software especialista de cada uno de los campos mencionados, se han analizado y aplicado técnicas avanzadas de \textbf{mecanizado} (fresado de aluminio mediante control numérico), \textbf{fabricación} (soldadura utilizando técnicas de \textit{reflow}) o \textbf{caracterización} de diferentes dispositivos (baterías de litio, células solares de silicio...), entre otros.
%
%    El resultado de todo lo expuesto culmina con la obtención de un entorno de simulación completo y funcional, que cumple con los requisitos definidos en etapas iniciales, y con el cual se cierra la etapa universitaria de Máster.

\end{singlespace}

\vspace{1.25cm}

\newpage

~

\vspace{-1.15cm}


\begin{otherlanguage}{english}

\begin{center}
%\begin{adjustwidth}{-9pt}{0pt}
\textbf{\Large Sound limiter for music venues}
%\end{adjustwidth}
\par\end{center}{\Large \par}

\begin{center}
\textbf{\large Alejandro Ruiz Becerra}
\par\end{center}{\large \par}

\vspace{0.75cm}
\begin{doublespace}
\noindent \textbf{KEYWORDS:}
\end{doublespace}

\begin{singlespace}

    \noindent GranaSAT, Acoustics and audio, Acoustical engineering, Reverse engineering, Noise control, Equalization, Electronics.

%\noindent \glsname{cubesat}, \glsname{altium}, Aerospace design, \glsname{solid}, \glsname{ground}, \acrshort{EDA}, Electronic, \acrshort{OBC}, \acrshort{ADCS}, \acrshort{EPS}, \acrshort{OBDH}, \acrshort{PCB} Design, \acrshort{matlab}.

\end{singlespace}

\begin{doublespace}
\noindent \textbf{ABSTRACT:}
\end{doublespace}

\begin{singlespace}

%    \noindent The complexity and the mutildisciplinary scope of this Bachelor's Thesis allows to cover, not only some of the different specialties of the Bachelor Degree in Informatics Engineering, but also to adquire knowledge and transversal habilities from other fields of the Engineering, such as \textbf{Electronics} and \textbf{Acoustics}.
%
    \noindent El objetivo del presente proyecto es diseñar e implementar el software necesario para la construcción de un limitador de sonido para locales de ocio, de forma que se cumplan las especificaciones legales y ordenanzas exigidas por las instituciones en éste ámbito, siendo beneficiara del presente trabajo la empresa \textbf{Heimdal Sound Control}.

    \noindent El proyecto puede dividirse en tres grandes bloques: ingeniería inversa, diseño e implementación. Durante la primera fase se estudian y analizan limitadores de sonido de la competencia, ya presentes en el mercado; para luego diseñar el sistema en base a los requisitos extraídos del proceso de ingeniería inversa, y finalmente desarrollar y probar el software del producto.

    \noindent Este Trabajo de Fin de Grado se sitúa en el ámbito de un proyecto mayor, ambicioso y de largo recorrido, y se apoya en el trabajo realizado por otros alumnos pertenecientes a diversas competencias. Por tanto, el presente trabajo no debe verse como un todo, sino como un gran engranaje dentro de una máquina mayor, el cuál permite que el conjunto de componentes interaccionen entre ellos.

    \noindent La complejidad y el ámbito multidisciplinar de este Trabajo de Fin de Grado permite cubrir, no sólo algunas de las diferentes especialidades del Grado en Ingeniería Informática, sino también adquirir conocimientos y habilidades transversales o específicos de otros campos de la Ingeniería, como la \textbf{Electrónica} y la \textbf{Acústica}.

    \noindent El resultado de todo lo expuesto culmina con un equipo real de limitación de sonido completo y funcional, que cumple con los requisitos definidos en las etapas iniciales del proyecto, y con el cual se cierra la etapa universitaria de Grado.


%\noindent The main purpose of this project is developing a multidisciplinary Simulation Platform for \textbf{CubeSats}. It will be composed of three differentiated blocks, around which the project is structured: a mechanical simulation platform, a \glsname{ground} management software and a \glsname{cubesat} prototype that will be the base for the future \textbf{GranaSAT-I}.
%
%This Master's Thesis is addressed from a double perspective: on the one hand, the development of a Simulation Platform of great usefulness in an academic environment, as a way to get students from multiple degrees closer to the aerospace world, and particularly to CubeSats, given its current context of peak, being fostered by institutions such as \textbf{European Space Agency} (\acrshort{ESA}); on the other hand, in a research environment, providing with a mean to implement new communication algorithms, orbit controllers, and generally speaking, for the development and test of new technologies and techniques, before launching.
%
%The development and implementation of this project is performed following methodologies of System Engineering contrasted in the aerospace industry, giving realism and getting the student closer to professional techniques, widely recognized in the job market. Furthermore, the complexity and multidisciplinary scope of this Master's~Thesis allows covering not only the different specialties of the Master~in~\textbf{Telecommunication}~Engineering but also acquiring knowledge and transversal abilities from other fields of the Engineering, such as \textbf{Mechanical} or \textbf{Aerospace}. Besides specific software of each of the mentioned areas, advanced techniques of \textbf{machining} (aluminum milling), \textbf{manufacturing} (solder reflow) or \textbf{characterization} of different devices (lithium batteries, silicon solar cells...) among others, have been analyzed and applied.
%
%The result of the exposed culminates with the obtention of a complete and functional simulation environment, which  complies with the requirements defined in the preliminary stages, and supposes the finalization of the Master.
\end{singlespace}

\newpage
\thispagestyle{empty}
%

\end{otherlanguage}

~

%Begin ----  Para que funcione bien el TOC en PDF
\clearpage
\phantomsection
\thispagestyle{empty}
\addcontentsline{toc}{chapter}{Dedication}

\vspace*{8cm}

\begin{quotation}
\noindent \begin{flushright}
%\textbf{\emph{\Large Dedicado a}}\textbf{\emph{\large }}\\
\textbf{\emph{\large \lq Tough and competent\rq}}\\
%\textbf{\emph{  Eugene F. Kranz}}\\
%\textbf{\emph{\large Answer}}\\
%\textbf{\emph{\large \lq --Whatever it takes\rq}}\\
%\textbf{\emph{\large Mis padres, Paco y Encarni, y mi hermano Alberto, porque sin ellos y sin su apoyo, llegar hasta aquí hubiera sido imposible.}}
%\textbf{\emph{\large .....Esto lo último.....}}
%Todos aquellos que no tuvieron la oportunidad.
\par\end{flushright}{\large \par}
\end{quotation}
\newpage
\thispagestyle{empty}


% Agradecimientos
%Begin ----  Para que funcione bien el TOC en PDF

%\cleardoublepage

\textbf{\emph{\large }}\\
\vspace{4cm}

\begin{quotation}
\centering 
\textit{{\emph{1969 --- 2019}}}\\
\textit{{\emph{$50^{\text{th}}$ anniversary of the first human landing on the moon}}}

\end{quotation}
\begin{figure}[H]
	\centering
		\includegraphics[scale=0.55]{figurastfm/Chapter1/Imagenes/2.png}
	\label{fig:vostok1}
\end{figure}

\begin{quotation}
\centering 
\textbf{\emph{Science-fiction yesterday, fact today, obsolete tomorrow}}\\
\textit{{\emph{Otto O. Binder}}}\\

\end{quotation}

\afterpage{\blankpage}


\clearpage
\phantomsection
\addcontentsline{toc}{chapter}{Acknowledgments}


\begin{quotation}
\noindent \begin{center}
\textbf{\emph{\Large Acknowledgments:}}\textbf{\emph{\large }}\\
\textbf{\emph{\large }}\\
\textbf{\emph{\large }}\\
\textbf{\emph{\large }}
%Todos aquellos que no tuvieron la oportunidad.
\par\end{center}{\large \par}
\end{quotation}

This work has been possible thanks to a very small amount of people, but to whom I owe a great acknowledgment; although the one I can show you in these lines will undoubtedly be insufficient, serve this brief space as such.

The first of them must be for my family, my parents, José and Amparo, and my brother Francisco Javier. Ever since I can remember, their unconditional support in every aspect has been of primary importance for every of my undertakings to come to fruition. The support provided during this project has been only another proof of this.

To Luis and Alberto, for the technical support, the laughs, endless coffees and warnings to stop doing what you are doing. To Pablo, for the countless hours of advice, for the vision where mine got clouded, for the important things in life, for the wine and the debate near a bar.

This Master's Thesis would not have been born without the persistence of my tutor, Andrés Roldán Aranda, in the face of my initial reticence. His ambition has got me closer to numerous projects throughout all these years, and this work unbeatably closes this unforgettable stage.

And of course, to Natalia, for being with me during the course of this long stage, and keeping supporting me at the worst moments, for the understanding and the patience, thank you very much.

Everyone of the mentioned is part of this Thesis. To all of them, and to those who could not see it, thank you.
\afterpage{\blankpage}


\newpage
\afterpage{\blankpage}



\begin{quotation}
\noindent \begin{center}
\textbf{\emph{\Large Agradecimientos:}}\textbf{\emph{\large }}\\
\textbf{\emph{\large }}\\
\textbf{\emph{\large }}\\
\textbf{\emph{\large }}
%Todos aquellos que no tuvieron la oportunidad.
\par\end{center}{\large \par}
\end{quotation}

Este trabajo ha sido posible gracias a un número muy reducido de personas, pero a las cuales debo un gran agradecimiento; aunque el que pueda mostrarle en estas líneas será sin duda insuficiente, sirva este breve espacio como tal.

El primero de ellos ha de ir destinado a mi familia, mis padres, José y Amparo, y mi hermano Francisco Javier. Desde siempre, su apoyo incondicional en todos los aspectos ha resultado de importancia capital para que toda empresa acometida llegara a buen puerto. El apoyo mostrado durante el periodo de desarrollo de este proyecto no ha sido sino una muestra más del mismo.

A Luis y Alberto por el apoyo técnico, las risas, los cafés interminables y los avisos para dejar lo que estás haciendo. A Pablo, por las incontables horas de asesoramiento y visión allí donde la mía se nublaba, por las cosas  importantes de la vida, por el vino y el debate cerca de una barra.

Este Trabajo Fin de Máster no habría nacido sin la persistencia de mi tutor, Andrés~Roldán~Aranda, frente a mi reticencia inicial. Su ambición me ha acercado a numerosos proyectos a lo largo de todos estos años, y este trabajo cierra de manera inmejorable esta inolvidable etapa.

Y por supuesto, a Natalia, por acompañarme durante el transcurso de esta larga etapa, y seguir aún sirviendo de apoyo en los peores momentos, por la comprensión y la paciencia, muchas gracias.


Cada uno de los mencionados son partícipes de este Trabajo. A todos ellos, y a los que no pudieron verlo, muchas gracias.

\afterpage{\blankpage}

%\clearpage
%
%
%\bigskip
%
%
%\thispagestyle{empty}



% Lista de FIGURAS

%\pagenumbering{roman}%

\pagestyle{fancy}%
\cleardoublepage{\pagestyle{plain}\cleardoublepage}%

%Begin ----  Para que funcione bien el TOC en PDF
\cleardoublepage
\phantomsection \label{tableofcont}
%END  ---- Para que funcione bien el TOC en PDF

\begin{singlespacing}%onehalfspacing}
\tableofcontents
\end{singlespacing}%onehalfspacing}




%Begin ----  Para que funcione bien el TOC en PDF
\cleardoublepage
\phantomsection \label{listoffig}
%END  ---- Para que funcione bien el TOC en PDF
 
\begin{singlespacing}%onehalfspacing}
\listoffigures
\end{singlespacing}%onehalfspacing}


%\fancyhead{}
%\fancyhead[RO]{\small{\nouppercase{Índice de Figuras}}}
%\fancyhead[LE]{\small{\nouppercase{Índice de Figuras}}}

%\addcontentsline{toc}{chapter}{Lista de Figuras}%

\clearpage{\pagestyle{fancy}\cleardoublepage}%


% Lista de Videos
%Begin ----  Para que funcione bien el TOC en PDF
\cleardoublepage
\phantomsection \label{listofvid}
\addcontentsline{toc}{chapter}{List of Videos}
%END  ---- Para que funcione bien el TOC en PDF
 
\begin{singlespacing}%onehalfspacing}
\listof{videoFloat}{List of Videos}
\end{singlespacing}%onehalfspacing}


%\fancyhead{}
%\fancyhead[RO]{\small{\nouppercase{Índice de Figuras}}}
%\fancyhead[LE]{\small{\nouppercase{Índice de Figuras}}}

%\addcontentsline{toc}{chapter}{Lista de Figuras}%

\clearpage{\pagestyle{fancy}\cleardoublepage}%


% Lista de TABLAS
%Begin ----  Para que funcione bien el TOC en PDF
\cleardoublepage
\phantomsection \label{Tablass}
%\addcontentsline{toc}{chapter}{List of Tables}
%END  ---- Para que funcione bien el TOC en PDF
\begin{singlespacing}%onehalfspacing}
\listoftables
\end{singlespacing}%onehalfspacing}

%
%\fancyhead{}
%\fancyhead[RO]{\small{\nouppercase{Índice de Tablas}}}
%\fancyhead[LE]{\small{\nouppercase{Índice de Tablas}}}

%\addcontentsline{toc}{chapter}{Lista de Tablas}%

\clearpage{\pagestyle{fancy}\cleardoublepage}%




% Glosario
\printglossaries

%%%Símbolos: se incluyen en el listado de símbolos
      %\gls{pi}, \gls{ohm}, \gls{angstrom},
      %\gls{QIPrima},    \gls{QBPrima},    \gls{QGPrima}
      %\gls{QI},    \gls{QB},    \gls{QG}
      %\gls{qI},    \gls{qB},    \gls{qG}
      %\gls{Idb},    \gls{gm},    \gls{gmb},    \gls{gsd}
      %\gls{gbg},    \gls{gbs},    \gls{gbd}
%
%Acrónimos
%Enlace al glosario
%\glsname{HISIM} (\glstext{HISIM})
%%
%%%%%%% ERROR  %%%%%%%
%\acrshort{EKV} %(\acrlong{CRA})
%\glsname{EKV} (\glstext{EKV})

%%Citas: hace un enlace a la Bibliografía
%\cite{rytter1993vibration}
%\cite{frietzen2002CFRP}
%\cite{chang1999SHM}
%
%\refFigure{fig:matchgroup}
%%
%%
%%%Símbolos: se incluyen en el listado de símbolos
      %\gls{pi}, \gls{ohm}, \gls{angstrom},
      %\gls{QIPrima},    \gls{QBPrima},    \gls{QGPrima}
      %\gls{QI},    \gls{QB},    \gls{QG}
     % \gls{qI},    \gls{qB},    \gls{qG}
      %\gls{Idb},    \gls{gm},    \gls{gmb},    \gls{gsd}
     % \gls{gbg},    \gls{gbs},    \gls{gbd}
%
%Acrónimos
%\acrshort{SMD} (\acrlong{SMD})
%%Enlace al glosario
%\glsname{GranaSAT} (\glstext{GranaSAT})
%
%%%%%% ERROR  %%%%%%%
%\acrshort{CRA} %(\acrlong{CRA})
%\glsname{CRA} (\glstext{CRA})

%%Citas: hace un enlace a la Bibliografía
%\cite{definitionsACDS}
%\cite{AdafruitUSBtiny}
%\cite{ballaerospace}

%%Contador de PDF incluidos a pagina completa
\newcounter{includepdfpage}
\hypersetup{
    %bookmarks=true,% show bookmarks bar?
    %unicode=true, %allows to use characters of non-Latin based languages in Acrobat’s bookmarks
    bookmarksnumbered=true,
		pdfstartview={FitH},
		      %– FitH : Fit whole width of page
		      %– FitV : Fit whole height of page
		      %– FitB : Fit whole ”Bounding Box” page
		      %– FitBH: Fit whole width of ”Bounding Box” of page
		      %– FitBV: Fit whole height of ”Bounding Box” of page
		pdfpagemode=UseOutlines,
    pdftitle={Design of a multidisciplinary 1U CubeSat Simulation Platform},    % title
    pdfauthor={Jose Carlos Martinez Durillo},     % author
    pdfsubject={Master's Thesis},   % subject of the document
    %pdfkeywords={ambiental,station,noise,iot}, % list of keywords
    colorlinks=true,       % false: boxed links; true: colored links
    linkcolor=blue,       % color of internal links
    citecolor=green,        % color of links to bibliography
    filecolor=black,       % color of file links
    urlcolor=black         % color of external links.
		}


% Capítulo 1: Introducción

\chapter{Introducción}\label{chap:chapter1}
\pagenumbering{arabic}%

El presente proyecto se desarrolla en el marco del Trabajo de Fin de Grado con el cual se finaliza el Grado Universitario en Ingeniería Informática, con especialidad en Ingeniería del Software, cursado por el alumno Alejandro Ruiz Becerra en la Escuela Técnica Superior de Ingenierías Informática y de Telecomunicación de la Universidad de Granada. El objetivo principal del presente trabajo es demostrar los conocimientos y habilidades adquiridos por el alumno durante la realización de dicho grado. Para ello, se ha realizado el software de un limitador de sonido para locales con equipos de música.

Este TFG se realiza en colaboración con el proyecto académico \glsname{GranaSAT}, un proyecto multidisciplinario, dirigido por el Profesor Andrés María Roldán Aranda, que reúne a personas de diferentes campos que quieren adquirir conocimientos relativos a la Ingeniería Aeroespacial y a la Ingeniería Electrónica.


\begin{figure}[H]
    \centering
    \includegraphics[scale=0.9]{FigurasTFM/Chapter1/PDF/granasat.pdf}
    \caption{Logotipo de \glsname{GranaSAT}}
    %\vspace{-1cm}
\end{figure}


This Master's Thesis is presented as a compilation of the knowledge acquired throughout the years of the bachelor's and master's degree and specially, during this project period. It aims to reflect the engineering process behind the design, development, prototyping and verifying stage of a product. The overall goal of the project is developing an integrated 1U~\glsname{cubesat}~Simulation~Platform, to be used as a testing and training environment.

This ambitious project has been performed seeking a \textbf{double perspective} which will be detailed later, applying professional standards and methodologies used in space sector, in order to successfully overcome the project, minimizing risks and costs. It has supposed not only a hard design process but also a great labour of integration in order to maximize the use of existing equipment and procedures, available at \glsname{GranaSAT} laboratory. The final result is a fully-integrated simulation environment useful in a variety of situations: algorithms implementation, physical measurement or hardware characterization, among many others; its flexibility and expansion possibilities are countless. To make it possible, a conjunction of different fields of knowledge has been necessary; not only related to Telecommunication Engineering but also with completely different fields such as Mechanical or Aerospace Engineering. In addition to an academically and experimentally interesting platform, this conjunction has allowed developing a potentially marketable product, implying additional considerations to be taken into account.

In order to keep an structure, the project has been developed considering the different parts of the platform, i.e., \textbf{Mechanical Platform, Ground Station and Simulation CubeSat}. This document also follows that division, which eases understanding and allows a natural progress to the reader.


There is a generalized consensus regarding the best practices and methodologies around the design process in engineering; this project will follow one of the most extended methodologies in space sector, proposed by Wertz and Everett \cite{smad}, simplified when needed. As usual, it is characterized by an iterative refinement of the requirements which leads to the design and finally the verification of the objectives.


%The Product Development Process is composed of different stages: Technical~Requirements~Definition, System~Analysis, System~Design, Verification~and~Testing. The diagram of the figure \ref{proc} shows how that process flows since the initial step, which using a Market analogy could be the client proposal, until obtaining the final product. This process follows the planning proposed as a consequence of the Analysis, which can be seen at figure \ref{fig:gantt}. Once Verification and Testing stage has been adequately passed, should the product comply with functional and economics requirements, the product and therefore the Project will be finished.
%
%
		%\begin{figure}[H]
			%\centering
			%\includegraphics[scale=0.75]{FigurasTFM/Chapter1/PDF/designprocess.pdf}
			%\caption{Design Process} \label{proc}
		%\end{figure}



This Bachelor's Thesis fits within \href{http://granasat.ugr.es/}{\glsname{GranaSAT} Project}, a multidisciplinary project which gathers people from a variety of fields who are committed to acquiring new knowledge related to Electronics and Aerospace Engineering. Since its origins, one of its main purposes has been getting a \glsname{cubesat} in orbit, one of the reasons of this project; however, today its goals goes far beyond, and a wide range of devices and projects are developed in collaboration with different students and enterprises.

			\begin{figure}[H]
			\centering
			\includegraphics[scale=0.9]{FigurasTFM/Chapter1/PDF/granasat.pdf}
			\caption{\glsname{GranaSAT} Logo}
			%\vspace{-1cm}
			\end{figure}


\section{Motivation}\label{sec:mot}

%\begin{wrapfigure}{R}{0.3\textwidth}
	%\centering
	%\includegraphics[width=0.45\textwidth]{figurastfm/Chapter1/Imagenes/workshop.jpg}
	%\caption{\label{workshop}At \acrshort{ESA} CubeSats Workshop in Redu, Belgium (2018)}
%\end{wrapfigure}


After investing almost three years working on \glsname{GranaSAT} Project, dealing with hundred of problems and a variety of technologies, as well as being aware of its participation in Bexus Project \cite{gsatbexus} in 2014, I had an overall knowledge of the \glsname{cubesat} concept, main subsystems and typical objectives. However, after being selected to attend CubeSats Workshop organized by \textbf{European Space Agency \acrshort{ESA}} in Belgium, I realized that I needed to go in depth in the study of these fascinating satellites. \autoref{workshop} shows a photo of that great experience.

\begin{figure}
\centering
\includegraphics[width=75mm]{figurastfm/Chapter1/Imagenes/workshop.jpg}
\caption{At \acrshort{ESA} CubeSats Workshop in Redu, Belgium (2018)}
\label{workshop}
\end{figure}


Attending the aforementioned workshop made me realize the usefulness of using a simulated environment to acquire new concepts and increase technical knowledge, which is particularly interesting in an academic atmosphere such as the University. So my first motivation to develop this Master's Thesis was to provide the \glsname{GranaSAT} laboratory with a multidisciplinary Simulation Platform for CubeSats, taking the advantage of their little size, in contrast with bigger satellites, much more complex to test in a non-professional environment; as it is widely known, after Bexus Project (see \cite{gsatbexus}) \glsname{GranaSAT}~Group has kept working on real developments and missions in collaboration with \acrshort{ESA}, particularly around CubeSats, with the first \glsname{cubesat} of the University of Granada, \textbf{GranaSAT-I} intending to participate in the \acrshort{ESA} programme, \textbf{\textit{Fly Your Satellite!}} \cite{flyyour}, introduced in \autoref{fys}.

For all these reasons, I decided to focus my Master's Thesis on a \textbf{double perspective}: on the one hand, a multidisciplinary academic-oriented \textbf{Simulation Platform} for CubeSats, which allows testing its different subsystems as well as implementing innovative solutions in a physical scenario; on the other hand, a complete project-oriented platform to be used by future \glsname{GranaSAT} members	on its way to the \textit{Fly Your Satellite!} programme \cite{flyyour}, including a \glsname{cubesat} \glsname{engmodel} with different subsystems to be tested and fully integrable with the countless technology modules and procedures which \glsname{GranaSAT} Project has been provided with throughout the years, as result of the work of voluntary students and researchers; precisely, some of them are used in this Master's Thesis, so interoperability and cooperation with the rest of the technical staff of the group have also been important incentives to me; this combination of existing technologies and equipment is intrinsic to complex project and aerospace design in particular: a large group of professionals working towards a same goal throughout time, in this case, \textbf{flying a \glsname{cubesat}}. \autoref{perspc} depicts these perspectives.


			\begin{figure}[H]
			\centering
			\includegraphics[scale=1]{FigurasTFM/Chapter1/PDF/doublepersp2.pdf}
			\caption{Double orientation of the project} \label{perspc}
			%\vspace{-cm}
			\end{figure}


This double orientation of the project will entail innumerable mentions through this document, as some decisions will have to be taken in order to comply with both perspectives. Generally speaking, my intention has been keeping an adequate balance which allows the Simulation Platform to \textbf{satisfy both testing and academic needs.}

After an arduous Bachelor's Thesis \cite{tfg} \textbf{awarded twice} by the \textit{Colegio Oficial De Ingenieros Técnicos De Telecomunicación} \cite{coitt} and \textit{Escuela Técnica Superior de Ingenierías Informática y de Telecomunicación} at the University of Granada \cite{etsiit}, this project has supposed a definitely larger and exciting challenge, substantially increasing my professional experience and knowledge in aerospace field by dealing with space standards, learning different systems architecture and putting into practice numerously experimental techniques.

The Simulation Platform which emerges from this Master's Thesis is the result of a painstaking process of engineering, study and comparison of the cutting-edge technology used in space which I sincerely hope that can be useful for the \glsname{GranaSAT} Project in the future.


\section{Prior art. Problem Statement}\label{sec:iot}

\subsection{Background study: CubeSats} \label{cubesatdefin}

\glsname{cubesat} emerged as a result of increasing costs in space field, incredible amounts of money had to be spent in missions with a considerable risk. They allowed to trade some functionalities off in exchange for a significant reduction in costs and therefore a greater number of missions; higher risk of failure is accepted because of the possibility to perform more and more frequent launches, as illustrated in \autoref{increase}.

	\begin{figure}[H]
			\centering
			\includegraphics[scale=0.75]{FigurasTFM/Chapter1/PDF/increse.pdf}
			\caption{Number of CubeSats launched over the last years \cite{usage}} \label{increase}
			\vspace{-1.5\baselineskip}
			\end{figure}

As for the academic scope, \glsname{cubesat} allowed people with little or no experience in the design of space missions to incorporate new ideas and concepts into missions with fewer restrictions. Within university projects, the challenge is not only accomplishing a successful mission but also get it at the lowest cost. All of which, of course, supposes the risk to be inherently high, but it is not a problem itself.

\glsname{cubesat} concept were first developed at Stanford and California Polytechnic universities in 1998. It was precisely an Spanish engineer, \textbf{Jordi Puig-Suari} along with Robert Twiggs, who created them. The primary intent of the development of the CubeSat was to provide a standard set of dimensions for external physical structure, in order to make it compatible with a standardized orbital deployer. The final concept of the CubeSat structural standard is shown in \autoref{cubesatstd}. It is based on a \textbf{10 x 10 x 10 cm cube and a mass of 1.33~kg} maximum (see Figure\autoref{c1u}) and is referred to as a one unit `1U' and is scalable in 1U increments, as shown in Figure\autoref{stds}. This standardization facilitates a wide variety of designs to be launched and deployed using a common system; the most used is Poly-Picosatellite Orbital Deployer (\acrshort{PPOD}), having launched all CubeSats since 2007. Figure\autoref{ppodexam} shows a typical \acrshort{PPOD} launcher while Figure\autoref{csatlaunching} depicts three 1U CubeSats being launched.




\newpage
\makeatletter
\setlength\@fptop{0pt}
\setlength\@fpsep{8pt plus 1fil}
\setlength\@fpbot{0pt}
\makeatother
\begin{figure}[t]
			\centering
			\subfloat[1U standard\label{c1u}]{\includegraphics[page=1,trim={0cm 0cm 0cm 0cm},clip=true,width=60mm]{figurastfm/Chapter1/Imagenes/c1u.jpg}}
			 			 \quad
			\subfloat[1U, 2U and 3U standards comparison\label{stds}]{\includegraphics[page=1,trim={0cm 0cm 0cm 0cm},clip=true,width=75mm]{figurastfm/Chapter1/Imagenes/variosstd.png}}

			\caption{\glsname{cubesat} physical structure \cite{nasa}}
\end{figure}

\begin{figure}[H]
			\centering
			\subfloat[3U capacity example \cite{xinlin}\label{ppodexam}]{\includegraphics[page=1,trim={0cm 0cm 0cm 0cm},clip=true,width=70mm]{figurastfm/Chapter1/Imagenes/ppod.png}}

			\subfloat[\glsname{cubesat} launching \cite{nasa}\label{csatlaunching}]{\includegraphics[page=1,trim={0cm 0cm 0cm 0cm},clip=true,width=105mm]{figurastfm/Chapter1/PDF/cubesatlaunching.pdf}}
			\caption{\acrshort{PPOD} Launcher}
\end{figure}

			%\afterpage{
\begin{landscape}
\begin{figure}
			\centering
			\includegraphics[page=1,trim={0cm 0cm 0cm 0cm},clip=true,width=230mm]{figurastfm/Chapter3/PDF/cubesatstd.pdf}
			\caption{Standard for CubeSat Mechanical Structure \cite{calpoly}} \label{cubesatstd}
\end{figure}
\end{landscape}
%}

CubeSats are typically composed of different subsystems, which are more or less standard in literature. This project will follow the division used by \acrshort{ESA}, depicted in \autoref{cubesatsub}; besides, the most relevant components and functions of each subsystem are also included. These subsystems are briefly analyzed next.


\begin{figure}[H]
				\centering
				\includegraphics[width=165mm]{figurastfm/Chapter1/PDF/csatsubs.pdf}
				\caption{\glsname{cubesat} typical subsystems}
				\label{cubesatsub}
\end{figure}

\begin{itemize} [noitemsep]

	\item \textbf{On-board Computer (\acrshort{OBC}):} the main subsystem of a \glsname{cubesat}, it is in charge of the core processing, reading and polling sensors, R/W operations on memory, data management, etc. It is typically based on a central microprocessor or computing core. \\

	\item \textbf{Communications (\acrshort{COMMS}):} in charge of sending and receiving all the information related to a mission. It is sometimes taken as the most critical subsystem, given that \textbf{without contact the mission is over,} even if the rest of the system functions flawlessly. It is composed of the transmission and receiving antennas, amplifiers, \acrshort{RF} electronics, filtering, signal couplers, among many others elements. \\

	\item \textbf{Attitude Determination \& Control System (\acrshort{ADCS}):} its main purpose is assuring the correct orientation of the CubeSat and dealing with external torques. Composed of sensors and actuators such as \glsname{magnetorquers}, it usually implements an independent processing unit. As for its software, it can be based on different \textbf{control laws} and predictive filters such as \textbf{Kalman filter}.\\

	\item \textbf{Electrical Power System (\acrshort{EPS}):} this module deals with the generation, distribution and control of the spacecraft electrical power. It must supply a continuous source of power in adequate condition for each element of the system, and deal with possible faulty situations. It is composed of a long chain of elements, from solar arrays to voltage regulators or batteries.\\

	\item \textbf{Thermal Subsystem (\acrshort{THS}):} in charge of setting and maintaining the temperature range of the \glsname{cubesat} and its components, which are all designed to operate over a defined temperature range. This is accomplished by balancing the flow of heat energy. Generally, in CubeSats it is not required an active thermal control system.\\

		\item \textbf{On-Board Data Handling (\acrshort{OBDH}):} it handles all data sent and received by the CubeSat, including \acrshort{TCC}. Although it is typically implemented on the \acrshort{OBC} and connected to the \acrshort{COMMS} subsystem, because of its increasing importance during the last few years, it is nowadays considered a proper subsystem by itself.

\end{itemize}

Additionally, sometimes the \textbf{payload} is taken as another subsystem of the \glsname{cubesat}. It is actually the most important part of it given that it constitutes the reason of the mission itself: \textbf{the payload provides the expected results and measurements}. There is a huge variety of payloads, with equal amount of purposes; cameras, sensors, radiometers... the list is endless. In sum, anything able to provide the team with some kind of result or just even functioning correctly is liable to be a payload.

This project will focus on \acrshort{OBC}, \acrshort{ADCS} and in the power source area of the \acrshort{EPS}. Also, some basic functions related to \acrshort{OBDH} will be implemented. \autoref{cubesatanaly} deeply analyzes these subsystems.


\subsection{ESA \textit{Fly Your Satellite!}} \label{fys}

Because of the increasing interest within university students in CubeSats, \acrshort{ESA} decided to boost this initiative in 2013 by organizing the first \textit{Fly Your Satellite!} program in collaboration with European Universities, with the goal to be a complementary preparation for the future professionals of the space sector.

During the program, the accepted student teams are supported in the development of their \glsname{cubesat} and their previously defined mission. \acrshort{ESA} experts support and mentor the teams through the different phases, focusing on assembly, integration, verification, testing and operations. At a point, the students are allowed to use the Agency test facilities and take part in training sessions. This priceless opportunity allows the students to gain great experience in the lifecycle of an actual space project, using working methods and standards from professional space programmes which significantly help to build their own professional skills \cite{ESA}.

ESA \textit{Fly Your Satellite!} programme is divided into different phases with pre-defined objectives and reviews conducted by \acrshort{ESA} specialists. Gaining access to the next phase requires to successfully complete the previous one. These phases are depicted in the flow diagram of \autoref{flyflow}.

\begin{figure}[H]
				\centering
				\includegraphics[width=\textwidth]{figurastfm/Chapter1/PDF/cubesat_flow.pdf}
				\caption{\textit{Fly Your Satellite!} phases} 	\label{flyflow}
				\vspace{-1cm}
\end{figure}


As illustrated in the diagram, application to the program is constrained to those teams which already have a complete detailed design. As previously mentioned, one of the main goals of the \glsname{GranaSAT} Project is successfully participating in this program. This Master's~Thesis lays the foundations for a future \glsname{cubesat} design and application to the \textit{Fly~Your~Satellite!}.


\subsection{European Cooperation for Space Standardization}\label{sec:ECSS}


The European Cooperation for Space Standardization (\glsname{ECSS}) is an initiative established to develop a coherent, single set of user-friendly standards for use in all European space activities. It started officially in the autumn of 1993, when the European space community realized that differences in space standards resulted in higher costs, lower effectiveness and a less competitive industry.

\begin{figure}[H]
				\centering
				\includegraphics[width=60mm]{figurastfm/Chapter2/Imagenes/ECSS_large.jpg}
				\vspace{0.3cm}
				\caption{\glsname{ECSS} logo \cite{ecss}}
				%\label{fig:minitio}
								\vspace{-1cm}

\end{figure}

Numerous ECSS standards have been released with the aim to be applied altogether for the management, engineering and product assurance. Standards and recommendations for the development of the space technology are split up in four disciplines: \textit{\textcolor{red}{Space project management branch (M)}}, \textcolor{blue}{Space product assurance branch (Q)}, \textcolor{green}{Space engineering branch (E)}, and \textcolor{orange}{Space sustainability branch (U)}. \autoref{ecssdisciplines} shows the different \glsname{ECSS} disciplines \cite{ecssdescrip}.

		\begin{figure}[h]
			\centering
			\includegraphics[page=1,trim={0cm 0cm 0cm 0cm},clip=true,width=150mm]{figurastfm/Chapter2/PDF/ecssdisciplines.pdf}
			\caption{ECSS Document Tree \cite{ecsstree}}
			\label{ecssdisciplines}
		\end{figure}



As shown in \autoref{ecsseng}, the documentation is organized into four branches, each one defined to the detail level required to differentiate major functions, disciplines and activities. Besides, four hierarchical levels are considered:

\begin{itemize}

\item{\textbf{Level 0 (ECSS-P-00)} – The document at Level 0 describes the policy and objectives of the ECSS system and its architecture as well as the principle rule for the creation, validation and maintenance of documents.}
\item{\textbf{Level 1 (ECSS-M-00, ECSS-Q-00, ECSS-E-00, ECSS-U-00)} -The documents at Level 1 describe the strategy in the specific domain. They give a global view of the requirements and outline the interfaces between the elements (and the documents) at Level 2.}
\item{\textbf{Level 2 (ECSS-M-10, ECSS-Q-10 …)} – The documents at Level 2 describe the required objectives and functions for all aspects in the individual domain (project organization, quality assurance, system engineering, etc.).}
\item{\textbf{Level 3} – The documents at Level 3 describe methods, procedures and recommended tools to achieve the requirements of Level 2 documents. In addition they define the constraints and requirements for interfaces and performance of the specified product or activity. The Level 3 documents are guidelines and are allowed to be adapted to the projects’ needs.}

\end{itemize}

Also, as a reference, and given that a large number of them will be considered for this project, space engineering branch is broken down and included in \autoref{ecsseng}.

%\afterpage{
\begin{landscape}
			\begin{figure}
				\centering
				\includegraphics[page=1,trim={0cm 0cm 0cm 0cm},clip=true,width=230mm]{figurastfm/Chapter2/PDF/ecssengbranch.pdf}
				\caption{ECSS Engineering branch breakdown \cite{ecsstree}}
				\label{ecsseng}
			\end{figure}
\end{landscape}%}




\section{Introduction to System Engineering. EDP in Space.}\label{engprocess}

Systems engineering is an interdisciplinary field which involves engineering management and focuses on how to design and manage complex systems over their life cycles \cite{wiki}.

A system is an integrated set of elements to accomplish a defined objective. These elements include  hardware, software, firmware, human resources, information, techniques, facilities services, and other support elements.

 \autoref{fig:sysengfig} shows, according to ECSS-E-ST-10C \cite{ESASYSENG}, the boundaries of system engineering, its relationship with production, operations, product assurance and management disciplines and its internal partition into the following system engineering sub-functions:

\begin{itemize}

	\item{\textbf{Requirement engineering}, which consists on requirement analysis and validation, requirement allocation, and requirement maintenance;}
	\item{\textbf{Analysis}, which is performed for the purpose of resolving requirements conflicts, decomposing and allocating requirements;}
	\item{\textbf{Design and configuration}, which results in a physical architecture, and its complete system functional, physical and software characteristics;}
	\item{\textbf{Verification}, which objective is to demonstrate that the deliverables conform to the specified requirements;}
	\item{\textbf{System engineering integration and control}, coordinating the various engineering disciplines and participants;}

	\end{itemize}

			\begin{figure}[h]
			\centering
			\includegraphics[page=1,trim={0cm 0cm 0cm 0cm},clip=true,width=150mm]{figurastfm/Chapter2/PDF/syseng.pdf}
			\caption{System engineering, sub-functions and boundaries \cite{ESASYSENG}}
			\label{fig:sysengfig}
			\end{figure}

Through this process the system engineering function performs a multidisciplinary functional decomposition to obtain logical lower level products.

Along with the variety of disciplines in engineering, several design processes have been proposed throughout time whether in a general way or a field-specific. Space field is not an exception to that and there are different models applicable in order to accomplish a systematic and effective design process. It is the so-called \textbf{Engineering Design Process} or \acrshort{EDP}.

Within space sector, the highest level is often called \textbf{\textit{mission level}}, which usually consists of one or more segments (space, ground, user...). From this perspective, requirements originate from the next upper level (the customer) and elements are procured from the next lower level (the suppliers).

This project will apply one of the most extended design methodology, proposed by Wertz and Everett, see \cite{smad}. It divides the design process into a series of phases or steps, which coincides with the structure followed not only during the development of this project but also in this document itself.

As previously stated, any design process is in some way focused on \textbf{requirements}. Requirements definition is a complex and iterative process which begins with one or more objectives and continues defining a system which meets them as much as possible, with the \textbf{minimum cost and risk}. While overall objectives normally remain the same through the whole process, they can be accomplished in many different ways according to the technology or solution applied. The Design process followed in Space field breakdowns into different phases:

\begin{itemize} [noitemsep,topsep=0pt]

	\item \textbf{Objectives/goals definition.} Expectable results from the project, not only in technical terms but also economic, social, political, etc. This is the very first stage of the design process and defined goals here are not supposed to vary during the process. They are treated in \autoref{sec:obj}. \\

	\item \textbf{Space Mission Engineering.} Inherently iterative process in charge of refining both the requirements and methods to accomplish them. Each iteration will lead to a more detailed and defined mission concept. Primary and secondary objectives are specified, \glsname{stakeholders} interests are considered, the timeline is defined, etc. It is treated in \autoref{chap:chapter2}. \\

	\item \textbf{Space System Engineering.} The previously preliminary requirements defined are translated into a set of specific requirements on every element of the system. This translation will determine performance, cost and risk. It is extensively treated in \autoref{chap:chapter3}. \\

	\item \textbf{System Design.} The project is designed according to the requirements stated in previous stages, comparing different technologies and implementing the optimal one for each subsystem. It is treated in \autoref{chap:chapter4}. \\

	\item \textbf{Verification \& testing.} The project is concluded by conducting different functionality tests which verifies the results and checks the compliance with the requirements. It is exposed in \autoref{chap:chapter5}. \\

\end{itemize}

\autoref{sysengflow} briefly illustrates this flow. This project will follow a methodology as close as possible to this system engineering perspective, according to \cite{ecss} standards.

	\begin{figure}[H]
			\centering
			\includegraphics[scale=0.75]{FigurasTFM/Chapter1/PDF/syseng.pdf}
 		\caption{Engineering Design Process (\acrshort{EDP}) in space} \label{sysengflow}
				\vspace{-2cm}
		\end{figure}


\section{Project Goals and Objectives}\label{sec:obj}


After introducing the problem statement which motivates this Master's Thesis, this section outlines the main top-level non-technical goals of the project. The reason to limit this section to those is to stick to the \textbf{Engineering Design Process} described in \autoref{engprocess} and followed throughout this Master's Thesis, which considers objectives definition part of a much more complex and iterative process called \textit{Requirements Definition}. Besides, considering again the double perspective (technical testing and academic) explained in previous sections, and the substantial extent of this project, I decided to breakdown the objectives and functional requirements of each subsystem in \autoref{chap:chapter2} \textbf{(Mission Engineering)}, while formal requirements are addressed in \autoref{chap:chapter3} \textbf{(System Engineering)}. Then, \autoref{chap:chapter4} addresses the design of the system in accordance with the exposed in previous chapters.

Therefore, objectives listed in \autoref{topobj} must be understood as the author's expected results in academic and professional terms of the execution of this project.

\begin{table} [H]
\centering

\begin{tabularx}{\linewidth}{lX}

\multicolumn{1}{c}{\textbf{Ref.}}                      & \multicolumn{1}{c}{\textbf{Objective}}                    \tabularnewline \specialrule{1.1pt}{1pt}{1pt}
Obj.1                                              & To gain experience with different standards typically used in aerospace field such as \acrshort{ECSS} or \acrshort{CCSDS}.                                \tabularnewline \midrule
Obj.2                                               & To familiarize with professional Requirements Refinement techniques by applying the Engineering Model Process, particularly oriented to space field.                                                                    \tabularnewline \midrule
Obj.3                                             & To generate enough and clear documentation of the whole process, which may be required during the same project or just useful for future references or designs.                \tabularnewline \midrule
Obj.4                                                    &  To dynamically acquire multidisciplinary knowledge of different specialties as needed during the execution of the project, featuring Mechanical and Aerospace Engineering concepts and techniques.    \tabularnewline \midrule

Obj.5                                                    &  To demonstrate the knowledge acquired during both, Bachelor's and Master's degree in Telecommunication Engineering, as well as multidisciplinary abilities gathered during the execution of this Master's Thesis.      \tabularnewline \midrule

Obj.6                                                    &   To successfully overcome the subject ‘Master's Thesis’.     \tabularnewline \midrule

\end{tabularx}
\caption{Top-level objectives of the project}
%\vspace{-1.3\baselineskip}
\label{topobj}

\end{table}




\newpage
\section{Project Structure}\label{sec:chapterDescription}

This project, divided into six chapters and an addendum, progressively analyzes the system under development from different points of view, addresses the design tasks and finalizes with the successful completion of the product.

These chapters are:

\begin{itemize} [noitemsep,topsep=0pt]
\item \hyperref[chap:chapter1]{\textbf{Chapter one}}. This chapter, which is intended to be an introduction and show the general objectives and the reasons which motivate this project. The prior art, as well as a brief introduction to the system engineering methodology applied in this development, are also included in this chapter.

\item \hyperref[chap:chapter2]{\textbf{Chapter two}}. The second chapter addresses the so-called stage \textit{Space~\textbf{Mission}~Engineering}, in which the technical objectives of the system are remarked from a high abstraction level. These objectives lead to the \textbf{Functional~Requirements~Definition}.


\item \hyperref[chap:chapter3]{\textbf{Chapter three}}. According to the system engineering methodology applied, the third stage of the project corresponds to the \textit{Space \textbf{System} Engineering}, so it is treated in this chapter. It is included a short introduction to the testing procedures needed to comply with a real mission constraints, including several \acrshort{ECSS} standards. Besides, in this chapter the project is breakdown into three separated blocks; each of them is analyzed, and the \textbf{Formal Requirements} are defined, as a previous step to the design stage.

\item \hyperref[chap:chapter4]{\textbf{Chapter four}}. The fourth chapter deals with \textbf{system design}. It translates the technological solutions analyzed in the previous chapter to actual systems able to execute the tasks required. The blocks structure of the project introduced in the previous chapter is followed again, and each of them is extensively treated, including details at all levels of the design task.
\item \hyperref[chap:chapter5]{\textbf{Chapter five}}. Once the design is completed, chapter five addresses a series of validation tests, in order to check that the system meets the \textbf{Formal~Requirements} defined and, consequently, the Functional Requirements. It is worth noting that this validation is performed using one of the tools developed within this project.
\item \hyperref[chap:chapter6]{\textbf{Chapter six}}. Finally, chapter six includes the main conclusions extracted from this Master's Thesis, as well as some future lines of work which have naturally emerged during the design process.

\item \hyperref[cap:budget]{\textbf{Addendum}}, in which the budget and associated cost of this project are detailed.
\end{itemize}

			\begin{figure}[H]
			\centering
			\includegraphics[page=1,trim={0cm 45cm 49cm 0cm},clip=true,width=170mm]{FigurasTFM/Chapter1/PDF/gantt.pdf}
			\caption{Gantt Chart of the Project}
			\label{fig:gantt}
			\end{figure}

















% Capítulo 2: Estado del arte
\chapter{Space Mission Engineering} \label{chap:chapter2}

\textbf{Space Mission Engineering} is a broad process in charge of determining the needs of the \textit{client} or goals of the system. Those initial needs are not likely to change and can be considered the \textbf{objectives of the mission}. In that line, this chapter will address top-level objectives, analyzing them and culminating with the task \textbf{Functional Requirements Definition} for each of the subdivisions of the system. 

Furthermore, a brief introduction to this stage of the project is included, in relation with both, generally speaking missions and this particular project. Regarding this point, it may be needed to remember that the double scope of this Master's Thesis is at the halfway between a real mission design and the development of a simulation platform, which will imply taking balanced decisions when an excessively mission-oriented requirement is needed, i.e., space-qualified materials will not be treated in this project. Those assumptions shall be taken into account for future developments specifically intended to fly in orbit. 


\section{Introduction to Mission Engineering} \label{missioneng}%no sé si merece la pena este titulo

Mission Engineering typically begins with one or more general objectives and constraints which build up the base to later define the functional requirements of the system, i.e., which tasks should the system be able to perform, from the \textit{client} (mission) point of view. It is usually an iterative process, consisting in refining the requirements expected, being aware of the \textbf{capabilities} and \textbf{constraints} of the technical team and equipment. 
Documenting the results of this iterative process must also be part of the philosophy of work, so decisions can be reexamined and justified if needed. 

In a real mission, setting excessively optimistic or ambitious objectives (with the aim of getting the greatest amount of results, for example) may be counterproductive, taking to a delayed delivery of the system (increasing costs) or, in the worst scenario, to a failure of the project. Indeed, not only too optimistic goals are risky, overdoing details of the mission in any sense may lead to unexpected problems once the mission is analyzed as a whole. At this point, the main focus must be to translate the needs and objectives of the end user into functional requirements.

As analyzed in \cite{smad}, this stage typically implies a series of steps through which the team must wonder themselves the questions that will lead to the definition of the needs and requirements. Some of these are:

\begin{itemize}

	\item{What are the qualitative goals seeked and why?}
	\item{Who are the \glsname{stakeholders}?}
	\item{Which is the timelime?}
	\item{Considering needs, available technology, previous know-how and so forth, how well should the broad objectives be achieved?}
	
\end{itemize}

Therefore, Mission Engineering is a crucial stage of the \textbf{Engineering Design Process} (\acrshort{EDP}) in space, in order to get the most performance possible for the inversion made. Its importance is directly related to the early point of the project in which they are defined; because they are stated as the goals to achieve, the design will point to those from the beginning, implying therefore seeing a significant cost increase in the project as a result of any eventual change to them. as well as additional, avoidable risk. All in all, the final goal is to successfully fulfill the needs of the mission by meeting the end-clients' objectives in time at minimum risk and cost.

Considering the double perspective of this project, as previously exposed in \autoref{chap:chapter1}, these recommendations apply specially for the prototype-oriented focus, as developing a system intended to be in orbit. While the Simulation Platform branch allows relaxing them in some ways, being an academic project with a certain \textbf{deadline} makes it recommendable to keep risks under control in order to stick to the schedule and conclude it in time. In any case, as stated in \autoref{cubesatdefin}, because of the nature of CubeSats itself, a certain amount of risk is accepted in these kind of projects; trying new ideas or concepts inherently \textbf{includes that risk}, even more so when low cost is intended. Accomplishing the mission under that condition will be the real challenge.


\section{System Objectives. Functional Requirements Definition}

After introducing the foundations of the Mission Engineering Process, the technical objectives of the system are defined, in contrast with the ones defined in \autoref{sec:obj}. They fit the top-level needs and will lead to the definition of the \textbf{functional requirements} of this project, after performing the iterative process stated in \autoref{missioneng}; firstly, the breakdown of the parts of this project is shown. The proposed simulation platform will be composed of three different subsystems, illustrated in \autoref{decomp}, in which the main purposes and functions of each are featured.

\begin{figure}
				\centering
				\includegraphics[width=160mm]{figurastfm/Chapter2/PDF/decomp.pdf}
				\caption{Simulation Platform Subsystems}      		
				\label{decomp}
\end{figure}


As seen in \autoref{decomp}, this project involves three clearly differentiated parts, which will be broken down in later chapters. Following the philosophy of this stage, now they are considered from a top-level mission's view, determining the technical objectives of the project for each one.


\subsection{Mechanical Platform}

Although, as illustrated in \autoref{decomp} this project as a whole has the final purpose to develop a complete Simulation Platform framework, one of the main differentiated parts is specifically a physical simulation platform which allows testing a CubeSat prototype, as well as interacting with it in every possible way. Therefore, this  part of the platform will require an extensive work of \textbf{mechanical design}. Also, this subpart of the system answers both orientations of this Master's Thesis: on the one hand, it allows technical testing of functionalities, which can be useful at early stages of the development, before going through official testing (see \autoref{testing}) so the approximate performance of the prototype can be analyzed; additionally, it allows determining the capabilities of the prototype, feasibility of the design, changes needed, etc. On the other hand, it also satisfies the educational aspect of this project, by allowing students to implement different solutions in a space-oriented platform, easing the acquisition of specific concepts and putting it into practice. Taking this into account, the set of objectives related to the mechanical platform is shown in  \autoref{techobjmechplat}. Note that both objectives and functional requirements are identified by a code which eases a quick reference when needed; this notation will be used through all this project when dealing with requirements.

\begin{table} [H]
\centering

\begin{tabularx}{\linewidth}{lX}
 & \multicolumn{1}{c}{\textbf{Technical Objectives}}    
\tabularnewline \specialrule{1.1pt}{1pt}{1pt}

\multicolumn{1}{c}{\textbf{Ref.}}                      & \multicolumn{1}{c}{\textbf{Primary}}                    \tabularnewline \specialrule{1.1pt}{1pt}{1pt}
MP.Obj.1                                              & To allow testing at least 1U \glsname{cubesat}, physically rotating, emulating in orbit
flight.                                                                                  \tabularnewline \midrule
MP.Obj.2                                              & To allow integration of external elements to the platform, such as one or more lightning source emulating the Sun, magnetic field presence, etc.                                                                                 \tabularnewline \specialrule{1.1pt}{1pt}{1pt}
\multicolumn{1}{c}{\textbf{Ref.}}                      & \multicolumn{1}{c}{\textbf{Secondary}}                    \tabularnewline \specialrule{1.1pt}{1pt}{1pt}
MP.Obj.3                                            & To keep it as lightweight as possible, maximizing portability.                   \tabularnewline \midrule

MP.Obj.4                                                   & To build it using a low-cost philosophy, specially in prototyping stage. \tabularnewline \midrule
MP.Obj.5                                                   & The subsystem shall use inexpensive materials in prototyping stage.    \tabularnewline \midrule
\end{tabularx}

\caption{Mechanical Platform - Technical Objectives}
%\vspace{-1.3\baselineskip}
\label{techobjmechplat}   

\end{table}


While \textbf{Primary Objectives} are those to be fulfilled in order to consider a certain level of success, \textbf{Secondary Objectives} are additional goals which would complement the primary list. All in all, they express the needs of the client and/or mission. From these, functional requirements can be extracted, as shown in \autoref{functmechplat}.

\begin{table}[H]
\centering

\begin{tabularx}{\linewidth}{lX}

\multicolumn{1}{c}{\textbf{Ref.}}                      & \multicolumn{1}{c}{\textbf{Functional Requirements}}                    \tabularnewline \specialrule{1.1pt}{1pt}{1pt}
MP.FR.1                                              & The subsystem shall have enough capacity to host at least 1U \glsname{cubesat}.                                    \tabularnewline \midrule
MP.FR.2                                              & The subsystem shall have low friction in order to emulate in orbit conditions.                                                                      \tabularnewline \midrule
MP.FR.3                                            & The subsystem shall weight as less as possible (without \glsname{cubesat} unit).                 \tabularnewline \midrule
%MP.FR.4                                                   & The subsystem shall use inexpensive materials in prototyping stage.    \tabularnewline \midrule

\end{tabularx}
\caption{Mechanical Platform - Functional Requirements}
%\vspace{-1.3\baselineskip}
\label{functmechplat}

\end{table}
\subsection{Ground Station} \label{gstation}

When talking about a satellite link, \glsname{ground} is the Earth's-end in charge of dealing with the communications tasks from and to the spacecraft. In the same way, this project will make use of the \glsname{ground} to establish a communication link with the \glsname{cubesat} under test, allowing testing \acrshort{TCC} System, but will also function as central server of a hypothetical network of CubeSats, which would be certainly useful in a classroom environment, in which each student would have their \glsname{cubesat} but only one \glsname{ground} is needed. Once again, this subpart satisfies both orientations of this Master's Thesis. \autoref{techobjgstation} and \autoref{functgstation} shows the technical objectives and functional requirements of this part, respectively.


\begin{table} [h]
\centering

\begin{tabularx}{\linewidth}{lX}
 & \multicolumn{1}{c}{\textbf{Technical Objectives}}    
\tabularnewline \specialrule{1.1pt}{1pt}{1pt}

\multicolumn{1}{c}{\textbf{Ref.}}                      & \multicolumn{1}{c}{\textbf{Primary}}                    \tabularnewline \specialrule{1.1pt}{1pt}{1pt}
GS.Obj.1                                              & To allow \glsname{ground} operators to effectively communicate with GranaSAT-I \glsname{cubesat}.                                                                                 \tabularnewline \midrule
GS.Obj.2                                              & To allow sending control commands to GranaSAT-I.                                                                                                         \tabularnewline \midrule
GS.Obj.3                                              & To allow receiving telemetry and payload data from GranaSAT-I.        \tabularnewline \midrule
GS.Obj.4                                              & To allow defining new control commands and telemetry packages. \tabularnewline


\specialrule{1.1pt}{1pt}{1pt}
\multicolumn{1}{c}{\textbf{Ref.}}                      & \multicolumn{1}{c}{\textbf{Secondary}}                    \tabularnewline \specialrule{1.1pt}{1pt}{1pt}
GS.Obj.5                                            & 
To function as central \glsname{ground} for a number of compatible CubeSats, potentially a swarm, in a local environment. It may be useful as approach to the real case stated in \cite{paperamericano} in which I collaborated with Professor Stakem from \textbf{Goddard Space Flight Center} (\acrshort{NASA}). \tabularnewline \midrule

GS.Obj.6                                                   & To count with software analysis in charge of managing transmitted and received data. \tabularnewline \midrule
GS.Obj.7                                                   & To communicate with GranaSAT-I using industry standard formats. \tabularnewline \midrule
\end{tabularx}
\caption{Ground Station - Technical Objectives}
%\vspace{-1.3\baselineskip}
\label{techobjgstation}   

\end{table}



\begin{table}[h]
\centering

\begin{tabularx}{\linewidth}{lX}

\multicolumn{1}{c}{\textbf{Ref.}}                      & \multicolumn{1}{c}{\textbf{Functional Requirements}}                    \tabularnewline \specialrule{1.1pt}{1pt}{1pt}
GS.FR.1                                              & The subsystem shall have a communication system compatible with \glsname{cubesat} GranaSAT-I,  in terms of frequency, technology and whatever other aspect needed to accomplish an effective communication.                                                       \tabularnewline \midrule
GS.FR.2                                              & The subsystem shall have a software module in charge of packetization and de-packetization.                                                                      \tabularnewline \midrule
GS.FR.3                                            & The subsystem shall have a software module in charge of keeping logs, plotting received data, managing alerts and any other capability liable to implemented. It shall allow data analysis in both, real time and post-processing stage.                   \tabularnewline \midrule
GS.FR.4                                                   & The subsystem shall be able to function in a distributed local network of compatible CubeSats, managing communications, connections, etc. \tabularnewline \midrule
GS.FR.5                                                   & To communicate with GranaSAT-I using \acrshort{CCSDS} (see \ref{ccsds}). \tabularnewline \midrule
\end{tabularx}
\caption{Ground Station - Functional Requirements}
%\vspace{-1.3\baselineskip}
\label{functgstation}   

\end{table}


\subsection{Simulation CubeSat}

The third top-level subpart of the system is the simulation \glsname{cubesat} itself, the GranaSAT~-~I. It is intended to be used as the basis to keep working in the development of an university \glsname{cubesat} in the University of Granada. As for this project, it must be able to integrate within the physical simulation platform, described in \autoref{gstation}, as well as communicating with it in the terms stated before. It must comply with \glsname{cubesat} standard~\cite{calpoly} and being designed in such a way future designs and subsystems can be integrated. Particularly, this project will deal with \textbf{On-board~Computer}~(\acrshort{OBC}), \textbf{Attitude~Determination~\&~Control~System}~(\acrshort{ADCS}) the photovoltaics module of the \textbf{Electrical Power System}~(\acrshort{EPS}) and basic \textbf{On-Board Data Handling} \acrshort{OBDH} functions, integrated in the \acrshort{OBC}. Additionally, as part of the software subsystem of the CubeSat, at least \acrshort{OBC} and \acrshort{ADCS} shall have a minimum working set of functions, which allow testing basic functions of the platform (mission simulation approach) and can be used as guide for students to implement and try new features (academic approach).

This is the largest and most complex subpart of the whole system. It involves tasks belonging a wide range of the technical spectrum, from \textbf{Mechanical Engineering} (GranaSat-I physical design and characterization and inertial orbit simulator design) to \textbf{Telecommunication Engineering} (electronic design, communication, software development) or \textbf{Aerospace Engineering }(in orbit dynamics, space representations). This fact makes this project an actual representation of what a mission design or space-related project is like: a multidisciplinary challenge which	combines work and knowledge of completely different fields in order to reach a common objective. 

These end user objectives yields to the functional requirements listed next. They contain not only those derived from the iterative analysis of the mission objectives, but also taking into account \acrshort{ECSS} applicable standards, featuring ECSS-E-ST-60-30C \cite{ESAAOCS} or \linebreak ECSS-E-ST-70-41C \cite{ESATLM} among others. Once more, the intention is keeping the design process as close as possible to a real mission design, easing future improvements of this project intended to be in orbit.

 
%%para que las tablas que están solas en una pagina no se centren verticalmente
%\makeatletter
%\setlength{\@fptop}{0pt}
%\makeatother

\newpage

\autoref{techobjsimcub} lists Technical Objectives for the Simulation \glsname{cubesat}.

\begin{table} [ht!]
\centering

\begin{tabularx}{\linewidth}{lX}
 & \multicolumn{1}{c}{\textbf{Technical Objectives}}    
\tabularnewline \specialrule{1.1pt}{1pt}{1pt}

\multicolumn{1}{c}{\textbf{Ref.}}                      & \multicolumn{1}{c}{\textbf{Primary}}                    \tabularnewline \specialrule{1.1pt}{1pt}{1pt}
SC.Obj.1                                              & To be an \glsname{engmodel} for GranaSAT-I design, which  allows testing both provided functionalities and new ones added in subsequent projects.                                                                                 \tabularnewline \midrule
SC.Obj.2                                              & To allow the complete testing possible of \acrshort{OBC} module developed, from management of the GranaSAT-I electronics, to dealing with the different functioning modes, or coping with critical situations.                                                                                                         \tabularnewline\midrule
SC.Obj.3                                              & To allow the complete testing possible of \acrshort{ADCS} module developed, from management of the GranaSAT-I maneuvering system, including motors or \glsname{magnetorquers}, detumbling mode, send sensors data to the \acrshort{OBC} dynamically, etc. etc.                                                                                                           \tabularnewline\midrule
SC.Obj.4                                              & To allow space and magnetic attitude determination, and control, if needed.                                                                                                                                  \tabularnewline\midrule
SC.Obj.5                                            & To allow photovoltaics module testing, when used with a proper \acrshort{EPS}. This module must be able to provide enough energy to the GranaSAT-I in every  necessary point of the mission/simulation.                                                                                                       \tabularnewline\midrule
SC.Obj.6                                              & To allow wired communications, when used in the simulation platform, in order to change any internal configuration or download data.                                                                                                                 \tabularnewline


\specialrule{1.1pt}{1pt}{1pt}
\multicolumn{1}{c}{\textbf{Ref.}}                      & \multicolumn{1}{c}{\textbf{Secondary}}                    \tabularnewline \specialrule{1.1pt}{1pt}{1pt}
SC.Obj.7                                            & 
To keep it as lightweight as possible, maximizing portability.                                                                          \tabularnewline \midrule

SC.Obj.8                                                   & To measure different ambient conditions, such as pressure, temperature or similar. \tabularnewline \midrule
SC.Obj.9                                                   & To allow wireless communications, when used in the simulation platform, in order emulate radio-link communication used in a real in orbit mission. \tabularnewline \midrule

SC.Obj.10                                              & To deal with the changing conditions \textit{in orbit} and act consequently.                                           \tabularnewline\midrule
SC.Obj.11                                              & To determine its angular velocity.                                                                                                               \tabularnewline
\end{tabularx}
\caption{Simulation CubeSat - Technical Objectives}
%\vspace{-1.3\baselineskip}
\label{techobjsimcub}   

\end{table}

\newpage
\autoref{funreqsimcub} shows the Functional Requirements determined.


\begin{table} [ht!]
\centering

\begin{tabularx}{\linewidth}{lX}

\multicolumn{1}{c}{\textbf{Ref.}}                      & \multicolumn{1}{c}{\textbf{Functional Requirements}}                    \tabularnewline \specialrule{1.1pt}{1pt}{1pt}
SC.FR.1                                              & The subsystem shall follow applicable interconnection standards, so \acrshort{OBC} and any other subsequent can be tested. \tabularnewline \midrule
SC.FR.2                                              & The subsystem shall have movement, magnetic and ambient conditions sensors, as well as a tachometer-like sensor.                                                                    \tabularnewline \midrule
SC.FR.3                                            & The subsystem shall have integrate \acrshort{ADCS} and lightning sensors capabilities in order to determine power source attitude or irradiation level.              \tabularnewline \midrule
SC.FR.4                                                   & To ensure that the attitude guidance and pointing specified by the mission requirements, during the mission operational phase, are met. This includes tracking of a fixed point or Sun pointing, as well as any other specific attitudes needed for system purposes, like communications for instance.    \tabularnewline \midrule
SC.FR.5                                                   & To perform the attitude measurement, estimation, guidance and control needed for the mission, autonomously.    \tabularnewline \midrule
SC.FR.6                                                   & The subsystem shall count with \acrshort{IEEE}802.11b \cite{ieee80211} connection capability, integrated with the \acrshort{OBC}.    \tabularnewline \midrule
SC.FR.7                                                   & The subsystem shall count with \acrshort{IEEE}802.3 \cite{ieee8023} (Ethernet) connection capability, integrated with the \acrshort{OBC}.    \tabularnewline \midrule
SC.FR.8                                                   & The subsystem shall count with \acrshort{USB} \cite{usborg} connection capability, integrated with the \acrshort{OBC}.    \tabularnewline \midrule
SC.FR.9                                                   & The whole system shall provide \glsname{housekeeping} telemetry to enable the verification of the nominal behaviour of sensors, actuators and on-board functionalities, on ground.    \tabularnewline \midrule

\end{tabularx}
\caption{Simulation \glsname{cubesat} - Functional Requirements}
\label{funreqsimcub}

\end{table}


% Capítulo 3: Ingeniería Inversa
\chapter{Space System Engineering}\label{chap:chapter3}

This third chapter tackles another phase related to analysis and requirements definition. It resembles the more general 'System Analysis` and will cover the whole Platform, as well as the most common \textbf{testing procedures} demanded in space according to \acrshort{ECSS}. The Platform will be analyzed following the previously defined structure, divided into Mechanical Simulation Platform, Ground Station and Simulation \glsname{cubesat}.

After deeply studying every subsystem and technologies associated to each, \textbf{formal~requirements} are defined (in contrast with the objectives and preliminary requirements defined in \autoref{chap:chapter1} and \autoref{chap:chapter2}) according to the Design Process followed in this project, described in \autoref{engprocess}. At that point, the project may go through the \textbf{System Design} process, detailed in \autoref{chap:chapter4}.

%%%%%%%%%%%%%%%%%%%%%%%%%%%%%%%%%%%%%%%%%

%%\section{Introduction to System Engineering}  \label{sec:syseng}
%
%Systems engineering is an interdisciplinary field of engineering and engineering management that focuses on how to design and manage complex systems over their life cycles \cite{wiki}.
%
%A system is an integrated set of elements to accomplish a defined objective. These elements include  hardware, software, firmware, human resources, information, techniques, facilities services, and other support elements.
%
%OJO AJUSTAR Within space sector, as seen in \autoref{missioneng}, the highest level is often called \textit{mission level}, which usually consists of one or more segments (space, ground, user...). From this perspective, requirements originate from the next upper level (the customer) and elements are procured from the next lower level (the suppliers). 
%
%Figure \autoref{fig:sysengfig} shows, according to ECSS-E-ST-10C \cite{ESASYSENG}, the boundaries of system engineering, its relationship with production, operations, product assurance and management disciplines and its internal partition into the following system engineering sub-functions:
%
%\begin{itemize}
%
	%\item{Requirement engineering, which consist on requirement analysis and validation, requirement allocation, and requirement maintenance;}
	%\item{Analysis, which is performed for the purpose of resolving requirements conflicts, decomposing and allocating requirements;}
	%\item{Design and configuration which results in a physical architecture, and its complete system functional, physical and software characteristics;}
	%\item{Verification, which objective is to demonstrate that the deliverables conform to the specified requirements;}
	%\item{System engineering integration and control, coordinating the various engineering disciplines and participants;}
%
	%\end{itemize}
	%
	%
			%\begin{figure}[h]
			%\centering
			%\includegraphics[page=1,trim={0cm 0cm 0cm 0cm},clip=true,width=150mm]{figurastfm/Chapter2/PDF/syseng.pdf}
			%\caption{System engineering, sub-functions and boundaries \cite{ESASYSENG}}
			%\label{fig:sysengfig}
			%\end{figure}
%
	%The system engineering process is in turn applied by each system engineering function of each supplier of the elements of the product decomposition.This process consists of activities to be performed by the system engineering function within each project phase according to the designated lifecycle model. The objective is to obtain a product which satisfies the customer technical requirements within pre-established objectives of cost time and quality. Through this process the system engineering function performs a multidisciplinary functional decomposition to obtain logical lower level products.	
	%
	%The functional decomposition defines, for each level of the system, the technical requirements for the procurement of subassemblies or lower level products as well as the requirements for the verification of the final characteristics of each product. The system engineering process uses the results of these lower level verification activities to build a bottom-up multi-layered evidence that the customer requirements have been met.
%
	%This project will follow a methodology as close as possible to this system engineering perspective, according to \cite{ecss} standards.

%%%%%%%%%%%%%%%%%%%%%%%%%%%%%%%%%%%%%%%%5
\section{Testing procedures} \label{testing}

CubeSats intended to be in orbit must pass a series of comprehensive testing programs in order to get the qualification to fly. Through this analysis, the different tests performed are introduced, so the constraints can be taken into account in \autoref{chap:chapter4}, and especially in future designs. In a real mission project, in which a qualification to fly is needed, it is vital to know which are the technical goals the system shall comply with as soon as possible.

This project does not intend to perform a complete testing-oriented design, because its double perspective does not allow to set space-related requirements specifically, needed to pass the demanding testing stage. Instead, the most relevant tests are introduced; they can be considered when setting the requirements and the design can stick to them as much as possible, enabling to design in complete accordance with them in the future; the aforementioned work will suppose less requirements change and their associated costs when that re-engineering process is performed.

Testing environment and requirements are given by standard ECSS-E-ST-10-03C \cite{ESAtest}. Particularly, this standard splits the process into two levels of decomposition: requirements for \textbf{space segment equipment test} and \textbf{space segment element test}. While \textit{space segment} refers to a part of a space system intended to fulfil the goals of the mission, \textit{space segment elements} are elements within a space segment.

\subsection{Space segment equipment test requirements} \label{spacesegment}

\subsubsection{General tests}
\paragraph{Humidity test}

In order to prevent the space segment of the system from the humidity effects (see more information on this in ECSS‐Q‐ST‐70‐01 \cite{ESAhumid}). This test shall be performed if space segment equipment can be exposed to humidity level above 65 \% during its life time. To qualify humidity test, the space segment equipment shall be installed in a chamber with temperature at room ambient conditions. Temperature shall be increased to +35\textdegree C and humidity to at least 95 \% over an hour. This process must be repeated during 4 different cycles according to the mentioned standard. Finally, the equipment must be visually inspected for deterioration or damage. An example of humidity chamber is shown in \autoref{tempch}.

			\begin{figure} [H] 				
				\centering
				\includegraphics[page=1,trim={0cm 0cm 0cm 0cm},clip=true,width=85mm]{figurastfm/Chapter3/Imagenes/temp_chamber.png}
				\caption{Temperature and humidity chamber \cite{tempchamber}} \label{tempch}
			\end{figure}


\paragraph{Life test}

The life test for space segment equipment qualification shall be designed to demonstrate the ability of the space segment equipment to withstand the maximum operating time and the maximum number of predicted operational cycles during the “product lifetime” by providing the required performance at the end of life. Obviously, the test must be performed under the environmental conditions expected during actual operation, including ambient, thermal, vacuum or a combination of these. 

\subsubsection{Mechanical and structural integrity tests}
\paragraph{Physical properties measurements}

Mechanical integrity is one of the most crucial aspects in every spatial system. Therefore, according to ECSS‐Q‐ST‐70‐01 \cite{ESAhumid}, there must be determined the following physical properties:

\begin{itemize} [noitemsep,topsep=0pt]
	
	\item{Dimensions and interfaces} \\
	\item{Mass} \\
	\item{Centre of gravity with respect to a given coordinate system for three mutually perpendicular axes} \\
	\item{Momentum of inertia with respect to the given coordinate system}

\end{itemize}
%\vspace{-\topsep}

Because of the intermediate character of the simulation platform designed in this project, between real simulation and training platform, this testing will be applied not only to the CubeSat itself, but to the whole physical simulation platform, as widely analyzed in \autoref{chap:chapter4}.

\paragraph{Acceleration test}

In order to ensure uniform force distribution on the space segment equipment, it is centrifuged with an arm whose length shall be at least five times the dimension of the space segment equipment measured along the arm. 

%meter imángenes como https://www.google.es/url?sa=i&source=images&cd=&cad=rja&uact=8&ved=2ahUKEwj7q4iLpI7iAhWNlhQKHe3BAF0Qjxx6BAgBEAI&url=http%3A%2F%2Fshiken.jaxa.jp%2Fen%2Ffacility13_e.html&psig=AOvVaw3h76l0b7_1W2AS1ktcZjk7&ust=1557485469670442

\paragraph{Random and sinusoidal vibration test}

Random and sinusoidal tests shall be both conducted in the launch configurations for all axes. In order to evaluate the space segment equipment integrity, a resonance search shall be performed before and after both tests. The equipment shall not suffer a shift in frequency above 5 \% and 10 \% for amplitude shift, respectively. \autoref{vibration} shows a typical vibration test equipment from \acrshort{ESA}.

			\begin{figure} [H] 				
				\centering
				\includegraphics[page=1,trim={0cm 0cm 0cm 0cm},clip=true,width=100mm]{figurastfm/Chapter3/Imagenes/vibration_test.jpg}
				\caption{Vibration test \cite{ESA}} \label{vibration}
			\end{figure}
			
\paragraph{Shock test}

The shock tests demonstrate the ability of the space segment equipment to withstand the shocks encountered during its lifetime, e.g.: fairing separation or solar arrays and antennas deployment. The equipment shall be powered during the test and the selected shock test method shall achieve the specified \textbf{ Shock Response Spectrum }with a representative transient, comparable in shape and duration to the expected in‐flight shock.

Once again, visual inspection and hardware integrity review shall be performed after test. \autoref{shock} shows a standard shock test equipment.

	\begin{figure} [H] 				
				\centering
				\includegraphics[page=1,trim={0cm 0cm 0cm 0cm},clip=true,width=100mm]{figurastfm/Chapter3/Imagenes/SHOCK_TEST.jpg}
				\caption{Shock test equipment \cite{isispace}} \label{shock}
			\end{figure}

\paragraph{Leak test}

The leak test shall demonstrate the ability of sealed or pressurized space segment equipment to conform to the leak rates stated in the specifications. They are performed only on sealed or pressurized space segment equipment, sensitive to loss of pressure or vacuum. Additionally, leak tests shall be performed prior to and following the completion of space segment equipment thermal and mechanical tests. \autoref{leaktest} shows a leaking test performed to a CubeSat at \acrshort{ESA}.

	\begin{figure} [H] 				
				\centering
				\includegraphics[page=1,trim={0cm 0cm 0cm 0cm},clip=true,width=100mm]{figurastfm/Chapter3/Imagenes/vcam.jpg}
				\caption{Leak test \cite{ESA}} \label{leaktest}
			\end{figure}

\subsubsection{Thermal-Vacuum tests}

While ECSS-E-ST-10-03C \cite{ESAtest} is the general standard for testing, ECSS-E-HB-31-03A \cite{ESAthermal1} and ECSS-Q-ST-70-04C \cite{ESAthermal2} particularly address Thermal-Vacuum testing and analysis; this kind of test subjects the system to a series of temperature cycles in vacuum conditions as present in space. The latter of them include among the deleterious effects to be anticipated during the thermal cycling test under vacuum:  

\begin{itemize}

	\item{\glsname{outgassing}}
	\item{Cracking, contraction or fracture of materials or assemblies due to sudden dimensional changes by expansion.}
	\item{Overheating of materials or assemblies due to change in convection and conductive heat transfer characteristics.}


\end{itemize}

This standard also specifies conditions for work area, such as the need to be nominally clean with minimum dust, filtered ventilation or ambient conditions relative to temperature (25 \textdegree	C $\pm$ 3) or humidity (25 \%	$\pm$ 10). A typical thermal cycling test has the following steps:

\begin{itemize} [noitemsep,topsep=0pt]
	

	\item{The initial characteristics are analyzed at (22 \textdegree C	$\pm$ 3) or humidity (55 \%	$\pm$ 10) within six hours prior to the beginning of the thermal cycling.} \\
	\item{Thermal cycling begins after a working vacuum of \SI{1e-5}{\pascal} to $\pm$ 5 \% has been reached.} \\
	\item{The chamber is thermally cycled between temperatures of ($\pm$100 \textdegree C $\pm$ 5) unless otherwise specified, at a nominal heating or cooling rate of +(10~\textdegree C~$\pm$~2) per minute.} \\
	\item{A minimum of 100 thermal cycles shall be performed on each sample.} \\
	\item{Within six hours after the completion of thermal cycling in vacuum, final inspection and testing of the sample shall be conducted at the initial ambient conditions.} \\

\end{itemize}

This test is commonly performed in a so-called \textbf{Thermal Vacuum Chamber} or \acrshort{TVAC}, as shown in \autoref{tvac}. 

\vspace{1cm}
	\begin{figure} [H] 				
				\centering
				\includegraphics[page=1,trim={0cm 0cm 0cm 0cm},clip=true,width=100mm]{figurastfm/Chapter3/Imagenes/tvac.jpg}
				\caption{Thermal Vacuum Chamber \cite{isispace}} \label{tvac}
			\end{figure}

\newpage
\glsname{GranaSAT} laboratory counts with a self-made \acrshort{TVAC} with a full \textit{in-house} design; details are available at \cite{juanma}. \autoref{tvacjuanma} shows some photos of the chamber. In the future, the testing procedures described above will be performed using that equipment, allowing to determine the accuracy of a certain system to the standard, before going through the official tests of the \textit{Fly Your Satellite!} \cite{flyyour}.

			\begin{figure}[H]
			\centering
			\subfloat[Front]{\includegraphics[page=1,trim={0cm 0cm 0cm 0cm},clip=true,width=100mm]{figurastfm/Chapter3/Imagenes/tvac2.png}}
			 \quad
			\subfloat[Side]{\includegraphics[page=1,trim={0cm 0cm 0cm 0cm},clip=true,width=100mm]{figurastfm/Chapter3/Imagenes/tvac3.png}}
			\caption{Thermal Vacuum Chamber at \glsname{GranaSAT} laboratory} \label{tvacjuanma}
\end{figure}


\newpage
\subsubsection{Electrical/RF Tests}

The following Electrical/RF Tests are crucial for an adequate functioning of the whole system and are specified in ECSS-E-ST-20-07C Rev. 1 \cite{ESAmagnetic}.


\paragraph{EMC}

Electromagnetic compatibility (\acrshort{EMC}) of a space system or equipment is the ability to function satisfactorily in its electromagnetic environment without introducing intolerable electromagnetic disturbances to anything in that environment. For acceptance stage, the space segment equipment shall be subjected to the
following tests:

\begin{itemize} [noitemsep,topsep=0pt]
	

	\item{\textbf{Bonding verification:} also called ground continuity test, consists in testing whether the ground points of a device are well connected in between each other.} \\
	\item{\textbf{Power lines isolation.}} \\
	\item \textbf{Inrush current:} consisting in the surge or momentary burst of current that flows into the device when powered on, due to the high initial currents required to charge the capacitors and inductors or transformers. These currents can be as high as 20 times the steady state currents. Even though it only lasts for about 10 ms it takes between 30 and 40 cycles for the current to stabilize to its steady state. It must be limited in order to prevent the equipment from any damage. \autoref{inrush} illustrates this phenomena. \\
	
	\begin{figure} [H] 				
				\centering
				\includegraphics[page=1,trim={0cm 0cm 0cm 0cm},clip=true,width=100mm]{figurastfm/Chapter3/Imagenes/inrush_current.jpg}
				\caption{Inrush current phenomena \cite{sunpower}} \label{inrush}
				\vspace{-0.3cm}
	\end{figure}	
	
	\item Conducted emission time \& frequency domain on power lines in the operating mode.
	
\end{itemize}
\vspace{-1cm}


\paragraph{Magnetic test} \label{magnetotest}

This test is performed in order to obtain an estimation of the \glsname{mmoment} of the device under test. A common constraint is measuring the magnetic field at distances typically more than three times the size of the device, as well as assigning to its geometric center a right‐handed orthogonal coordinate system XYZ. In order to perform this test, the device should be placed in an \textbf{earth field compensated area} providing zero‐field conditions for the intrinsic moment determination i.e., a set-up which is able to \textit{cancel} Earth's magnetic field, so it does not alter the measurement. This is usually ensured by using a \textbf{Helmholtz Cage}. \glsname{GranaSAT} Laboratory counts with one of these systems developed by one of its members throughout the years, see \cite{ramsel}. Figure\autoref{cage1} shows the developed Helmholtz Cage while Figure\autoref{cagedetail} shows the detail of the wires along the cage.

			\begin{figure}[H]
			\centering
			\subfloat[Developed Helmholtz Cage\label{cage1}]{\includegraphics[page=1,trim={0cm 0cm 0cm 0cm},clip=true,width=70mm]{figurastfm/Chapter3/Imagenes/cage_side.jpg}}
			 \quad
			\subfloat[Wires detail\label{cagedetail}]{\includegraphics[page=1,trim={0cm 0cm 0cm 0cm},clip=true,width=60mm]{figurastfm/Chapter3/Imagenes/cage_detail.jpg}}
			\caption{GranaSAT Helmholtz Cage}
			\vspace{-0.5cm}
\end{figure}

Once the setup is ready, the following test sequence must be followed:

\begin{itemize}  [noitemsep,topsep=0pt]
	
	\item{\acrshort{DUT} not operating, initial measurements on the six semi‐axes at the reference distances.} \\
	 \item Deperm or demagnetization, in order to get an earth field compensated area in the Helmhotz Cage, at a frequency of 3 Hz, maximum amplitude between \SI{4000}{ \micro\tesla} and \SI{5000}{ \micro\tesla}, successively on each XYZ axis of the \acrshort{DUT}. \\
	\item{It is applied a perm field of \SI{300}{ \micro\tesla} on each axis and measurements are repeated.}

\end{itemize}

Once again, this procedure can be tested using our own equipment, shown before. \autoref{deperm} shows graphically how this procedure shall be performed. Once the device is inside a demagnetizated area, measurements are taken at reference points.

	\begin{figure}[H]
			\centering
			\includegraphics[page=1,trim={0cm 0cm 0cm 0cm},clip=true,width=160mm]{figurastfm/Chapter3/PDF/deperm.pdf}
			\caption{Smooth deperm procedure \cite{ESAmagnetic}}
			\label{deperm}
		\end{figure}




\subsection{Space segment element test requirements}

Sometimes, a satellite is used as service module on one end and as payload module test on the other. In these cases when it is not feasible to test a space segment element as a single entity, it may be tested individually. On the other hand, sometimes testing as \textit{segment} is not possible just due to the size of the \acrshort{DUT}, which exceeds the capacities of the testing facility needed. When some of the described situations apply, instead of performing a \textbf{segment equipment testing}, an \textbf{element equipment testing} is chosen. As far as the scope of this project concerns, descriptions of \autoref{spacesegment} apply. %However, as a reference again, some interesting tables regarding this procedures are included in AÑADIR REF AL ANEXO QUE SEA, YA VEREMOS.

To finish this section, allowable tolerances and test accuracies allowed by this standard are included for the most representative tests, \autoref{tabletolerances} and \autoref{tableaccuracies}.


	\begin{figure}[H]
			\centering
			\includegraphics[page=1,trim={0cm 0cm 0cm 0cm},clip=true,width=135mm]{figurastfm/Chapter3/PDF/tabletolerances.pdf}
			\caption{Allowable tolerances \cite{ESAtest}}
			\label{tabletolerances}
		\end{figure}


\begin{figure}[H]
			\centering
			\includegraphics[page=1,trim={0cm 0cm 0cm 0cm},clip=true,width=135mm]{figurastfm/Chapter3/PDF/tableaccuracies.pdf}
			\caption{Tests accuracies \cite{ESAtest}} 			\label{tableaccuracies}
			\vspace{-1cm}
		\end{figure}

This finalizes the analysis of the most relevant testing procedures needed to qualify a design to fly, according to \acrshort{ECSS}. It must be used as basis not only for this project but especially for future designs, which shall be performed sticking to them as much as possible since the early stages of the design, taking advantage of the previous work developed here.

At this point, where the main challenges to overcome for the design have been outlined, the main subsystems that integrate this project are analyzed next, following with the formal requirements definition.

%\section{Inertial 2D Orbit Simulator (I2DOS)} este nombre en diseño it has been called \textbf{Inertial~2D~Orbit~Simulator} or \acrshort{I2DOS}

\section{Mechanical Platform Analysis} \label{mechplat}

As seen in \autoref{testing} CubeSats can be tested in a large variety of situations and circumstances. Since their outbreak, different kinds of non-professional platforms have emerged, which allow enthusiasts and university developers to test their designs before official assessment,  as seen in previous section. As commented before, this Master's Thesis proposes a planar platform to emulate in orbit dynamics in a 2D plane; the formal requirements will be treated along this section. Previously, some of the most featured kinds of platforms commented are reviewed, by way of \textit{Prior Art} in this particular subsystem.

\subsection{Air bearing}

Air bearing platforms have been used for satellite \acrshort{ADCS} testing for decades. The primary objective of air bearing tests is the faithful representation of spacecraft dynamics, so they try to get a minimal-torque environment, as in space, which is particularly difficult to duplicate.  This frictionless nature partially simulates a zero-g environment, allowing the pitch, roll, and yaw  (see section \ref{sec:eulerangles}) control systems of the satellite to work as they would in space.


It allows manipulating the system under more realistic conditions. It is one of the most popular technologies for this goal. \autoref{testbedtfm} illustrates a 3D model example of an air bearing testbed.

\begin{figure}[H]
			\centering
			\includegraphics[page=1,trim={0cm 0cm 0cm 0cm},clip=true,width=75mm]{figurastfm/Chapter3/PDF/testbed.pdf}
			\caption{Spherical air bearing testbed 3D model \cite{tfmtestbed}}
			\label{testbedtfm}
			\vspace{-1cm}
		\end{figure}
		
		Additionally, air bearings are intended to enable \acrshort{DUT} to experience some level of rotational and translational freedom. Pressurized air passes through small holes in the grounded section of the bearing and establishes a thin film that supports the weight of the moving section. These planar air-bearing systems provide one rotational and two translational degrees of freedom \cite{airbearingreview}.
		
		One of the most common types of air bearing is the \textbf{spherical}, as the one depicted in \autoref{testbedtfm}. As shown, the two sections of the bearing are portions of concentric spheres, one of which rotates on an air film bounded by the other section in three degrees of freedom. Another \glsname{GranaSAT} member developed one of these systems, see in detail in \cite{testbedgranasat}. \autoref{tbedgs} shows that implementation.

\begin{figure}[h]
			\centering
			\subfloat[3D model]{\includegraphics[page=1,trim={0cm 0cm 0cm 0cm},clip=true,width=145mm]{figurastfm/Chapter3/PDF/testbed3d.pdf}} \\ 
			\subfloat[Front]{\includegraphics[page=1,trim={0cm 0cm 0cm 0cm},clip=true,width=89mm]{figurastfm/Chapter3/PDF/testbedburgos1.pdf}}
			 \quad
			\subfloat[With \glsname{cubesat} structure]{\includegraphics[page=1,trim={0cm 0cm 0cm 0cm},clip=true,width=70.5mm]{figurastfm/Chapter3/PDF/testbedburgos2.pdf}}
			\caption{Air bearing testbed developed at \glsname{GranaSAT} \cite{testbedgranasat}} \label{tbedgs}
\end{figure}


\subsection{Inertial platforms}

Another technology to simulate orbital dynamics is an inertial platform. It must allow \glsname{cubesat} to rotate in three axes. Just as air bearings could be spherical or other type, inertia platforms design can also vary. One of the most extended is the so-called gyroscopic mechanical platform, because of its simplicity. Figure\autoref{gyroscopicplat} depicts a 3D model example. The platform is designed and constructed to be adjusted to the \glsname{cubesat}size. 


			\begin{figure}[H]
			\centering
			\subfloat[Gyroscopic mechanical inertia platform \cite{inertiaplat} \label{gyroscopicplat}]{\includegraphics[page=1,trim={0cm 0cm 0cm 0cm},clip=true,width=80mm]{figurastfm/Chapter3/PDF/inertiaplat.pdf}}
			 \quad
			\subfloat[Rotational inertia platform \cite{rotinet} \label{rotinertia}]{\includegraphics[page=1,trim={0cm 0cm 0cm 0cm},clip=true,width=80mm]{figurastfm/Chapter3/PDF/rotinertia.pdf}}
			\caption{Inertial platforms examples} 
			%\vspace{-2cm}
\end{figure}


%\begin{figure}[H]
			%\centering
			%\includegraphics[page=1,trim={0cm 0cm 0cm 0cm},clip=true,width=90mm]{figurastfm/Chapter3/PDF/inertiaplat.pdf}
			%\caption{}
			%\label{gyroscopicplat}
		%\end{figure}
		
		

Another typical inertia platform is the \textbf{rotational} one. During testing, the \glsname{cubesat} is aligned with the centre of the rotation axis of the board. Figure\autoref{rotinertia} shows another 3D model example.

%\begin{figure}[H]
			%\centering
			%\includegraphics[page=1,trim={0cm 0cm 0cm 0cm},clip=true,width=90mm]{figurastfm/Chapter3/PDF/rotinertia.pdf}
			%\caption{Rotational inertia platform \cite{rotinet}} \label{rotinertia}
		%\end{figure}


\subsection{Formal Requirements Definition}

As described when analyzing \acrshort{EDP}, after an iterative refinement of the preliminary and technical requirements stated in previous chapters, the formal requirements can be defined. These are the ones which will allow the system to comply with the needs of the mission and therefore, satisfy the objectives of all the \glsname{stakeholders}. This section is included for each subsystem. Once again, they are shown in tabular format, see \autoref{frmech}.

\begin{table}  [H]
\centering

\begin{tabularx}{\linewidth}{lX}

\multicolumn{1}{c}{\textbf{Ref.}}                      & \multicolumn{1}{c}{\textbf{Formal Requirements}}                    \tabularnewline \specialrule{1.1pt}{1pt}{1pt}
MP.FoR.1                                              & The platform shall be 10x10x10 cm dimensions and have adequate attachment capabilities, in order to secure the \glsname{cubesat}. \tabularnewline \midrule
MP.FoR.2                                              & The platform shall allow a consistent and balanced rotation movement, initiated by hand.     \tabularnewline \midrule
MP.FoR.3                                            & The platform shall weight less than 3~kg without the \glsname{cubesat} unit. \tabularnewline \midrule
MP.FoR.4                                                   & The platform shall take up a surface inferior to \SI{0.4}{m^2}.  \tabularnewline \midrule
MP.FoR.5                                                   & The platform shall have three pair of holes separated by 1.50 cm, in order to incorporate external elements. The nearest to the center shall be at least 20~cm far.    \tabularnewline \midrule
MP.FoR.6                                                   & The rotating part of the platform shall exhibit a diagonal inertia matrix.   \tabularnewline \midrule

\end{tabularx}
\caption{Mechanical Platform - Formal Requirements}
\label{frmech}

\end{table}

\newpage
\section{Ground Station Analysis}

Every mission must count with a \glsname{ground}. As described in section \ref{gstation}, in a real mission, it is in charge of communicating with the spacecraft, assuring the correct functioning of the system and so forth. As for this project, its tasks will be similar, but in the local environment designed for the simulation platform. Besides, the system must be expandable enough so it can be used in a real mission.

The concept \glsname{ground} is so wide that this project cannot go into detail as it comprises a huge amount of equipment, software, procedures, communications system and so on. Instead, formal requirements will be defined based on the expected objectives and technical requirements stated. This is one of the examples in which this Master's~Thesis must focus on its academic orientation; nevertheless, as mentioned, the designed \glsname{ground} shall suppose the basis for future improvements which allow to comply with the requirements for a real mission scenario. Additionally, a couple of relevant standards are briefly introduced next. %\autoref{forgstation} contains the formal requirements needed.

\subsection{Data-management standards in Ground Station}

The packetization format of the information exchanged between the spacecraft and the \glsname{cubesat} is also standardized. Indeed, not only the packet format is based on a standard but also the mechanism to exchange the definition of the control commands and telemetry itself. Both are addressed next.

\subsubsection{Consultative Committee for Space Data Systems} \label{ccsds}

This is a governmental organism founded in 1982 to develop standards for space data and foster interoperability. It is composed of the most important space agencies in the world, \acrshort{NASA}, \acrshort{ESA}, \acrshort{JAXA}...

\begin{figure}[H]
			\centering
			\includegraphics[page=1,trim={0cm 0cm 0cm 0cm},clip=true,width=75mm]{figurastfm/Chapter3/PDF/ccsdslogo.pdf}
			\caption{\acrshort{CCSDS} logo} 
		\end{figure}


One of its most important contributions is the so-called space packet procotol, typically known just as the organism, \acrshort{CCSDS}. It specifies a communications protocol to be used by space missions to transfer space application data over a network that involves a ground-to-space or space-to-space communications link. It is defined in \cite{ccsds1}, \cite{ccsds2} and \cite{ccsds3}. \autoref{ccsdspacket} shows an example of \acrshort{CCSDS} telemetry packet format.

\begin{figure}[H]
			\centering
			\includegraphics[page=1,trim={0cm 0cm 0cm 0cm},clip=true,width=165mm]{figurastfm/Chapter3/PDF/ccsdspacket.pdf}
			\caption{\acrshort{CCSDS} format example \cite{ccsds2}} \label{ccsdspacket}
			\vspace{-1cm}
		\end{figure}


\subsubsection{XML Telemetric and Command Exchange} \label{xtce}

XML Telemetric and Command Exchange is an XML based data exchange format and spacecraft telemetry and command meta-data \cite{wiki}. For a given mission there are a number of lifecycle phases supported by a variety of systems and organizations; telemetry and command definitions must be exchanged among all of these phases, systems, and organizations.

A typical example of this process is between the spacecraft manufacturer and spacecraft-operating agency. The first defines the telemetry and command data in a format much different than the one used in the ground segment. This creates the need for database translation, increasing customization and probability of error. Using a common exchange format streamlines the process of transferring definitions from the satellite integrator to the operations team and between ground systems supporting the same satellite. This reduces the need to develop mission-specific database import/export tools and enables the reutilization of command and telemetry database tools \cite{xtce}. \acrshort{XTCE} has also been adopted as a recommendation by the \acrshort{CCSDS}.

\subsection{Formal Requirements Definition}


\begin{table} [H]
\centering

\begin{tabularx}{\linewidth}{lX}

\multicolumn{1}{c}{\textbf{Ref.}}                      & \multicolumn{1}{c}{\textbf{Formal Requirements}}                    \tabularnewline \specialrule{1.1pt}{1pt}{1pt}
GS.FoR.1                                              & The \glsname{ground} shall have a centralized system \acrshort{CCSDS} and \acrshort{XTCE} compatible. \tabularnewline \midrule
GS.FoR.2                                              & The \glsname{ground} shall allow connectivity with standard radio stations. \tabularnewline \midrule
GS.FoR.3                                            & The \glsname{ground} shall be able to use both \acrshort{TCP} or \acrshort{UDP} for the local environment operation.  \tabularnewline \midrule
GS.FoR.4                                                   & The \glsname{ground} shall have a text-based logging system which continuously registers events, along with timing information.\tabularnewline \midrule
GS.FoR.5                                                   & The \glsname{ground} shall allow connectivity with both  \acrshort{IEEE}802.3 and  \acrshort{IEEE}802.11 standards.   \tabularnewline \midrule
GS.FoR.6                                                   & The \glsname{ground} shall count with a \acrshort{LED} lamp exhibiting an irradiance of at least 200 \SI{200}{W/m^2}.   \tabularnewline \midrule

\end{tabularx}
\caption{\glsname{ground} - Formal Requirements}
\vspace{-0.5cm}
\label{forgstation}

\end{table}



\section{Simulation CubeSat Analysis} \label{cubesatanaly}

CubeSats were introduced at \autoref{cubesatdefin}. In this section, the subsystems addressed in this Master's Thesis are profusely analyzed and just as done before, formal requirements are defined. It is important to keep in mind the big picture illustrated in \autoref{cubesatsub} when analyzing a \glsname{cubesat} as a whole system. As it was specified at that point, this Master's Thesis will focus on \acrshort{OBC}, \acrshort{ADCS} and in the power source area of the \acrshort{EPS}.

\subsection{On-board computer (OBC)}\label{sec:OBC}

On-board computer is probably the most important subsystem of a \glsname{cubesat}. It is in charge of controlling the whole system, controls I/O and coordinates the different subsystems to successfully perform the tasks needed at every moment. As seen in \ref{cubesatdefin}, sometimes \acrshort{OBC} inherently includes another subsystem such as \textbf{On-board~Data~Handling}~(\acrshort{OBDH}), while in others systems it tackles functions related to a different subsystem (e.g. \acrshort{TCC}) out of simplicity. In order to perform all these tasks, the \acrshort{OBC} can be composed of more than one processing core, which can be used to free the main one. This wide scope makes this subsystem to require a strong effort of integration with different specialists.

Additionally, sensors can also be considered part of the \acrshort{OBC} given that as a last resort it is in charge of processing the information received from them. Besides, a significant number of them are physically placed at the \acrshort{OBC}

Just as previous subsystems analyzed, \acrshort{OBC} can be subdivided into different modules, which are broken down next.

\subsubsection{Central Processing Unit}

The Central Processing Unit  can be seen as the microprocessor of a standard computer. It is typically an \textbf{embedded system} and the most complex module of the system. It works as interface with the rest of the modules and \acrshort{HID}.

It is recommendable for this module to be as flexible as possible and to allow a wide range of programming possibilities. It must be able to deal with hardware interruptions, act as \textbf{Scheduler}, manage Write/Reading operations on memory, prevent the system from failures and recover it from faulty situations... in sum it is the global coordinator of the system.

Central Processing Units can be either a single \acrshort{SOC} or a complete module \acrshort{COTS} and in aerospace sector they are usually able to have a complete \textbf{Operating System} installed, normally Linux-based. %The diagram in FIGURE illustrates some of the duties and relations with the rest of the system that this module usually has.

For all the reasons exposed, a high-level analysis of this subsystem is impossible due to the vast amount of possibilities. Therefore, although the main requirements are listed at \ref{obcreqs} just as with the rest of the sections, to go deeper on the possibilities and distinctive features of this kind of modules requires the scope to be somewhat more constrained; that task will be performed and exposed in \autoref{obcdesig}.

\subsubsection{Co-processing Programmable Core} \label{copropo}

One of the featured characteristics of the cutting-edge designs in a variety of sectors is including a reconfigurable or programmable hardware module, which can be used to act as auxiliary processor, performing a variety of tasks. \textbf{Field-Programmable Gate Arrays} (\acrshort{FPGA}) are usually used for this purpose. The task to be developed by this module is generally specified using some Hardware Description Language such as \acrshort{VHDL} or Verilog.

An \acrshort{FPGA} can be used to deal with any computable problem, featuring its speediness for some applications because of their parallel nature. Some of the common usages nowadays are aerospace, digital processing, wireless communications or image processing. All of those usages are potentially useful in a \glsname{cubesat} like the one under development; that is the reason why Co-processing Programmable Core are increasingly implemented, its versatility allows to solve different problems with a single module.

This brief analysis will focus on \acrshort{FPGA} because of their wide usage for this purpose, as mentioned before. The most common \acrshort{FPGA} architecture consists of an array of \textbf{logic blocks} and routing channels with I/O capabilities. A certain application circuit must be mapped into an \acrshort{FPGA} with the adequate resources, which may vary considerably. Generally speaking, a logic block consists of a number of logical cells which are typically composed of four input \textbf{Lookup tables} \acrshort{LUT}, a \textbf{full adder} and \textbf{D-type flip-flop}. \autoref{fpgafig} shows this structure.	Also, as depicted in \autoref{fpgafig}, a clock signal is needed as most of the electronics inside of an \acrshort{FPGA} is synchronous.
 


			\begin{figure} [H] 				
				\centering
				\includegraphics[page=1,trim={0cm 0cm 0cm 0cm},clip=true,width=85mm]{figurastfm/Chapter3/Imagenes/fpga.png}
				\caption{Typical logic cell \cite{wiki}} \label{fpgafig}
			\end{figure}
			
			Among the concepts introduced before, \textbf{Lookup tables} (\acrshort{LUT}) have special importance. They are arrays which replace runtime computation with simple array indexing which is time-saver in processing terms, as retrieving a value from memory is often faster than I/O operation \cite{wiki}. \acrshort{LUT} are in charge of storing the truth table of the any boolean function. It is a parameter sometimes used to estimate the \textit{capacity} of an \acrshort{FPGA}.
			
			Regarding the programming, the user normally provides a certain design using some Hardware Description Language, as mentioned. A netlist is generated by a procedure called \textbf{synthesis}, which will be fitted to the \acrshort{FPGA} architecture using a process called \textbf{place-and-route}. As a result of this process, the user will obtain a performance report exposed via timing analysis or simulation. Finally, the generated binary file is used to program the \acrshort{FPGA}.
			
		\acrshort{FPGA} are particularly suitable for this Co-processing module because of their versatility; the variety of implementations possible is countless; for example, one of the most featured applications is \textbf{embedding a processor} inside an \acrshort{FPGA}. It has many advantages, e.g.: specific peripherals can be chosen based on the application, mitigates obsolescence, reduces costs and allows impressive customization. For instance, there is a complete flexibility to select any combination of peripherals and controllers which can be directly connected to the processor's bus. This point allows meeting non-standard requirements, for example a \acrshort{COTS} processor with ten \acrshort{UART} may be impossible to find, but it is easy to implement in an \acrshort{FPGA} embedded processor \cite{fpgaproc}. An embedded processor implementation example can be found on \cite{fpgacalvo}. Another possible implementations using \acrshort{FPGA} are generating digital interfaces such as \acrshort{I2C} or generating \acrshort{PWM}.
	
\subsubsection{Communications}

For obvious reasons, Communications subsystem or, more specifically, Telecommand~\&~Telemetry~(\acrshort{TCC}) are vital in CubeSats. Communications must be understood not only as the ones between the \glsname{cubesat} and the \glsname{ground} but also inside the \glsname{cubesat} itself. Next, a brief analysis of the useful technologies for CubeSats is introduced.

\paragraph{Wired}

Normally, wired connections are necessary even in CubeSats intended to fly to be able to program data, perform test and so forth. In this project, because of the simulation parcel, wired communications will be specially important.

\subparagraph{Serial}

This in an historical wired interface. Serial communication is based on sequentially sending data one bit at a time. Serial ports are typically identified as such which comply with RS-232 standard. One of the main drawbacks of serial communications is their slowness.

Although it is considered to be deprecated, it is still usual in many electronics systems, as debugging port or just as a contingency way of communication. In fact, these have could its usefulness in a \glsname{cubesat}. Throughout years, serial ports have led to new derived technologies, such as \acrshort{USB}. \autoref{serial} depicts DE-9 connectors, to be used with RS-232 standard.

			\begin{figure} [H] 				
				\centering
				\includegraphics[page=1,trim={0cm 0cm 0cm 0cm},clip=true,width=65mm]{figurastfm/Chapter3/Imagenes/serial.jpeg}
				\caption{DE-9 connector \cite{wiki}} \label{serial}
				\vspace{-1cm}
			\end{figure}


\subparagraph{Universal Serial Bus}

One of the most extended interfaces in industry. Derived from serial ports, they were initially intended to standardize the connection of peripherals, and finally have largely replaced its preceding interfaces.

\acrshort{USB} has plenty of advantages, improved ease of use as it is self-configurable and hot pluggable. It is also much more faster than serial ports, reaching 5 Gb/s in the 3.0 version, one of the latest \cite{wiki}. They are also extensively used as power ports. \autoref{usbfig} shows some of the standard connectors pinout.

			\begin{figure} [H] 				
				\centering
				\includegraphics[page=1,trim={0cm 0cm 0cm 0cm},clip=true,width=80mm]{figurastfm/Chapter3/Imagenes/usb.png}
				\caption{USB standard connectors pinout \cite{wiki}} \label{usbfig}
				\vspace{-0.5cm}
			\end{figure}

Because of their versatility, \acrshort{USB} interfaces may have numerous applications on CubeSats, from programming to power ports or to downloading data.

\subparagraph{Ethernet}

First standardized in 1983 as \acrshort{IEEE}802.3 is by far the most used wired technologies in \acrshort{LAN}. Systems communicating over this interface divide streams of data into \textbf{frames} which contain source and destination addresses and allow error-checking. 

In CubeSats, Ethernet interfaces may be useful to allow communications using high-level protocols such as \acrshort{SSH}, which ease communications and configuration of the system.

			\begin{figure} [H] 				
				\centering
				\includegraphics[page=1,trim={0cm 0cm 0cm 0cm},clip=true,width=60mm]{figurastfm/Chapter3/Imagenes/ethernet.png}
				\caption{RJ45 connector \cite{ethernet}}
				\vspace{-2cm}
			\end{figure}


\paragraph{Wireless}

Wireless communications are vital in real CubeSats, for obvious reasons. In orbit CubeSats communicates with \glsname{ground} using radio-links and dedicated antennas. It is normally a whole subsystem within Communications system. It goes beyond the scope of this project analyzing that kind of implementations; extended radio-link analysis and calculation can be found on \cite{tfg}. In this section, it will be briefly introduced local area wireless communications, which can be used as base for future developments focused on real long-distance links.

\subparagraph{Radio-frequency} \label{radiofr}

Communications over the air are the only choice when dealing with space communications. There is a variety of frequencies used in radio-links, depending on distance, power and numerous additional factors. Dependingo on the frequency used, the link will have certain particularities which shall be taken into account. \autoref{tfgradio} schematically shows an example of radio path between a \glsname{ground} based on Granada and \acrshort{ISS}; it illustrates some of the aspects to be considered in a long-distance radio-frequency communication such as elevation angle or effective distance, see \autoref{tfgradio}.

			\begin{figure} [H] 				
				\centering
				\includegraphics[page=1,trim={0cm 0cm 0cm 0cm},clip=true,width=130mm]{figurastfm/Chapter3/PDF/radio_l.pdf}
				\caption{Radio-path basic elements \cite{tfg}} \label{tfgradio}
				\vspace{-0.5cm}
			\end{figure}


Although this project will not tackle this kind of communications, short-distance radio link may be considered useful for the Simulation Platform branch of the project. Contrary to the links designed for larger distances, short-distance communications over the air exhibits way lower complexity and cost due to the minor requirements in relation with power, losses constraints, among others.


\subparagraph{IEEE802.11}

Commonly known as Wi-Fi (a trademark from the Wi-Fi Alliance), it is based on \acrshort{IEEE}802.11 standard and it is the most extended wireless communication protocol for local areas. Once again, this makes it specially suitable for this Simulation Platform because of its ease of use and extended use. On the other hand, it cannot be used for space radio-links.

Since its origins, back in 1997, the standard has gone through a significant number of reviews. Nowadays the most usual versions are \acrshort{IEEE}802.11g/n/ac reaching data rates up to 1300 Mbit/s in the 5 GHz band \cite{wiki}.

As a result of the wide use of \acrshort{IEEE}802.11, the majority of the systems related to it are highly standardized (\acrshort{IC}, antennas, amplification subsystems, etc.) which reduces complexity and eases integration.

\subsubsection{Payload. Sensors.}

The \glsname{payload} is generally known as the amount of cargo capacity of an aircraft, including fuel and people. However, it may also refer to the equipment specifically intended to perform a certain mission while in orbit, for instance a camera or a \glsname{startracker}.

In CubeSats field, a common \glsname{payload} example are \textbf{sensors}: they are typically a crucial part of any mission, not only as part of the \glsname{payload} but also a necessary part for the correct functioning of the whole system. Normally they are some kind of electronic device which takes some measurement or perform a certain action depending on the inputs it receives. Some examples of sensors have already been treated before, see \ref{imus}. Another possible sensors are: barometer, thermometer, magnetometer, lightmeter, tachometer... the list is countless. Many of them will be analyzed and implemented in \autoref{chap:chapter4}. \autoref{sensores} shows some examples of them.

While this analysis considers sensors as part of the \acrshort{OBC}, given that eventually the data will be treated there, they can belong to a different subsystem such as \acrshort{ADCS}, see section \ref{imus}. 
			\vspace{-0.5cm}

			\begin{figure}[H]
			\centering
			\subfloat[\acrshort{IMU}]{\includegraphics[page=1,trim={0cm 0cm 0cm 0cm},clip=true,width=55mm]{figurastfm/Chapter3/Imagenes/MPU.png}}
			 \quad
			\subfloat[Barometer and thermometer]{\includegraphics[page=1,trim={0cm 0cm 0cm 0cm},clip=true,width=55mm]{figurastfm/Chapter3/PDF/BMP.pdf}}
			\caption{Sensor examples \cite{hobbyking}} \label{sensores}
			\vspace{-2cm}
\end{figure}	


\subsubsection{Flight Software. On-Board Data Handling} \label{flight}

In charge of all the exposed, there must be a \textbf{governor}, a coordinator which controls the hardware and deals with the different stages and situations of a mission, this is the \textbf{Flight Software} (FSW). As in a significant number of the subsections analyzed through this Master's Thesis, Flight Software is such a large field that it allows a complete project focused on it. Therefore, hereby only a basic classification is introduced; this will be an important issue at the design stage and also for future improvements of the system.

As far as this project concerns, considered flight software can be subdivided into the following.

\paragraph{Non-real-time Operating System}

Non-real-time Operating Systems is basically a general purpose \acrshort{OS} to be used with personal computers, servers, etc. The main difference with \acrshort{RTOS} is \textbf{determinism}. Non-\acrshort{RTOS} are not deterministic as tasks will not run at a certain time and for a certain time, there are no guarantees for critical tasks, exhibits high latency because of using unpredictable virtual memory, as well as \glsname{jitter}.

The previous characteristics translate into the use of \textbf{non-preemptive schedulers}. Examples of these kind of \acrshort{OS} are the widely used Windows\,\textsuperscript{\textregistered}, Debian, etc.

\paragraph{Real-time Operating System}

In contrast with the latter, \acrshort{RTOS} are completely deterministic, this is, how and when a task will run given whatever conditions defined for it to do so, it is \textbf{guaranteed}. They are intended to process data as it comes in, without buffer delays. There are a huge amount of systems nowadays which must use \acrshort{RTOS}: cars, \textbf{spacecrafts}, avionics, critical systems... In sum, \acrshort{RTOS} must be able to compute a task in a \textbf{limited or predictable amount of time} i.e., is time-bounded. However, this behaviour has nothing to do with processing speed but with a known deadline (a second, an hour or a month) and reduced \glsname{jitter}.

From a technical perspective, this implies not to use virtual memory, strict scheduling (preemptive) and avoid non-deterministic elements. \acrshort{RTOS} are usually much smaller than general purpose \acrshort{OS} in order to ease maintenance and find sources of delay. One of the most used \acrshort{RTOS} is \textbf{VxWorks} \cite{vxworks}; Linux non-preemptive kernel can also be \textbf{patched} to allow real-time behaviour.


			\begin{figure} [H] 				
				\centering
				\includegraphics[page=1,trim={0cm 0cm 0cm 0cm},clip=true,width=70mm]{figurastfm/Chapter3/Imagenes/vxworks.png}
				\vspace{0.3cm}
				\caption{VxWorks logo \cite{vxworks}} 
				\vspace{-2cm}
			\end{figure}

\subsubsection{Formal Requirements Definition} \label{obcreqs}


\begin{table} [H]
\centering

\begin{tabularx}{\linewidth}{lX}

\multicolumn{1}{c}{\textbf{Ref.}}                      & \multicolumn{1}{c}{\textbf{Formal Requirements}}                    \tabularnewline \specialrule{1.1pt}{1pt}{1pt}
OBC.FoR.1                                              & The \acrshort{OBC} shall have an \acrshort{I2C} interface which allows communication with the rest of subsystems and components. \tabularnewline \midrule
OBC.FoR.2                                              & The \acrshort{OBC} shall have an \acrshort{SPI} interface which allows communication with the rest of subsystems and components. \tabularnewline \midrule
OBC.FoR.3                                            & The \acrshort{OBC} shall have a programmable \acrshort{FPGA} of at least 2500 \acrshort{LUT}. \tabularnewline \midrule
OBC.FoR.4                                                   & The \acrshort{OBC} shall have one RJ-45 connector. \tabularnewline \midrule
OBC.FoR.5                                                   & The \acrshort{OBC} shall have at least two \acrshort{USB} Micro-B device-ports and capability to deal with four.  \tabularnewline \midrule
OBC.FoR.6                                                   & The \acrshort{OBC} shall be programmable with an external computer using a \acrshort{USB} Micro-B connector. \tabularnewline \midrule
OBC.FoR.7                                                   & The \acrshort{OBC} shall have a \acrshort{LED} system to check the correct functioning of the \acrshort{EPS}  \tabularnewline \midrule
OBC.FoR.8                                                   & The \acrshort{OBC} shall have at least two \acrshort{USB} device-ports and capability to deal with four.  \tabularnewline \midrule
OBC.FoR.9                                                   & The \acrshort{OBC} shall have a barometer with a resolution of at least 0.1 hPa.  \tabularnewline \midrule
OBC.FoR.10                                                   & The \acrshort{OBC} shall have a thermometer with a resolution of at least 0.1\textdegree~C.  \tabularnewline \midrule
OBC.FoR.11                                                   & The \acrshort{OBC} shall have at least 8 \acrshort{ADC} channels. \tabularnewline \midrule
OBC.FoR.12                                                   & The \acrshort{OBC} shall have 6 analog sun sensors, one per each face of the \glsname{cubesat}. \tabularnewline \midrule
OBC.FoR.13                                                   & The \acrshort{OBC} shall have a Real Time Clock (\acrshort{RTC}) which provides the whole system with time information. \tabularnewline \midrule
OBC.FoR.14                                                   & The \acrshort{OBC} shall have an \acrshort{SD} card holder, to store telemetry data.  \tabularnewline \midrule
OBC.FoR.15                                                   & The \acrshort{OBC} shall have an integrated \acrshort{IEEE}802.11g device which allows wireless communications. \tabularnewline \midrule
OBC.FoR.16                                                   & The \acrshort{OBC} shall implement a linux-based \acrshort{RTOS}. \tabularnewline \midrule

\end{tabularx}
\caption{\acrshort{OBC} - Formal Requirements}
\vspace{-0.5cm}
\label{forobc}

\end{table}

\newpage

\subsection{Attitude Determination and Control System (ADCS)}\label{sec:ADCS}

%mirar que son esas peliculas
%$Magnetic Controllers: B-dot, Y-spin and Cross-Product$
%Wheel Controllers: Y-Momentum and Quaternion Feedback
%Estimators: Rate Kalman Filter (RKF), TRIAD algorithm, Extended Kalman Filter (EKF)

The Attitude Determination and Control System (\glsname{ADCS}) (or \glsname{AOCS}) is the subsystem in charge of assuring that the CubeSat is correctly oriented and is able to deal with external torques and forces properly. It consists of a series of sensors and actuators, which work along with different algorithms, allowing proper pointing to the objective, like sun pointing (power) or receiving antennas pointing (communications). All these factors will be extensively treated along this and \autoref{chap:chapter4}, but first it is necessary to analyze and put into context the need of this subsystem, as well as the basic concept regarding \textbf{Spacecraft dynamics} and mathematical tools required to deal with it.

\subsubsection{Spacecraft dynamics}

Spacecraft dynamics is again an extensive topic which cannot be addressed in a few lines in this project. Nevertheless, this section will introduce some of the most important concepts used in this Master's Thesis, vital for the design stage; particularly, rotation representation and concepts such as angular momentum or \textbf{inertia~matrix} are briefly introduced. Out of simplicity, no demonstrations will be included and the analysis will be kept as simple as possible.

\paragraph{Rotation representations}\label{sec:rotationrepres}

The attitude of a three-dimensional body is most conveniently defined with a set of axes fixed to the body. This set of axes is generally a triad of orthogonal coordinates, and is normally called a \textit{body coordinate frame}. The attitude of a body is thought of as a coordinate transformation that transforms a defined set of reference coordinates into the body coordinates of the spacecraft. 

Below, the most important three-axis attitude frames are summarily exposed. For further details, see \cite{sabroff} and \cite{wertz}.

\vspace{-0.2cm}
\subparagraph{Direction cosine matrix}\label{sec:directioncosine}

The basic three-axis attitude transformation is based on the direction cosine matrix, also called \textbf{attitude matrix}. Any attitude transformation in space is actually converted to this essential form. It has the important property of mapping vectors from the reference frame to the body frame, describing the transformation from coordinate system \textit{a} to \textit{b}.

This system has no singularities, which is its main advantage, but on the other hand, it supposes propagating nine elements (three unit vectors, each with three components) defining an orthogonal right-handed triad.

\subparagraph{Euler angles}\label{sec:eulerangles}

The Euler angle rotation is defined as successive angular rotations about the three orthogonal axes of the body frame. Typically, these are defined by \textit{i}, \textit{j}, and \textit{k}, and those of the reference frame by \textit{I}, \textit{J}, and \textit{K}. There is a multitude of order combinations by which the rotation can be performed.  It is common to define the Euler~\textbf{roll~angle} ($\phi$) as a rotation about the $X$ body axis, the \textbf{pitch angle} ($\theta$) about the $Y$~body~axis, and the \textbf{yaw angle} ($\psi$) about the $Z$ body axis. However, any other definition is acceptable as long as it remains consistent with the analytical development. \autoref{fig:eulerangles} shows an illustration of this rotation representation system. 

			\begin{figure}[h]
				\centering
				\includegraphics[page=1,trim={0cm 0cm 0cm 0cm},clip=true,width=90mm]{figurastfm/Chapter3/PDF/euler_angles.pdf}
				\caption{Three-axis Euler angles around \glsname{cubesat} \cite{ESA}}
				\label{fig:eulerangles}
			\end{figure}

Euler angles are intuitive and often relevant for requirements specification, but computer implementation is not straightforward and presents numerics singularities with 90\textdegree{} rotations, a phenomena called \textbf{gimbal lock}; it consists in the loss of one degree of freedom in a three-dimensional gimbal system, when the axes of two of the three gimbals are driven into a parallel configuration, \textit{locking} the system into rotation and degenerating into a two-dimensional space \cite{wiki}. \autoref{gimbal} illustrates this problem.
%\vspace{0.5cm}
	
Besides, attitude rotations derived on the basis of Euler angles necessitate dealing with nine elements of the direction cosine matrix, and each element may include several trigonometric functions. These are some of the reasons for the wide use of quaternions, explained next.

	\begin{figure}[H]
				\centering
				\includegraphics[page=1,trim={0cm 0cm 0cm 0cm},clip=true,width=140mm]{figurastfm/Chapter3/Imagenes/gimballock.jpg}
				\caption{Gimbal lock pheonomena \cite{gimbal}}
				\label{gimbal}
				%\vspace{-2cm}
			\end{figure}


\subparagraph{Quaternions}\label{sec:quaternions}

Quaternions are a direct consequence of the properties of the direction cosine matrix. It can be shown with some linear algebra that a proper orthogonal 3x3 matrix has at least one \glsname{eigenvector} with eigenvalue of unity. That eigenvector $e_{i}$ has the same component along the body axes and along the reference frame axes. It can also be demonstrated that any attitude transformation by a series of consecutive rotation about the three orthogonal unit vectors of the coordinate system can be achieved by a single rotation about the eigenvector with unity value. The quaternion is defined as a vector as follows:

\begin{equation} 
q = q_{4} + iq_{1}+jq_{2}+zq_{3}
\end{equation}



With the unit vectors, $i, j, k$ satisfying:

\begin{equation} 
i^{2}=j^{2}=k^{2}=-1
\end{equation}


The main disadvantages of quaternions are being non-linear kinematics and not very intuitive. However, they \textbf{have no singularities} and are easy to implement in software, reaching a great efficiency in computations; while representing the attitude of a body in a reference frame by a direction cosine matrix requires knowing nine parameters $a_{ij}$, quaternions only require four $q_{i}$ parameters. Besides, the elements of the direction cosine matrix, in contrast to those of the quaternions, are trigonometric functions, which are much more cumbersome to compute. \autoref{quatplot} depicts the quaternions concept graphically.

			\begin{figure}[H]
				\centering
				\includegraphics[page=1,trim={0cm 0cm 0cm 0cm},clip=true,width=65mm]{figurastfm/Chapter3/Imagenes/quaternion.png}
				\caption{Quaternion graphical representation \cite{euclidean}}
				\label{quatplot}
			\end{figure}
			
\paragraph{Inertia matrix and angular momentum}\label{imatrix}

These are concepts constantly utilized when dealing with physical designs, and are closely related to another basic concept: the \textbf{center of mass}; given a certain distribution of mass in space, it is defined as the unique point where the weighted relative position of the distributed mass sums to zero \cite{wiki}. \autoref{centermass} illustrates where the center of mass (C) locates in a block toy.

	\begin{figure}[H]
				\centering
				\includegraphics[page=1,trim={0cm 0cm 0cm 0cm},clip=true,width=50mm]{figurastfm/Chapter3/PDF/centermass.pdf}
				\caption{Center of mass example \cite{wiki}}
				\label{centermass}
			\end{figure}
			
			On the other hand, the \textbf{moment of inertia} of a body determines the amount of \glsname{torque} needed to reach a certain angular acceleration about a rotational axes. It depends on the body's mass distribution and chosen axis; the larger the moment, the greater the \glsname{torque} needed to change the rotation rate. Depending on the number of axis of rotation, the different moments of inertia are arranged into a NxN matrix, with N that number of axis, mutually orthogonal; this is the so-called \textbf{inertia matrix}, and represents how the mass is distributed in the body. A desirable body design is that whose inertia matrix is diagonal; non-zero values are called the \textbf{principal axes of inertia}. They are always symmetric.
			
			These concepts are important in order to understand the basics of the rotational dynamics. Particularly, CubeSats attitude varies according to the fundamental equations of motion for rotational dynamics or \textbf{Euler equations}:

\begin{equation}\label{eulerrot}
\dot{H} = T-\omega\times H
\end{equation}

\autoref{eulerrot} represents the equations for the conservation of the angular momentum, denoted by \textbf{H}. Angular momentum is the rotational motion of a body that will continue unless changed by a \glsname{torque}, and it is calculated as the body's \textbf{moment of inertia} times its angular rate. It is clear from that expression that the body's angular momentum will remain constant in absence of external torques, \textbf{even if some parts of the body moves with respect to another}, for example, a reaction wheel (see \autoref{actuators}) spinning. Therefore, if that happens, the rest of the body will have to spin in the opposite direction in order to conserve the total angular momentum \cite{smad}. Further manipulations to \autoref{eulerrot} allow understanding how attitude can change due to multiple causes, closely related to concepts as \textbf{inertia matrix}, among others, described before.
			
			
			%This short introduction is enough for the reader to understand the use of these concepts along this Master's Thesis.

\subsubsection{Inertial Measurement 
Units} \label{imus}

If there is an element of particular importance in the \acrshort{ADCS}, it is the \textbf{Inertial Measurement Unit} (\acrshort{IMU}). They are electronic devices which measure gravitational acceleration, angular rate and orientation of the device, by using accelerometers, gyroscopes and magnetometers \cite{wiki}, normally with a set of these three per axis, see \autoref{fig:eulerangles}. \acrshort{IMU} are widely used in a variety of spacecrafts, including planes and satellites. However, they are also useful in everyday products such as mobiles phones or wearable devices, see \autoref{movil}.


\begin{figure}[H]
				\centering
				\includegraphics[page=1,trim={0cm 0cm 0cm 0cm},clip=true,width=60mm]{figurastfm/Chapter3/Imagenes/movil.jpg}
				\caption{Mobile phones also count with an \acrshort{IMU} \cite{sentinel}}
				\label{movil}
				\vspace{-2cm}
			\end{figure}
			
In spacecrafts field, an \acrshort{IMU} is usually part of the Inertial Navigation System, which uses its measurements to calculate attitude or angular speed. Specifically, angular speed is normally integrated to get angular position. In combination with the gravitational acceleration given by the accelerometers, attitude can be estimated using predictors such as a Kalman filter. If that estimation is used to transform acceleration measurements into a inertial reference frame and integrated once, linear velocity can be achieved, linear position if integrated twice. As it could not be any other way, under this complex electronic system basic physics works, stated at \autoref{eqmovim}.


%\begin{equation} \label{eqmovim}
%\begin{aligned}
%\text{Acceleration measurements} \rightarrow & a_\text{g}(t) \\
%&v_\text{linear}(t) = \int_{0}^{t} a_\text{g}(t) dt \\
%&r_\text{linear}(t) = \int_{0}^{t} v_\text{linear}(t) dt
%\end{aligned}
%\end{equation}

\begin{equation} \label{eqmovim}
%\begin{aligned}
a_\text{g}(t) \rightarrow v_\text{linear}(t) = \int_{0}^{t} a_\text{g}(t) dt \rightarrow r_\text{linear}(t) = \int_{0}^{t} v_\text{linear}(t) dt
%\vspace{-2cm}
%\end{aligned}
\end{equation}

\autoref{diagmov} illustrates this process. 

\begin{figure}[H]
				\centering
				\includegraphics[page=1,trim={0cm 0cm 0cm 0cm},clip=true,width=150mm]{figurastfm/Chapter3/PDF/attcalc.pdf}
				\caption{Position, velocity and attitude calculation process}
				\label{diagmov}
			\end{figure}

\newpage
\subsubsection{Control laws}

As described before, the \textbf{attitude determination} process consists in combining available sensor inputs to provide an accurate solution for the attitude state as a function of time. The term \textbf{control law} is the name normally used in \acrshort{ADCS} systems to refer to the \textbf{algorithms} which, using the data from the sensors, control the actuators available as needed in order to reach the target \textbf{attitude}. \autoref{controlaw} depicts this continuous process.

\begin{figure}[h]
				\centering
				\includegraphics[page=1,trim={0cm 0cm 0cm 0cm},clip=true,width=155mm]{figurastfm/Chapter3/PDF/controlaw.pdf}
				\caption{Control laws interaction with the rest of the \glsname{cubesat}}
				\label{controlaw}
			\end{figure}
			
			There is an enormous range of control laws, with different complexities and accuracies. Most spacecrafts use some kind of active control loop, as shown in \autoref{controlaw}. It depends on the actuator used for the attitude maneuver (see section \ref{actuators}). For systems in which spacecraft rates will be small, 3-axis control can normally be decoupled into three independent axes. Two of the most used ones are briefly introduced next.
			
			
\paragraph{Proportional-Integral-Derivative Controller} \label{pid3}

In the simplest form, each axis of the \glsname{cubesat} can be controlled by a \textbf{Proportional-Derivative controller}, with a control \glsname{torque} given by \autoref{pideq}.

\begin{equation} \label{pideq}
T_C=K_P\theta_E+K_D\omega_E
\end{equation}

Where $\theta_{\text{E}}$ is attitude error angle and $\omega_{\text{E}}$ is the attitude rate error. The most important design parameter is $K_{\text{P}}$ or proportional gain, representing the amount of control torque desired from a unit of attitude error. It determines the bandwidth of the system, closely related to the speed of response. To improve performance, sometimes an \textbf{integrator} is included, completing the Proportional-Integral-Derivative Controller or simply \acrshort{PID}. The error signals are normally measured as \textbf{Euler angles}. One of the keys for a correct functioning is \textbf{tuning}; there are different methods in the literature to get an optimal tuning depending on the requirements.

\paragraph{B-dot Controller} \label{bdotsec}

It is an specific control law dedicated to the \textbf{detumbling} stage. \glsname{tumbling} occurs right after deployment from the \acrshort{PPOD}. In that moment, \glsname{cubesat} is unusable because of the uncontrolled free spinning. Therefore, the first step before trying to control attitude using any control law such as the \acrshort{PID} seen before, must be reaching a controlled rotation speed. That is the detumbling process.

B-dot is based on the usage of \glsname{magnetorquers}, introduced in \ref{actuators}. They are used to generate a \glsname{torque} opposed to the rotation of the \glsname{cubesat}, by applying an alternating positive/negative current to the coil, which produces a certain magnetic moment. It only requires knowing the evolution of the magnetic field measured by the sensors. The functioning is simple, when the satellite is rotating around a given axis, progressively pointing in the same direction as the Earth's magnetic field, the magnetic field measured for that axis increases and so does the derivative of the magnetic field, being positive. As a response, an opposite signal to that derivative is sent through the coil, that is, a negative current. Once the magnetic field of that axis turns negative, the control signal is reversed, so the resulting torque is conserved. In sum, B-dot algorithm creates a magnetic moment in the opposite direction to the change in the magnetic field measured, expressed by \autoref{bdot}.

\begin{equation}\label{bdot}
M_\text{i}=-k_\text{i}\dot{B}_\text{i}
\end{equation}

Where $i$ is one of the axes, $M_\text{i}$ the generated magnetic moment, $B_\text{i}$ the measured magnetic moment at that axis and $k_\text{i}$ is the proportional constant calculated. 
			
\subsubsection{Actuators} \label{actuators}

CubeSats may implement a variety of actuators to be used to control attitude. The most important ones are \textbf{reaction wheels} and \textbf{magnetorquers}, both of which have been mentioned before. Next, they are briefly introduced.

\paragraph{Reaction wheels}

Reaction wheels are a particular kind of flywheel used actively in attitude control. It is normally operated at a constant rotation rate which makes the \glsname{cubesat} store a large angular momentum, tending to stabilize the satellite and allowing high pointing accuracy. They generate \glsname{torque} by turning the wheel in the opposite direction of the rotation. Reaction wheels are usually controlled using \acrshort{PWM} signals. They allow generating considerably more torque than \glsname{magnetorquers}, however, current consumption is way higher and increases the weight of the \glsname{cubesat}, as well as severely impacting the mass distribution, if it is not positioned correctly. \autoref{reactw} shows two typical \glsname{cubesat} reaction wheels.

			\begin{figure}[H] 
				\centering
				\includegraphics[width=75mm]{figurastfm/Chapter3/Imagenes/reactionw.png}
				\caption{Reaction wheels \cite{csatshop}}      		
				\label{reactw}
				%\vspace{-0.5cm}
  		\end{figure}
			
			\glsname{GranaSAT} counts with different reaction wheels designs from previous works, particularly \cite{testbedgranasat}. \autoref{3dreact} shows the manufactured reaction wheel.
			
				\begin{figure}[H] 
				\centering
				\includegraphics[width=85mm]{figurastfm/Chapter3/PDF/burgoswheel.pdf}
				\caption{Reaction wheel manufactured in \glsname{GranaSAT} laboratory \cite{testbedgranasat}}      		
				\label{3dreact}
				%\vspace{-0.5cm}
  		\end{figure}
			
\newpage
\paragraph{Magnetorquers}

Also known as torque rods, they are systems widely used to control attitude or going out of \glsname{tumbling}. They are electromagnetic coils, arranged in different ways, which create a magnetic dipole and produces a certain \glsname{torque} when interacting with Earth's magnetic field. Depending on the performance required and available area, it is built with a number of turns. They are reliable and energy-efficient but the \glsname{torque} that \glsname{magnetorquers} are able to provide is very limited. An important point is that the \glsname{torque} can be generated only perpendicularly to the Earth's magnetic field vector. As advanced in \ref{bdotsec}, they are applied an alternating current depending on the desired attitude, producing a \glsname{torque} $\tau$ given by \autoref{magnetoeq}.
		
		\begin{equation}\label{magnetoeq}
		\tau=nIA\times B
		\end{equation}
		
		Where $n$ is the number of turns, $I$ the current provided, $A$ the area of the coil and $B$ the magnetic field vector. \autoref{magnetor} shows two types of magnetorquers, the rod one, in which a copper wire is wrapped around a ferromagnetic core and an embedded coil based on the \acrshort{PCB} design.
			%\vspace{-0.5cm}

			\begin{figure}[H]
			\centering
			\subfloat[Rod format \cite{isispace}]{\includegraphics[page=1,trim={0cm 0cm 0cm 0cm},clip=true,width=85mm]{figurastfm/Chapter3/Imagenes/magnetorquer.png}}
			 \quad
			\subfloat[\glsname{GranaSAT} \acrshort{PCB} design]{\includegraphics[page=1,trim={0cm 0cm 0cm 0cm},clip=true,width=70mm]{figurastfm/Chapter3/PDF/pattern.pdf}}
			\caption{Magnetorquers} \label{magnetor}
			\vspace{-0.5cm}
\end{figure}

\newpage
\subsubsection{Formal Requirements Definition}

\begin{table} [H]
\centering

\begin{tabularx}{\linewidth}{lX}

\multicolumn{1}{c}{\textbf{Ref.}}                      & \multicolumn{1}{c}{\textbf{Formal Requirements}}                    \tabularnewline \specialrule{1.1pt}{1pt}{1pt}
ADCS.FoR.1                                              & Shall have a functional \acrshort{IMU}. \tabularnewline \midrule
ADCS.FoR.2                                              & Shall measure angular speed with an accuracy of at least 1\textdegree/s \tabularnewline \midrule
ADCS.FoR.3                                            & Shall measure gravitational acceleration with an accuracy of at least 0.1 g. \tabularnewline \midrule
ADCS.FoR.4                                                   & Shall measure magnetic field intensity with an accuracy of at least \SI{10}{\micro T}. \tabularnewline \midrule
ADCS.FoR.5                                                   & Shall implement at least a \acrshort{PID} control law.   \tabularnewline \midrule
ADCS.FoR.6                                                   & Shall count at least with a 4 cm diameter Z-axis reaction wheel.  \tabularnewline \midrule
ADCS.FoR.7                                                   & Shall include an adequate DC motor to accelerate the reaction wheel, between 3 V and 6 V supply.\tabularnewline \midrule
ADCS.FoR.8                                                  & Shall allow generating at least 5 independent \acrshort{PWM} signal with a minimum frequency of 10 Hz. \tabularnewline \midrule
ADCS.FoR.9                                                   & Shall include terminals which allow connecting the system to external magnetorquers. \tabularnewline \midrule

\end{tabularx}
\caption{\acrshort{ADCS} - Formal Requirements}
%\vspace{-2cm}
%\label{forgstation}

\end{table}


\subsection{Electrical Power System (EPS)} \label{sec:EPS}

Electrical Power System (\glsname{EPS}) is the subsystem which provides, stores, distributes, and controls spacecraft electrical power. It is the most crucial subsystem, as a lack of energy in orbit would inevitably result in the end of the mission. Some of the most typical top-level \acrshort{EPS} functions are listed below:

\begin{itemize}  [noitemsep,topsep=0pt]

	\item{Supply a continuous source of power during the mission life.} \\
	\item{Control and distribution of electrical power.} \\
	\item{Deal with both average and peak consumption requirements.} \\
	\item{Allow command and telemetry capabilities for \acrshort{EPS} status, remote control, etc.} \\
	\item{Suppress transient voltages or spikes in the bus, which may damage the system.} \\


\end{itemize}

Nowadays, the challenges around space power systems focus on maximizing efficiency and reliability while minimizing mass and costs. In this section, a functional breakdown of a typical spacecraft \acrshort{EPS} will be studied in detail in order to define technical requirements.

Once more, this project will address this part of the system from the double perspective mentioned (\autoref{perspc}). On the other hand, \autoref{fig:epsfunct} depicts the main components and functions of a typical \acrshort{EPS}.

			\begin{figure}[H] 
				\centering
				\includegraphics[width=160mm]{figurastfm/Chapter3/PDF/provisional_power_sub.pdf}
				\caption{Electrical Power System main components and functions}      		
				\label{fig:epsfunct}
  		\end{figure}

\subsubsection{Power Source}

As seen in \autoref{fig:epsfunct} power source is the first of the functions to be covered by \acrshort{EPS}. As every electrical system, spacecrafts will require enough energy in order to function properly. Typically, when designing a real space system different power sources are considered depending on the energy magnitude needed for the mission, which in turn will depend on its expected duration. 

\autoref{sourcesduration} illustrates that point, by plotting together the most usual power sources capabilities along with the expected duration of that technology. Some of the most common power sources depicted are briefly analyzed next, including some with an increasing interest in CubeSats sector in recent times.

			\begin{figure} [h] 
				\centering
				\includegraphics[page=1,trim={0cm 0cm 0cm 0cm},clip=true,width=100mm]{figurastfm/Chapter3/PDF/sourcesduration.pdf}
				\caption{Optimum energy sources for various power levels and mission durations \cite{spacecraftspower}}
				\label{sourcesduration}
			\end{figure}


\paragraph{Radio-isotope thermoelectric generator} 
\label{rtg}


Radio-isotope thermoelectric generators (\acrshort{RTG}) can be considered similar to a battery and they are based on \glsname{seebek}. They are usually used when large voltages are needed in a unmaintained term longer than common batteries or fuel cells (see section \ref{fuelcell}), chemical-based, allow; therefore, they have been used as main power source in a variety of situations, from satellites to simply provide isolated facilities with electricity.

\acrshort{RTG}, as most nuclear processes, make use of \textbf{thermocouples}, a device which is able to transform thermal energy into electrical energy, due to the \glsname{seebek} mentioned before; it makes this technology part of the so-called \textit{static sources}. Two different semiconductors are mutually connected, flowing an electric current when there is a temperature gradient (produced by decay of the radioactive source)  between the p-n junction of individual thermoelectric cells connected in a series-parallel arrangement to provide the desired DC electrical output from each converter. \autoref{thermocoup} shows an example of thermocouples.

			\begin{figure}[H] 
				\centering
				\includegraphics[width=50mm]{figurastfm/Chapter3/Imagenes/thermocouple.jpg}
				\caption{Thercouple \cite{wiki}}      		
				\label{thermocoup}
				\vspace{-2cm}
  		\end{figure}

Normally, plutonium-238 is used, providing 0.54 Watts/g \cite{rtgefec} and typically up to a few kW. \acrshort{RTG} are really inefficient, varying between 3-7 \%, and this technology is gradually falling into disuse. A notable curiosity is the numerous missions which flew to the moon using those systems\cite{rtgefec}. \autoref{apollo} shows a photography of one of these modules deployed on the Moon in Apollo 14 mission.

			\begin{figure}[H] 
				\centering
				\includegraphics[width=80mm]{figurastfm/Chapter3/Imagenes/apollo.jpg}
				\caption{\acrshort{RTG} used in Apollo 14 mission \cite{wiki}}      		
				\label{apollo}
				%\vspace{-2cm}
  		\end{figure}
\paragraph{Nuclear reactor}

Nuclear reactor is an example of \textit{dynamic source}, which uses a heat source and a heat exchanger to drive an engine in a thermodynamic power cycle. It is in charge of initiating a controlled self-sustained chain reaction which results into heat, transferred to a working fluid, which drives an energy-conversion heat engine. They are typically based on uranium-235 or plutonium-238 and one of its main advantages is avoiding thermal energy storage, as the source provides continuous heat. 

With an average efficiency of 35 \% and capability to produce up to a few hundred kW, it has been extensively used by Soviet Union and Russia, in contrast with the US.

\paragraph{Fuel cells} \label{fuelcell}
 
Fuel cells convert the chemical energy of an oxidation reaction to electricity, i.e., they are based on \textbf{REDOX reactions}, normally using hydrogen as fuel and oxygen as oxidizing agent \cite{wiki}. Contrary to solar cells (see section \ref{solarcellsexpl}) they can operate without sunlight, but must carry their own reactant supply, which allows them to produce electricity continuously for as long as fuel and oxygen are supplied.


			\begin{figure}[H] 
				\centering
				\includegraphics[width=85mm]{figurastfm/Chapter3/PDF/fuel_cell.pdf}
				\caption{Fuel cell chemical basis\cite{wiki}}      		
				\vspace{-0.5cm}
  		\end{figure}
			
			
A typical single cell produces a voltage of 0.8 VDC, however, when combined, a fuel cell unit is able to generate up to tens of kilowatts of power with efficiencies as high as 80 \%, with an average of 50-60 \% for larger currents, still way higher than the rest of the power sources analyzed. 

\paragraph{Solar Photovoltaic energy and Photoelectric effect} \label{solarcellsexpl}

Despite of the variety of power sources commonly used in spacecrafts, so far the vast majority of CubeSats missions (although there are proposals pointing to another systems, see \cite{glenn} or \cite{bristol}) have utilized \textbf{Photovoltaic energy}, mainly because of size and weight constraints, as well as a lower energy requirements than bigger spacecrafts which use any of the others power sources possibilities. Therefore, with the aim of keeping this simulation platform as real as possible, it will be the one chosen.  All the more reason, and contrary to the others, in this phase of the Engineering Design Process this technology is studied in depth in this section so the main drawbacks and constraints are outlined, in order to achieve more accurate formal requirements.	


Solar Photovoltaic energy functioning is based on \textbf{Photoelectric effect}. It was first described by Max Planck in 1900 and later in 1905 Albert Einstein went in depth, for which he was granted with the Nobel Prize. Basically, Photoelectric effect consists in the emission of electrons when light hits a material, with a variable kinetic energy depending on light frequency. The modern model which explains light behaviour was beyond the traditional conception of it as a wave, and proposed that light sometimes behaves as particles of electromagnetic energy called nowadays \textbf{photons}. According to Plank's equation:

	\begin{equation} \label{eq:plank}
			E_{\text{photon}}=\hslash v
	\end{equation}
	
In \autoref{eq:plank} $E_{\text{photon}}$ is the energy of a photon, $\hslash$ is the Planck's constant and $v$ is the frequency of the light in \SI{}{\hertz}. According to it, the energy of a photon is proportional to the frequency of the light \cite{photoelec}. Experimental tests show that if the incident light had a frequency lower than a certain frequency $v_{0}$ or \textbf{threshold frequency} no electrons were ejected regardless of the light amplitude, while for frequencies greater for $v_{0}$, they were. That value depends on the metal; \autoref{photoeffect} illustrates this point.


			\begin{figure}[H] 
				\centering
				\includegraphics[width=145mm]{figurastfm/Chapter3/Imagenes/photovolt.png}
				\caption{Threshold frequency in Photoelectric effect \cite{wiki}}      		
				\label{photoeffect}
  		\end{figure}

In the example of \autoref{photoeffect}, the red light frequency is lower than threshold frequency of the metal, so no electrons are ejected. On the other hand, green and blue lights do produce this ejection, given its higher frequency when compared to $v_{0}$. Besides, the higher energy of the blue light makes electrons be ejected with a higher \textbf{kinetic energy} than green does.

As it will be seen next, and also during \autoref{chap:chapter4}, this phenomena has a capital importance in functioning and performance of \textbf{solar cells}.

To conclude this subsection, \autoref{comptech} depicts a duration comparison of the different technologies analyzed.


			\begin{figure}[H] 
				\centering
				\includegraphics[width=105mm]{figurastfm/Chapter3/PDF/comp_tech.pdf}
				\caption{Duration comparison depending on technology \cite{mit}}      		
				\label{comptech}
				\vspace{-0.5cm}
  		\end{figure}
		

\subparagraph{Solar cells} \label{scellssect}


Solar cells are electrical devices that convert the energy of light directly into electricity due to the Photovoltaic effect previously described. They are composed of several P/N junctions monolithically connected in series. Briefly, in the n-type layer, there is an excess of electrons, and in the p-type layer, there is an excess of positively charged holes (which are vacancies due to the lack of valence electrons). Near the junction of the two layers, the electrons on one side of the junction (n-type layer) move into the holes on the other side of the junction (p-type layer). This creates an area around the junction, called the \textbf{depletion region}, in which the electrons fill the holes. \autoref{solarcells} illustrates this phenomena.

	\begin{figure}[h] 
				\centering
				\includegraphics[width=140mm]{figurastfm/Chapter3/Imagenes/solarcellphenom.jpg}
				\caption{Solar cell functioning \cite{solarcell}}      		
				\label{solarcells}
				%\vspace{-0.5cm}
  		\end{figure}

When all the holes are filled with electrons in the depletion zone, the p-type side of the depletion zone (where holes were initially present) now contains negatively charged ions, and the n-type side of the depletion zone (where electrons were present) now contains positively charged ions. The presence of these oppositely charged ions creates an internal electric field that prevents electrons in the n-type layer to fill holes in the p-type layer.

When light strikes a solar cell, electrons in the silicon are ejected, which results in the formation of 'holes' (the vacancies left behind by the escaping electrons). If this happens in the electric field, the field will move electrons to the n-type layer and holes to the p-type layer. When n-type and p-type layers are electrically connected the electrons travel from the n-type layer to the p-type layer by crossing the depletion zone and then go through the connection back to the n-type layer, creating a flow of electricity \cite{solarcell}.

			Although that is the basic behaviour, performance of solar cells varies depending on its type and material; the first noteworthy distinction is \textbf{Monocrystalline} and \textbf{Polycrystalline} Solar cells. While the first of them are composed of a single crystal which is fit into a solar panel, the second are made pouring the material into molds where it cools and solidifies. Monocrystalline are easily recognizable because of their darker colour and chopped off corners. They have a longer life compared to polycrystalline, which are also less space efficient than mono. \autoref{monopoly} shows an example of each of these solar cells.
			
			\begin{figure}[H]
			\centering
			\subfloat[Monocrystalline Solar Panel]{\includegraphics[page=1,trim={0cm 0cm 0cm 0cm},clip=true,width=50mm]{figurastfm/Chapter3/Imagenes/mono.png}}
			 \quad
			\subfloat[Polycrystalline Solar Panel]{\includegraphics[page=1,trim={0cm 0cm 0cm 0cm},clip=true,width=50mm]{figurastfm/Chapter3/Imagenes/poly.png}}
			\caption{Visual aspect solar panels comparison \cite{monopolyimg}} \label{monopoly}
			\vspace{-2cm}
\end{figure}
			
			
			Another important division is considered between \textbf{single-junction} and \textbf{multi-junction} solar cells, depending on the use of one or more materials for the p-n~junctions; the latter are composed of two or more junctions in layers one on top of each other, which allows to maximize the spectral range from which energy can be collected. This point is relevant because of the frequency dependency stated before, see \autoref{eq:plank}. \autoref{triplejunction} graphically illustrates the advantages of this technique: by using different alloys with \textbf{Gallium}, the spectral range is increased and so does efficiency. This figure is also useful to compare the spectral response of this triple-junction solar cell with Sun spectrum, by plotting AM1.5 spectrum; the Sun as irradiance source will be analyzed in \ref{irradsources}. Besides the ones mentioned, there are several more ways to classify solar cells (according to their generation, biohybrids, among others) but they are not relevant for this project.
			
				\begin{figure} [H] 
				\centering
				\includegraphics[page=1,trim={0cm 0cm 0cm 0cm},clip=true,width=120mm]{figurastfm/Chapter3/PDF/triplejunction.pdf}
				\caption{Structure of a GaInP-GaInAs-Ge solar cell. Spectral range covered \cite{phandphoto}}
				\label{triplejunction}
				\vspace{-1cm}
			\end{figure}
			
						
			
			Finally, regarding manufacturing materials, while single-junction \textbf{silicon} with efficiencies about 24.7 \% \cite{phandphoto} are the most used, later techniques make use of materials such as \textbf{GaAs}, which due to its higher light absorption coefficient is much more efficient than silicon, specially when using multi-junction panels. However, they also have higher cost, which makes them appropriate only when high efficiency is needed, e.g. in space applications \cite{wiki}. Particularly, \textbf{GaAs triple-junction} solar cells mentioned before are widely used in that sector, reaching efficiencies up to 30 \% \cite{azurspace}. Of course, those efficiencies are related to high quality solar panels, considerably diminishing when low cost manufacturing techniques are used (see IB-3 \cite{phandphoto}). 
 			 
			As a reference, \autoref{nrel} shows the historical evolution of solar cells efficiency depending on the materials used, issued by the \textbf{National Renewable Energy Laboratory} in April 2019. It is worth mentioning the higher efficiencies accomplished under laboratory conditions (up to \textbf{46 \%} using four-junction cells) when compared with the ones commercially available.  For a more realistic perspective of the solar panels market nowadays, see \autoref{efficienciesreal} \cite{phandphoto}; it also shows some parameters which will be introduced later. Intensive research and the advance in techniques and materials assures a favourable future for photovoltaic technology. 	
						\afterpage{
\begin{landscape}
\begin{figure}[H]
			\centering
			\includegraphics[page=1,trim={0cm 0cm 0cm 0cm},clip=true,width=230mm]{figurastfm/Chapter3/PDF/nrel.pdf}
			\caption{Best Research-Cell Efficiencies. NREL. \cite{nrelpage}} \label{nrel}
\end{figure}
\end{landscape}
}
			
			\begin{figure}[H]
			\centering
			\includegraphics[page=1,trim={0cm 0cm 0cm 0cm},clip=true,width=130mm]{figurastfm/Chapter3/PDF/solarcell_efficiencies.pdf}
			\caption{Typical optimum performance with the most common materials \cite{phandphoto}} \label{efficienciesreal}
\end{figure}
			

			Regardless of materials and manufacturing process, from a modelling point a view, a solar cell can be initially considered as a current source in parallel with a couple of diodes and a pair of resistors, as shown in \autoref{circuitmodelsolarcell}.
			
						
				\begin{figure} [H] 
				\centering
				\includegraphics[page=1,trim={0cm 0cm 0cm 0cm},clip=true,width=160mm]{figurastfm/Chapter3/PDF/solarcellcircuit.pdf}
				\caption{Solar cell circuit model} 		
				\label{circuitmodelsolarcell}
			\end{figure}
			
			Photons in striking light are absorbed by a semiconducting matrial, such as silicon, producing a current proportional to the amount of illumination, and adjusting the output voltage as necessary to provide that current. This is by definition the behaviour of a constant current source, so in \autoref{circuitmodelsolarcell} is represented by the current source ($I_{\text{Light}}$), which is diminished by parasitics effects given by diodes and resistances ($I_{\text{o1}},I_{\text{o2}},I_{\text{Sh}}$), which will be studied later. Let us analyze the circuit; by inspection, the current flowing out ($I_{\text{out}}$) is given by:
			
			\begin{equation} \label{eq:solarcelleqshort}
				I_{\text{out}}=I_{\text{Light}}-I_{\text{o1}}-I_{\text{o2}}-I_{\text{Sh}}
			\end{equation}
			
			Current through diodes can be substituted by \textbf{Shockley diode equation}, while $I_{\text{Sh}}$ is found using Ohm's law. Therefore, equation \autoref{eq:solarcelleqshort} yields to \autoref{eq:solarcelleq}.
			
			\begin{equation} \label{eq:solarcelleq}
		I_{\text{out}}=I_{\text{Light}}-I_{\text{o1}}(e^{qV_{\text{out}}/kT}-1)-I_{\text{o2}}(e^{qV_{\text{out}}/2kT}-1)-\frac{V_{\text{out}}+I_{\text{out}}R_{\text{S}}}{{R_{\text{Sh}}}}
		\end{equation}
			
			On the other hand, when the solar cell is short-circuited:
			
				\begin{equation} \label{ilightisc}
					V_{\text{out}}=0 \Rightarrow I_{\text{out}} \approx I_{\text{Light}}=I_{\text{SC}}
				\end{equation}
			
						
			So finally, there are four currents: $I_{\text{SC}}$ is the \textbf{short-circuit current}, $I_{\text{o1}}$ and  $I_{\text{o2}}$ are \textbf{dark saturation currents}, and $I_{\text{Sh}}$  is the \textbf{shunt current} due to ohmic losses. While the latter, as said before, is given by Ohm's law, the three first are analitycally given by complex expressions which go beyond the scope of this work (detailed demonstrations can be found at \cite{photovoltaichandbook}); however, they are not needed in order to understand the basic functioning of a solar cell operation. Briefly, the short-circuit current is the sum of the contributions from each of the three regions: \textbf{the n-type region, the depletion region, and the p-type region}. On the other hand, the dark saturation current is the current generated due to recombination in the quasi-neutral regions (apparently neutral regions where electric field is zero). Hence, $D_1$, in parallel, represents the recombination current in the quasi-neutral regions ($ I_{\text{o1}} \propto e^{qV/kT} $), while $D_2$ represents recombination current in the depletion region ($ I_{\text{o2}}~\propto~e^{qV/2kT} $); a common and reasonable assumption is to ignore the dark current due to the depletion region $ I_{\text{o2}} $. \autoref{ivmodel} shows at the front a graphical representation of \autoref{eq:solarcelleq} for typical values, while at the background the corresponding Power-Voltage curve.

\begin{figure} [h] 				
				\centering
				\includegraphics[page=1,trim={0cm 0cm 0cm 0cm},clip=true,width=120mm]{figurastfm/Chapter3/PDF/IV_model.pdf}
				\caption{Typical IV curve for standard solar cell} \label{ivmodel}
			\end{figure}
			
			At small applied voltages, diodes currents ($ I_{\text{o1}} $ and $  I_{\text{o2}}  $) are negligible and the solar cell behaves as a constant current source with an output current equivalent to the short-circuit current, $ I_{\text{SC}} $, as stated in \autoref{ilightisc}. When the applied voltage is high enough so that diodes currents (recombination current) become significant, the solar cell current drops quickly.
			
			Finally, shunt current is due to ohmic losses, modelled with a couple of resistors with a varying influence: as can also be seen in \autoref{eq:solarcelleq}, shunt resistance $R_{\text{Sh}}$ has no effect on the short-circuit current, but reduces the open-circuit voltage, $V_{\text{OC}}$. Conversely, the series resistance $R_{\text{S}}$ has no effect on the open-circuit voltage, but reduces the short-circuit current, $I_{\text{SC}}$; sources of series resistance include the metal contacts, particularly the front grid, and the transverse flow of current in the solar cell emitter to the front grid. This phenomena is illustrated in \autoref{rshuntseriescomp} which shows the behaviour of a solar cell depending on parasitics resistances; in\autoref{noseries}, shunt resistance value is varied assuming that series resistance is zero, while in\autoref{noshunt}, it is performed an analysis the other way around. Indeed, it is immediate to see how shunt resistances shift open-voltage circuit with respect to no shunt resistor, while series resistances shift short-circuit current. Notice that some extreme values have been used to illustrate the phenomena. 
			
					
			\begin{figure}
					\centering		
					\subfloat[Depending on $R_{Sh}$ $(R_{series} = 0)$\label{noseries}]{\includegraphics[page=1,trim={0cm 0cm 0cm 0cm},clip=true,width=135mm]{figurastfm/Chapter3/PDF/IV_Curves_noseries.pdf}} 
							\\[3ex]
					\subfloat[Depending on $R_{series}$ $(R_{Sh}\rightarrow\infty)$\label{noshunt}]{\includegraphics[page=1,trim={0cm 0cm 0cm 0cm},clip=true,width=135mm]{figurastfm/Chapter3/PDF/IV_Curves_noshunt.pdf}}	
					\caption{I-V Curves analysis with parasitics effects} \label{rshuntseriescomp}
			\end{figure}
			
			
 Going back to \autoref{ivmodel}, it illustrates several important \textbf{figures of merit} for solar cells, which help to understand its behaviour, some of which have already been tackled.

\begin{itemize}

	\item{\textbf{Open-circuit voltage ($V_{\text{OC}}$):} voltage $V_{\text{out}}$ across the output terminals when the cell is operated at open circuit. It is not possible to extract any power from the cell at this point.}

	\item{\textbf{Short-circuit current ($ I_{\text{SC}}$):} current $ I_{\text{out}} $ at the output when the cell is operated at short-circuit. It is not possible to extract any power from the cell at this point.}
	
	\item{\textbf{Maximum power point ($MPP (V_{\text{MP}},I_{\text{MP}})$):} point on the I-V curve where the power produced is at a maximum. For any given set of operational conditions, cells have a single operating point where the values of the current and voltage of the cell result in a maximum power output. These values correspond to a particular load resistance given by Ohm's law. This point defines a rectangle whose area is the largest for any point on the I-V curve, given by $P_{\text{MP}}=V_{\text{MP}}I_{\text{MP}}$.} The maximum power point is found by solving: \cite{photovoltaichandbook}
	

			\begin{equation} 
				\frac{\partial P}{\partial V}\Bigr|_{\substack{V=V_{\text{MP}}}}=\frac{\partial (IV)}{\partial V}\Bigr|_{\substack{V=V_{\text{MP}}}}=\Big[ I+V\frac{\partial I}{\partial V}\Big]\Bigr|_{\substack{V=V_{\text{MP}}}}=0
			\label{mppequation}
			\end{equation}
			
			The current at the maximum power point, $I_{\text{MP}}$, is then found by evaluating \autoref{eq:solarcelleq}

	\item{\textbf{Fill factor ($FF$):} measure of the \textit{squareness} of the I-V curve, given by the ratio of the areas of the two rectangles shown in figure \autoref{ivmodel}, always lower than one. Mathematically:}
	
		\begin{equation} \label{fillfactorequation}
			FF=\frac{V_{\text{MP}}I_{\text{MP}}}{V_{\text{OC}}I_{\text{SC}}}=\frac{P_{\text{MP}}}{V_{\text{OC}}I_{\text{SC}}}
	\end{equation}
	
	\item{\textbf{Packing Density:}  refers to the area of the panel which is covered with solar cells compared to that which is empty. It affects the output power of the module as well as its operating temperature; it mainly depends on the shape (round, squared...) of the solar cells used. It is usually taken to be about 0.8.	}
	
	\item{\textbf{Power conversion efficiency $(\eta)$:} relates the power obtained at the maximum power point with the incident power $P_{\text{in}}$, where the incident power is determined by the properties of the light spectrum incident upon the solar cell}:
	
	\begin{equation} 
		\eta=\frac{P_{\text{MP}}}{P_{\text{in}}}=\frac{V_{\text{MP}}I_{\text{MP}}}{P_{\text{in}}} \label{efficencysolarcell}
	\end{equation}

\end{itemize}

Obviously, the goal of a system getting energy from a solar cell, will be operating at \textbf{Maximum Power Point}. Apparently, it should not a problem: as illustrated in \autoref{ivmodel}, every solar cell will have a certain point in which power is maximized, according to \autoref{mppequation}; it may be higher or lower depending on the quality of the solar cell (i.e, parasitic effects importance, as seen before), but it is just a matter of finding it. The problem is, reality is not that simple; in a real scenario, there is a limit to how much voltage a certain panel can produce, i.e., it will act as a constant current source only as far as the connected load allows a voltage below that maximum. Now, let us suppose the illumination changes: according to the stated, for a given connected load, output current will vary, so does voltage; if for that load, the generated current tries to impose a certain voltage \textbf{above} the maximum the panel is able to produce, the \textbf{knee} of the curve is reached and generated power will turn unstable (see \autoref{ivmodel}).

 Using the concepts exposed, the latter means that \textit{MPP} will move constantly through the I-V curve, particularly depending on two factors: \textbf{irradiance} and \textbf{temperature}. While the first of them is obvious (as stated before, the more lightning striking the solar cell, the more current), the second is due to the temperature dependence of the intrinsic carrier concentration. It can be demonstrated (\cite{photovoltaichandbook}) that temperature rising yields to an increase of it, which then increases the \textbf{dark saturation currents}. Because open-circuit voltage is proportional to the reverse of them, it will decrease. This is usually \textbf{the most important performance loss.} \autoref{tempeffectimg} shows graphically this degradation in performance.

\begin{figure} [H] 				
				\centering
				\includegraphics[page=1,trim={0cm 0cm 0cm 0cm},clip=true,width=100mm]{figurastfm/Chapter3/PDF/temp_effects.pdf}
				\caption{Standard I-V curve depending on extreme temperatures \cite{spacecraftspower}} \label{tempeffectimg}
			\end{figure}

	It is immediate to see the degradation mentioned, which must be adequately taken into account when dimensioning a photovoltaic system. On the other hand, it is worth to mention the efficiency increase of a solar cell when exposed to extremely low temperatures; indeed, it is coherent with theoretical behaviour explained before: the lower the temperature, the lower dark saturation currents will be, and with them, the larger open-circuit voltage, as precisely shows \autoref{tempeffectimg}.

Another significant reason for performance degradation is the point in cell's life, that is, generated power is expected to be lower at its \textbf{End Of Life (\acrshort{EOL})} than at its \textbf{Beginning~Of~Life (\acrshort{BOL})}, as shown in \autoref{eoleffects}.

\begin{figure} [H] 				
				\centering
				\includegraphics[page=1,trim={0cm 0cm 0cm 0cm},clip=true,width=90mm]{figurastfm/Chapter3/PDF/eol_effect.pdf}
				\caption{I-V curves with cell's life \cite{spacecraftspower}} \label{eoleffects}
			\end{figure}


Finally, the last noteworthy reason of performance decrease is \textbf{radiation}. The impinging particles produce defects in the crystalline structure of the PV cells. The resulting defects degrade the voltage and current outputs of the cell. Low-energy particles create damage close to the surface, and therefore lower the open circuit voltage. On the other hand, high-energy particles penetrate deeper in the base and lower the lifetime of electron hole pairs, thus decreasing the short circuit current. \autoref{doses} plots an I-V typical curve affected by different radiation doses accumulated during its life service.

\begin{figure} [H] 				
				\centering
				\includegraphics[page=1,trim={0cm 0cm 0cm 0cm},clip=true,width=135mm]{figurastfm/Chapter3/PDF/ivradiation.pdf}
				\caption{I-V Curve degradation with radiation doses \cite{spacecraftspower}} \label{doses}
			\end{figure}

The power generation capability continues degrading as the radiation dose accumulates over the years \cite{spacecraftspower}. \autoref{lossesyears} depicts the consequences of radiation, showing that, although it oscillates, power generation have a decreasing tendency with time. 

\begin{figure} [H] 				
				\centering
				\includegraphics[page=1,trim={0cm 0cm 0cm 0cm},clip=true,width=100mm]{figurastfm/Chapter3/PDF/radiationyears.pdf}
				\caption{Losses throughout the years \cite{spacecraftspower}} \label{lossesyears}
				\vspace{-1cm}
			\end{figure}



	%\begin{figure}
					%\centering		
					%\subfloat[ \label{}]{\includegraphics[page=1,trim={0cm 0cm 0cm 0cm},clip=true,width=100mm]{figurastfm/Chapter3/PDF/radiationyears.pdf}} 
							%\\[3ex]
					%\subfloat[\label{}]{\includegraphics[page=1,trim={0cm 0cm 0cm 0cm},clip=true,width=135mm]{figurastfm/Chapter3/PDF/ivradiation.pdf}}	
					%\caption{Performance degradation due to radiation \cite{spacecraftspower}} \label{radiationeffects}
			%\end{figure}

	
			To finalize this comprehensive analysis of solar cells, there is another vital point in solar cells response, apart from the systems to optimize delivered power, intrinsic efficiency of the solar panels and so forth; this is \textbf{spectral response}. Indeed, solar cells behave differently depending on manufacturing materials, given that each one may cover different sections of the \textbf{spectrum}, affecting to efficiency and extracted power, as mentioned before. \autoref{negrous} shows spectral response of a typical solar cell manufactured with the most usual materials nowadays, some of which have been mentioned before. I must personally express my gratitude to Dr. Hamadani, from the \textbf{National Institute of Standards and Technology} in the US, for providing me with the data shown; additional information on that study can be found on \cite{hamadani}.
			
						
	\begin{figure} [h] 				
				\centering
				\includegraphics[page=1,trim={0cm 0cm 0cm 0cm},clip=true,width=120mm]{figurastfm/Chapter3/PDF/spect_material.pdf}
				\caption{Spectral response depending on solar cells material} \label{negrous}
				%\vspace{-2cm}
			\end{figure}			
			
				As stated in the legend, solar cells made of different materials are characterized: average silicon Si, high efficiency silicon, gallium arsenide GaAs, cadmium telluride CdTe, copper indium gallium selenide CIGS and organic solar cell. Although some of them have not been mentioned in this Master's Thesis because of its low utilization in space sector, they are included as general reference for photovoltaic technology. As plotted, while the organic solar cell presents the shorter range \textbf{350-630 nm} the solar cell manufactured with high efficiency silicon expands its spectral response between the \textbf{350-1150 nm} range, clearly the best out of the analyzed. It allows us to conclude that even the material which shows the best spectral response by its own, behaves worse than multi-junction architectures which in the end cover a larger portion of the spectrum.
			


	\subparagraph{Irradiance sources} \label{irradsources}
	
	As explained in previous section, spectral response of solar cells depends on the manufacturing material. In the same way, for a certain material, maximum efficiency will be accomplished only when the \textbf{whole spectral range of the solar cell} is covered by the irradiation source. This is usually not a problem when dealing with Sun spectrum (see AM1.5 series in \autoref{triplejunction}) as current technology allows a reasonable covering of it; however, it may not satisfy the simulation platform under development in this project, as the Sun is not likely to be used as irradiation source indoors. Therefore, it is worth considering whether solar cells behaviour will be the same when excited with a non-solar illumination source. Through this section, different irradiance sources, including the Sun, will be analyzed and compared with typical solar cells expected response.
	
			 \textbullet\ \textbf{Sun}

		The Sun is, for obvious reasons, the most common irradiance source when dealing with photovoltaic energy. Solar irradiance is the energy emitted by the Sun as a result of its nuclear fusion reactions in the range 250-2500 nm. The Sun radiates approximately as a \glsname{blackbody} at an effective temperature of 6000 K, with a total amount of energy \textbf{above the atmosphere} of about \SI[separate-uncertainty = true]{1366}{W/m^2 \pm{} \num{6.9}\percent} depending on the varying distance from the sun. That irradiance is distributed along a wide portion of the spectrum, including non-visible. However, a significant portion of solar radiation incident on the atmosphere is not received at the ground, being absorbed by atmospheric constituents such as \ce{H2O} vapour, \ce{CO2}, \ce{O3}, and \ce{O2}, and resulting on an average irradiance \textbf{at the Earth's surface} of about \SI{1120}{ W/m^2} \cite{sunspectral}. \autoref{spectrumsolaraltura} illustrates Sun's spectral behavior depending on the measurement point. 
		
		\begin{figure} [H] 				
				\centering
				\includegraphics[page=1,trim={0cm 0cm 0cm 0cm},clip=true,width=120mm]{figurastfm/Chapter3/PDF/Solar_spectrum_en.pdf}
				\caption{Solar spectrum above atmosphere and at surface \cite{wiki}} \label{spectrumsolaraltura}
			\end{figure}
			
			As stated, the terrestrial spectrum varies, so in order to count with a standard spectrum as representative as possible, it is commonly taken so-called \textbf{\acrshort{AM}1.5 spectrum}. Instead of a particular measured spectrum, it is calculated from the reference AM0 spectrum under representative geometric and atmospheric conditions. \acrshort{AM} stands for \textit{Air mass}; given that both absorption and scattering depend rather strongly on the path length of sunlight through the atmosphere, i.e. on the Sun’s elevation angle above the horizon. When the sun comes closer to the horizon, its light passes through more air. \acrshort{AM}1.5 refers to an spectrum taken with an air mass of 1.5 (Sun 41\textdegree{} above the horizon) with atmospheric conditions from the US standard atmosphere, representative for most middle latitudes, among others \cite{pvlight}. It will be used extensively during this project.
			
			In the same topic, solar irradiance is usually measured using devices called \textbf{pyranometers}; they translate incident irradiation into a voltage, which can be later mapped into irradiation units. Briefly, there are three different types: \textbf{thermopile pyranometers}, adequate for high accuracy, exhibiting flat spectral response on large range between 300-3000 nm; they are based on the thermocouple principle already mentioned on \ref{rtg}. Secondly, \textbf{Silicon-based pyranometers} also known as photodiode-based or photovoltaic are a more cost-effective option, with lower accuracy specially under cloudy conditions. Its main drawback is their limited spectral range, between 350-1100 nm. They are based on Photoelectric effect, also studied before in \ref{solarcellsexpl}. \autoref{comparpyranom} compares \acrshort{AM}1.5 reference with typical spectral response of each technology. 
			
				\begin{figure} [H] 				
				\centering
				\includegraphics[page=1,trim={0cm 0cm 0cm 0cm},clip=true,width=120mm]{figurastfm/Chapter3/Imagenes/solarpyranom.png}
				\caption{Solar spectrum \acrshort{AM}1.5 reference and pyranometers expected spectral range.} \label{comparpyranom}
			\end{figure}
			
				As mentioned, thermopile pyranometers exhibit a completely flat response, much wider than the rest of technologies.
			
		\textbullet\ \textbf{Xenon Arc Sun Simulator}
	

			Considering the simulation parcel of this project, it is needed to count with irradiance sources usable indoors. One of the most typically used are high power xenon lamps, as sun simulators. They are gas-discharge-based lamps which produce light when passing electricity through ionized xenon gas at high pressure \cite{wiki}. Although it depends on the lamp, life duration and some additional factors, xenon arc lamp may match closely solar spectrum, as desired. \autoref{xenonspect} illustrates the typical spectral behaviour of this kind of lamps, compared with \acrshort{AM}1.5 standard reference again; it shows a pretty matched spectrum considering \acrshort{AM}1.5 standard which confirms this technology as a really good candidate to simulate the Sun; however, as it will be analyzed later, that matching accuracy will depend on several factors (life of the lamp, quality, filters...) which typically will make it difficult to get such a good result.
			
				\begin{figure} [H] 				
				\centering
				\includegraphics[page=1,trim={0cm 0cm 0cm 0cm},clip=true,width=120mm]{figurastfm/Chapter3/Imagenes/xenon_sun.png}
				\caption{Xenon arc lamp typical spectrum compared with \acrshort{AM}1.5 reference} \label{xenonspect}
			\end{figure}
	
An example can be found at the Aerospace Engineering School of the University of Catalu\textbf{ñ}a \cite{catalan},  where it was characterized one of these simulators of the brand \textbf{Anmingli}, shown in figure\autoref{canoncat}, able to deliver up to 4 kW. Figure\autoref{xenondetail} shows the detail of the xenon arc lamp.
	
			\begin{figure}[h]
			\centering
			\subfloat[General view\label{canoncat}]{\includegraphics[page=1,trim={0cm 0cm 0cm 0cm},clip=true,width=83mm]{figurastfm/Chapter3/PDF/canon_cat.pdf}}
			 \quad
			\subfloat[Xenon lamp detail \label{xenondetail}]{\includegraphics[page=1,trim={0cm 0cm 0cm 0cm},clip=true,width=65mm]{figurastfm/Chapter3/PDF/bomb_cat.pdf}}
			\caption{Xenon arc lamp Sun simulator used \cite{catalan}} 
\end{figure}
			
			
		Using a silicon-based pyranometer, they characterized lightning distribution of the simulator. Based on that data the plot has been rebuilt and it is shown in \autoref{catorigin}. Despite the fact that the mentioned work only considered interpolated plot, original measurements have also been represented using the data available, i.e., while \autoref{catorigin} shows discrete measurements \textit{as are}, being difficult to interpret the distribution, \autoref{catinterp} shows an interpolated version of the data, which helps to understand in a more meaningful way how the simulator beam behaves. Indeed, interpolated plot \autoref{catinterp} allows determining that the beam is off-centred to the left, reaching a maximum irradiation on it of \SI{1500}{W/m^2} at a distance of 2.7~m. This is a usual problem with these kind of low-cost simulators, designed to work at larger distances, normally related to suboptimal lens design. This fact should be taken into account when used to characterize solar cells, trying to get maximum irradiation from the simulator. In the same way, it is usually a good practice to work with an irradiation value as close as possible to the expected in a real scenario; assuming \acrshort{AM}1.5 reference again, as stated before, a reasonable value to be expected at the surface is \SI{1120}{W/m^2}. In order to get that irradiation value, the distance must be increased in a factor given by \autoref{eqlight}:
	
	
	\begin{equation} 
		\frac{I_{1}}{I_{2}}=\frac{d^{2}_{2}}{d^{2}_{1}} \rightarrow d_{2} = \sqrt{d_{1}\frac{I_{1}}{I_{2}}}\label{eqlight}
	\end{equation}

	
	Effectively, recalling that the intensity of light is inversely proportional to the square of the distance, standard \acrshort{AM}1.5 irradiation value can be reached at \textbf{3.13 m}.
	
	In addition, in the same work \cite{catalan} spectral response was measured and compared with both, Sun's and a high quality simulator of the brand \textbf{MKS}, see \cite{oriel}. It allows them to conclude that their low-cost simulator matches close enough not only the more expensive one, but also Sun's spectrum. More information on this can be found on the cited work. 
	
	\afterpage{
\begin{landscape}
\vspace*{\fill}
\begin{figure}
			\centering
			\includegraphics[page=1,trim={0cm 0cm 0cm 0cm},clip=true,width=205mm]{figurastfm/Chapter3/PDF/original_catalan_memoria.pdf}
			\caption{Original Lightning distribution. Rebuilt from \cite{catalan}} \label{catorigin}
\end{figure}
\vfill
\end{landscape}
}
	
	
\afterpage{
\begin{landscape}
\vspace*{\fill}
\begin{figure}
			\centering
			\includegraphics[page=1,trim={0cm 0cm 0cm 0cm},clip=true,width=205mm]{figurastfm/Chapter3/PDF/interpolated_cat_2_7.pdf}
			\caption{Interpolated Lightning distribution. Rebuilt from \cite{catalan}} \label{catinterp}
\end{figure}
\vfill
\end{landscape}
}
	
	\newpage
	\begin{itemize} [noitemsep,topsep=0pt]
			 \item \textbf{LED Sun Simulator}
	\end{itemize}

	\acrshort{LED}-based lamps are another possibility to simulate the Sun. They are usually less similar to the Sun's spectrum but they are more power-efficient, which may make them appropriate to be used in a simulation platform like the one under development in this project. Additionally they are easier to buy and are available in smaller packages than xenon. Figure\autoref{classic} shows an example of a classic \acrshort{LED} bulb while\autoref{thbonded} depicts a \acrshort{COB} \acrshort{LED} thermally bonded. \acrshort{LED} lamps are based on \textbf{light-emitting diodes}. \acrshort{LED} lamps run on \acrshort{DC} with voltages way lower than usual \acrshort{AC} which entails the need of some electronics to convert usual \acrshort{AC} supply into an appropriate one for the \acrshort{LED}; these circuits are usually known as \textbf{\acrshort{LED} drivers}. Another important issue related to this technology is thermal management: high temperatures can lead to failure and output light reduction, so typically this kind of lamps also include heat dissipation elements or cooling fans \cite{wiki}. 
	
				\begin{figure}[H]
			\centering
			\subfloat[Classic bulb\label{classic}]{\includegraphics[page=1,trim={0cm 0cm 0cm 0cm},clip=true,width=65mm]{figurastfm/Chapter3/Imagenes/led1.jpg}}
			 \quad
			\subfloat[Thermally-bonded\label{thbonded}]{\includegraphics[page=1,trim={0cm 0cm 0cm 0cm},clip=true,width=78mm]{figurastfm/Chapter3/Imagenes/led2.jpeg}}
			\caption{\acrshort{LED} \cite{wiki}} 
		\vspace{-0.5cm}
\end{figure}	
	
	Regarding spectral behaviour and \acrshort{AM}1.5 comparison, \autoref{ledspec} shows both as well as another example of xenon arc lamp spectral response, which makes it possible to compare them at the same time.
	
			\begin{figure} [H] 				
				\centering
				\includegraphics[page=1,trim={0cm 0cm 0cm 0cm},clip=true,width=140mm]{figurastfm/Chapter3/PDF/led_Xenon_am_spec.pdf}
				\caption{\acrshort{LED} and xenon spectrum compared with \acrshort{AM}1.5 reference \cite{xenonled}} \label{ledspec}
							\vspace{-2cm}

			\end{figure}
			
			
			As previously mentioned, \acrshort{LED} spectral match differs from Sun's in a more significative way than xenon does. However, once more, this point will vary depending on each lamp and its conditions and should be adequately characterized before going into production stage. 
			
			Some research has been previously performed, studying the possibility to use a \acrshort{LED} solar simulator to characterized solar cells, for example \cite{indios} and \cite{xenonled}. The first of them concludes that their \acrshort{LED}-based simulator spectrum matches closely \acrshort{AM}1.5 for the visible spectral range. On the other hand, the second one characterizes I-V curves of certain solar cells using both \acrshort{LED} and xenon simulators; although spectral match was not perfect in any case (see \autoref{ledspec}), the author concludes that it did not affect the gross properties of the curve. Actually, he finishes featuring that \acrshort{LED} simulator behaves better than xenon, given that there were no discernable differences in the I-V response from a number of solar cell, maybe due to the extra radiation generated by the latter and therefore additional heat which, as seen in \ref{solarcellsexpl}, shifts the open circuit voltage of the I-V curve. Nevertheless, energy losses at certain parts of the spectrum are also mentioned (around 700 nm).


\subsubsection{Power Regulation and Control}

As deducible from past sections, spacecrafts usually carry an array of photovoltaic cells (i.e., a solar panel) which powers the load and supplies the whole system. Given the volatility of the power source in a real scenario (e.g., an eclipse) it is not hard to understand that energy must be somewhat stored so the system can keep working properly when the power source is not available. Regarding the simulation platform also under development, this point should also be considered to keep its accuracy. This need will be deeply analyzed in \ref{energystr} and in few words, it is satisfied using \textbf{batteries}, as shown in \autoref{fig:epsfunct} too.  

Therefore, during eclipse the battery will be in charge of powering the load. If the solar panel, the battery and the load were operated at the same constant voltage, there would be no need for any kind of power regulation. Nevertheless, as detailed in \ref{scellssect}, the solar panel output voltage is higher at the beginning of life (see \autoref{eoleffects}), and when the array is cold for several minutes after each eclipse (see \autoref{tempeffectimg}). Also the battery voltage changes. Moreover, typically the spacecraft will be composed of a variety of components with different voltage needs. Since the system is required to provide power to the load at a voltage regulated within specified limits, see \ref{testing}, a power regulation is always needed to match voltages of various power components during the entire operation time \cite{spacecraftspower}.

This power regulation is accomplished by battery charge and discharge converters, a shunt dissipator to control the bus voltage during sunlight and a controller in charge of adequately cope with the bus voltage error signal. The solar panel as power source, and the battery as load, have both their respective I-V curves, shown in \autoref{shuntcontrol}. The system can work at either of the two intersection points A or B. However, point A is inherently unstable because the load slope is lower than the source slope. Point B, on the other hand, is inherently stable. Without a shunt control, the system would operate at point B, producing a \textbf{lower power}. Nevertheless, with adequate shunts regulating the sunlight voltage, the system will pull back from point B to point C, \textbf{shunting} the excess current $I_{shunt}$ (difference between the source current at D and the load current at C) to ground, and producing more power.



\begin{figure} [H] 				
				\centering
				\includegraphics[page=1,trim={0cm 0cm 0cm 0cm},clip=true,width=88mm]{figurastfm/Chapter3/PDF/shuntcontrol.pdf}
				\caption{Stability of operating point and shunt control during sunlight \cite{spacecraftspower}} \label{shuntcontrol}
			\end{figure}


As for the controller, depending on the voltage error signal (the difference between actual bus voltage and desired), sends a control signal to one of the regulators in order to keep the bus voltage within the specified limits, allowing the battery to be charged (sunlight) or to provide energy (eclipse), with no damage for the system.

Although the latter is mainly a safety issue, power regulation is also closely related to performance. Indeed, as exposed in \ref{scellssect}, $V_{\text{out}}$ and $I_{\text{out}}$ will constantly vary, probably mismatching the \textit{MPP}. As \autoref{ivmodel} shows, for the majority of its useful curve, solar cells act as a constant current source, but when reaching \textit{MPP} boundaries (\textit{knee} of the curve), the curve has an approximately inverse exponential relationship between current and voltage. At every moment (depending on striking radiation, temperature etc.) these pair of values will imply a different load $R=V_{\text{out}}/I_{\text{out}}$ given by Ohm's law;  for an external device to draw maximum power from the solar cell, it	should \textit{see} a load with such a resistance equal to the inverse of this value; other way, and as \autoref{efficencysolarcell} states, \textbf{efficiency will decrease} and could eventually reach one of the \textit{No power} points. Thus by \textbf{varying the impedance seen by the panel}, the operating point can be moved towards the \textit{MPP}.

In sum, the desirable operation point varies considerably and this is not only an efficiency problem by itself, but it is also a problem for the following element in the power chain: \textbf{the batteries}. The following example will illustrate the problem: let us assume a set-up where the solar panel is directly connected to a standard 3.7 V Li-on battery. According to the manufacturer, the solar cell is able to provide a power of 1 W at the \textit{MPP}~(5~V,~0.2~A); however, given that it is directly connected to a 3.7 V source (the battery, whose voltage is taken as constant) the output voltage of the solar cell is fixed to that value, therefore operating at (3.7 V, 0.2 A) and producing 0.74 W in the best scenario, implying an \textbf{unexpected loss of 26 \%} (actually, it would be even worse given that the battery voltage is not constant but it will decrease along with its available energy, see section \ref{energystr}). This is another reason why an adequate power regulation is highly desirable.

%Therefore, it would be highly desirable some kind of system which tries to set the operating point as close as possible to the \textit{MPP} given by the load each moment, especially considering the pretty low efficiency described. Both scenarios are considered and analyzed in the following sections.

The photovoltaic power system, therefore, primarily consists of a solar array, a rechargeable battery, and a power regulator which regulates power flow between various components to control the bus voltage. This section deals with the power regulation stage, particularly the different architectures normally used and its efficiency.


	\paragraph{Direct Energy Transfer} \label{det}
	
	The set-up described before, in which the solar panel is directly connected to a battery, is known as \textbf{Direct Energy Transfer} (\acrshort{DET}) architecture. In this case, the power bus is said to be \textit{dominated} by the battery voltage, which implies that the solar panel must operate at the same voltage of the power bus, potentially not delivering the full power it is capable of at all times. Since the solar array is designed to never exceed a voltage past the \textit{MPP}, it will reach its full power power producing capability only when the battery is at its highest voltage, which occurs when the battery is completely charged. However, when the battery is at its minimum voltage, at the beginning of its charge cycle, the solar array will operate well below the \textit{MPP}. This phenomena is graphically illustrated in\autoref{ivexample} which shows the voltage shift produced when the \textit{MPP} is displaced due to mismatch load imposed by the battery. Figure\autoref{pexample} depicts the power losses due to that issue.

This architecture can be subdivided into two classes which are briefly analyzed next.


	\begin{figure} 
					\centering		
					\subfloat[I-V curves in open-circuit and loaded\label{ivexample}]{\includegraphics[page=1,trim={0cm 0cm 0cm 0cm},clip=true,width=140mm]{figurastfm/Chapter3/PDF/IV_Curves_example_lat.pdf}} 
					\\[3ex]	
					\subfloat[Power curves in open-circuit and loaded\label{pexample}]{\includegraphics[page=1,trim={0cm 0cm 0cm 0cm},clip=true,width=140mm]{figurastfm/Chapter3/PDF/P_Curves_example_lat.pdf}}	
					\caption{Solar panel performance directly connected to a battery} \label{compexample}
			\end{figure}

	\subparagraph{Sunlight regulated bus} \label{sunlightregulated}
	
	This sub-architecture is normally used when the objective is to minimize complexity; to achieve that power from both sources available --- the solar panel output and the battery~--- directly to the load. An in-depth analysis of this architecture goes beyond the scope of this project, however typical sun regulated bus architecture is depicted in \autoref{sunregul} so minimal references can be made. 
	
	\begin{figure} [H] 				
				\centering
				\includegraphics[page=1,trim={0cm 0cm 0cm 0cm},clip=true,width=105mm]{figurastfm/Chapter3/PDF/sunregul.pdf}
				\caption{Sun regulated \acrshort{DET} architecture \cite{spacecraftspower}} \label{sunregul}
			\end{figure}
	
	In this architecture, the bus is said to be 'regulated' by \textbf{shunt control} during sunlight which is accomplished by a battery charge regulator to control charge rate. On the other hand, during eclipse, the battery charges directly to the bus through the \textit{battery discharge diode}, $d$ in \autoref{sunregul}; this diode only allows discharge from the battery, blocking any uncontrolled charge current coming to the battery, leaving this task to the charge regulator, and disconnecting the battery from the bus during sunlight. Without this kind of regulation, the solar panel output voltage would settle at operating point B in \autoref{shuntcontrol}, which would be unsafe at some points of operation, for example at \acrshort{BOL} and after eclipse (the solar panel would be cold and produce a higher output voltage, as seen in \autoref{tempeffectimg}).
	
	It is simple and reliable but implies variations in bus voltage up to $\pm25 \%$ around nominal value \cite{spacecraftspower}. This architecture finds application mostly in relatively low power needs, such as CubeSats, for example.
	
	\subparagraph{Fully regulated bus} \label{fullyregulated}

	This sub-architecture is commonly known as \textit{regulated bus} and it is mainly characterized by a controlled bus voltage within a few percent during orbit period. Contrary to sunlight regulated, it requires a battery discharge converter which is expensive. On the other hand, it allows great flexibility in battery choice.
	
	It is typically used when the load has high power requirements or the spacecraft loads require a lower variation around bus voltage.


	\paragraph{Maximum Power Point Tracking} \label{mppt}
	
	Section \ref{scellssect} clearly exposed the advantages of a mechanism able to get the system to the optimum performance point continuously. It can be accomplished with the \textbf{Maximum Power Point Tracking} (\acrshort{MPPT}) architecture, which makes use of switching regulation between the solar panel and the load, as \autoref{mpptarc} shows.
	
		\begin{figure} [H] 				
				\centering
				\includegraphics[page=1,trim={0cm 0cm 0cm 0cm},clip=true,width=105mm]{figurastfm/Chapter3/PDF/mpptarc.pdf}
				\caption{Maximum Power Point Tracking architecture \cite{spacecraftspower}} \label{mpptarc}
			\end{figure}
	
	
	The peak power tracker senses the \textit{MPP} and is in charge of keeping the series regulator input voltage at the maximum power producing level; then the output voltage is stepped down to the load voltage,  i.e. the battery, by varying the duty cycle as needed.
	
	Noteworthy advantages of this architecture are that it maximizes the solar panel output power at every moment without neither shunt nor battery charge regulator. As main drawback, it may result in low efficiency due to power losses in the tracking process itself.
	
		\autoref{proscons} summarizes pros and cons of the architectures analyzed. 

	
			\begin{figure} [h] 				
				\centering
				\includegraphics[page=1,trim={0cm 0cm 0cm 0cm},clip=true,width=135mm]{figurastfm/Chapter3/PDF/proscons_mppt.pdf}
				\caption{Pros and cons of the analyzed architectures \cite{spacecraftspower}} \label{proscons}
			\end{figure}
			
				
\subsubsection{Energy storage. Batteries.} \label{energystr}


In an energy-limited scenario as in orbit, energy storage is a vital issue. For example, when in eclipse, as there is no energy input available, in order for the system to keep working, there must be some kind of reserve; also when the demand exceeds the power generation at
any time. This storage is usually accomplished by using \textbf{batteries} which store energy in an electrochemical form.

As in previous topics, some battery basic concepts are introduced before going into detail.

\begin{itemize}

	\item \textbf{Energy density:} electrical energy per unit mass (Wh/kg) available when fully
discharged from a fully charged state at a given rate and temperature.
	\item \textbf{Roundtrip efficiency:} also called `Wh efficiency', is the ratio of energy provided between full charge and the following full discharge at a given rate and temperature.
	\item \textbf{Depth of discharge (\acrshort{DoD}):} amount of energy taken out of the battery per cycle. Looking in the opposite direction, the \acrshort{SoC} (state of charge) is 100 - \%\acrshort{DoD}.
	\item \textbf{Nominal capacity (Ah):} typically stated by the manufacturer. This capacity may or may not represent the amount of Ah available upon discharge down to a technology dependent \acrshort{EOL} voltage. 
	\item \textbf{Dimensional abuse:} it is defined as the charge the cell can deliver at room temperature until it reaches a cut-off voltage of about 2/3 of the fully charged voltage. Charge and discharge currents (A) are expressed as a multiple or fraction of the nominal capacity which is called `C' and depends on physical size. (e.g. 2C, C/2, C/100...). It is widely used and its practical meaning is the possibility to deliver C amperes for 1 h or C/n amperes for n hours.
	\item \textbf{Cycle life:} defined as the number of charge/discharge (C/D) cycles the
battery can deliver while maintaining the cut-off voltage.
	

\end{itemize}

Batteries are made of numerous electrochemical cells assembled in a series–parallel circuit combination to obtain the required voltage and current. It has positive and negative electrode plates with insulating separator and a chemical electrolyte in-between. The two electrode plates
are connected to two external terminals mounted on the casing. The cell stores electrochemical energy at a low electrical potential. The cell voltage depends solely on the electrochemistry, and not on the physical size. \autoref{emfcell} shows the typical electrochemical cell construction \cite{spacecraftspower}
			\begin{figure} [H] 				
				\centering
				\includegraphics[page=1,trim={0cm 0cm 0cm 0cm},clip=true,width=105mm]{figurastfm/Chapter3/PDF/emfcells.pdf}
				\caption{Typical electrochemical cell construction \cite{spacecraftspower}} \label{emfcell}
				\vspace{-0.5cm}
			\end{figure}

Batteries cells can be thought of as voltage sources with small internal resistance which respectively decreases and increases linearly with the nominal capacity discharge. Ignoring parasitics, it is as a \acrshort{DC} source representing the \textbf{electromotive force} (\acrshort{EMF}); this is the typical term used when a voltage is generated by a battery and represents the \textbf{energy} per unit charge which has been made available by the generating mechanism, in this case, the electrochemical cell. \acrshort{EMF} depends on the \acrshort{SoC} and ageing. On the other hand, the cell resistance in series can be compared with the traditional output series resistance in power sources; besides \acrshort{SoC} and ageing, it also depends on temperature. Next, the most used batteries technologies are briefly introduced.

\paragraph{Nickel Hydrogen}

\ce{NiH2} batteries have been the most widely used in space sector during the last 25 years, which makes them highly reliable. They combine some of the best characteristics from other technologies, such as NiCd or the fuel cell (see \ref{fuelcell}) and can tolerate some over-charge or over-discharge without damage; it also has a greater charge/discharge cycle life and low internal resistances than older technologies such as NiCd and do not exhibit a noticeable \glsname{memeffect}. However, it has a high self-discharge rate, around 0.5 \%, and also high loss of capacity. In the typical configuration, each cell develops 1.25~V. \autoref{nih} shows a common \ce{NiH2} batteries stack.

				\begin{figure} [H] 				
				\centering
				\includegraphics[page=1,trim={0cm 0cm 0cm 0cm},clip=true,width=80mm]{figurastfm/Chapter3/Imagenes/nih.jpg}
				\caption{\ce{NiH2} batteries stack \cite{upsnih}} \label{nih}
			\end{figure}
	
\paragraph{Lithium-Ion}	

With a great expansion throughout the last years, Li-ion technology exhibits pretty higher energy density than \ce{NiH2}. It is corrosion-free but lithium is highly reactive and must be stabilized. Li-ion has a cut-off voltage of 2.7 V, average discharge voltages of 3.5 V and end-of-charge voltage of 4.2 V. The average discharge voltage of 3.5 V contrasts with the 1.25 V for \ce{NiCd} and \ce{NiH2} requires about one third of the cells in series for a given battery voltage, allowing a smaller assembly and test costs \cite{spacecraftspower}. Another significant advantage is its capacity to deliver peak power without negative consequences on life duration.

As for the negative side, Li-ion cells are sensitive to over-charge and over-discharge and exhibit low performance at low temperature, due to associated high internal resistance under those conditions. \autoref{lion} depicts a common Li-ion battery.

			\begin{figure} [H] 				
				\centering
				\includegraphics[page=1,trim={0cm 0cm 0cm 0cm},clip=true,width=75mm]{figurastfm/Chapter3/Imagenes/lion.jpg}
				\caption{Li-ion battery \cite{sparkfun}} \label{lion}
				\vspace{-2cm}
			\end{figure}
\newpage
\paragraph{Lithium Polymer}

This short comparison finalizes with lithium polymer batteries. These batteries are increasingly common these days, because of their wide use in drones and radio-controlled devices. Their electrochemical basis are similar to that of Li-ion but they are more delicate and have the same general problems of Li-ion e.g., over-charge, over-discharge, over-temperature, short circuit, crush and nail penetration may all result in a catastrophic failure.\autoref{lipo} shows a standard LiPo battery while\autoref{lionhinchada} illustrates the consequences of an incorrect usage.

			\begin{figure}[H]
			\centering
			\subfloat[Standard LiPo battery \cite{hobbyking}\label{lipo}]{\includegraphics[page=1,trim={0cm 0cm 0cm 0cm},clip=true,width=65mm]{figurastfm/Chapter3/Imagenes/lipo.jpg}}
			 \quad
			\subfloat[Expanded Lithium-based battery \cite{wiki}\label{lionhinchada}]{\includegraphics[page=1,trim={0cm 0cm 0cm 0cm},clip=true,width=65mm]{figurastfm/Chapter3/Imagenes/lipohin.jpg}}
			\caption{LiPo batteries} 
\end{figure}

\subsubsection{Formal Requirements Definition}
			
\begin{table} [H]
\centering

\begin{tabularx}{\linewidth}{lX}

\multicolumn{1}{c}{\textbf{Ref.}}                      & \multicolumn{1}{c}{\textbf{Formal Requirements}}                    \tabularnewline \specialrule{1.1pt}{1pt}{1pt}
EPS.FoR.1                                              & Solar cells shall produce at least 1 W of power. \tabularnewline \midrule
EPS.FoR.2                                              & Solar cells shall be silicon monocrystalline.\tabularnewline \midrule
EPS.FoR.3                                            & Batteries shall be Li-ion and have a capacity of at least 1500 mAh.  \tabularnewline \midrule
EPS.FoR.4                                                   & The \acrshort{EPS} shall have a contingency system to provide 3.3 V in absence of the main regulation subsystem.\tabularnewline \midrule
EPS.FoR.5                                                   & The \acrshort{EPS} shall have \acrshort{USB} connectors which allow charging batteries from external supply.   \tabularnewline \midrule

\end{tabularx}
\caption{\acrshort{EPS} - Formal Requirements}
\vspace{-0.5cm}
%\label{forgstation}

\end{table}

%\newpage
%\subsection{On-board computer (OBC)}\label{sec:OBC}
%
%On-board computer is probably the most important subsystem of a \glsname{cubesat}. It is in charge of controlling the whole system, controls I/O and coordinates the different subsystems to successfully perform the tasks needed at every moment. As seen in \ref{cubesatdefin}, sometimes \acrshort{OBC} inherently includes another subsystem such as \textbf{On-board~Data~Handling}~(\acrshort{OBDH}), while in others systems it tackles functions related with a different subsystem (e.g. \acrshort{TCC}) out of simplicity. In order to perform all these tasks, the \acrshort{OBC} can be composed of more than one processing core, which can be used to free the main one. This wide scope makes this subsystem to require a strong effort of integration with different specialists.
%
%Additionally, sensors can also be considered part of the \acrshort{OBC} given that as a last resort it is in charge of processing the information received from them. Besides, a significant number of them are physically placed at the \acrshort{OBC}
%
%Just as previous subsystems analyzed, \acrshort{OBC} can be subdivided into different modules, which are breakdown next.
%
%\subsubsection{Central Processing Unit}
%
%The Central Processing Unit  can be seen as the microprocessor of a standard computer. It is typically an \textbf{embedded system} and the most complex module of the system. It work as interface with the rest of the modules and \acrshort{HID}.
%
%It is recommendable for this module to be as flexible as possible and to allow a wide range of programming possibilities. It must be able to deal with hardware interruptions, act as \textbf{Scheduler}, manage Write/Reading operations on memory, prevent the system from failures and recover it from faulty situations... in sum it is the global coordinator of the system.
%
%Central Processing Units can be either a single \acrshort{SOC} or a complete module \acrshort{COTS} and in aerospace sector they are usually able to have a complete \textbf{Operative System} installed, normally Linux-based. %The diagram in FIGURE illustrates some of the duties and relations with the rest of the system that this module usually has.
%
%For all the reasons exposed, a high-level analysis of this subsystem is impossible due to the vast amount of possibilities. Therefore, although the main requirements are listed at \ref{obcreqs} just as with the rest of the sections, to go deeper on the possibilities and distinctive features of this kind of modules requires the scope to be somewhat more constrained; that task will be performed and exposed in section \autoref{obcdesig}.
%
%\subsubsection{Co-processing Programmable Core}
%
%One of the featured characteristics of the cutting-edge designs in a variety of sectors is including a reconfigurable or programmable hardware module, which can be used to act as auxiliary processor, performing a variety of tasks. \textbf{Field-Programmable Gate Arrays} (\acrshort{FPGA}) are usually used for this purpose. The task to be developed by this module is generally specified using some Hardware Description Language such as \acrshort{VHDL} or Verilog.
%
%An \acrshort{FPGA} can be used to deal with any computable problem, featuring its speediness for some applications because of their parallel nature. Some of the common usages nowadays are aerospace, digital processing, wireless communications or image processing. All of those usages are potentially useful in a \glsname{cubesat} like the one under development; that is the reason why Co-processing Programmable Core are increasingly implemented, its versatility allows to solve different problems with a single module.
%
%This brief analysis will focus on \acrshort{FPGA} because of their wide usage for this purpose, as mentioned before. The most common \acrshort{FPGA} architecture consists of an array of \textbf{logic blocks} and routing channels with I/O capabilities. A certain application circuit must be mapped into an \acrshort{FPGA} with the adequate resources, which may vary considerable. Generally speaking, a logic block consists of a number of logical cells which are typically composed of four input \textbf{Lookup tables} \acrshort{LUT}, a \textbf{full adder} and \textbf{D-type flip-flop}. \autoref{fpgafig} shows this structure.	Also, as depicted in \autoref{fpgafig}, a clock signal is needed as most of the electronics inside of an \acrshort{FPGA} is synchronous.
 %
%
%
			%\begin{figure} [H] 				
				%\centering
				%\includegraphics[page=1,trim={0cm 0cm 0cm 0cm},clip=true,width=85mm]{figurastfm/Chapter3/Imagenes/fpga.png}
				%\caption{Typical logic cell \cite{wiki}} \label{fpgafig}
			%\end{figure}
			%
			%Among the concepts introduced before, \textbf{Lookup tables} (\acrshort{LUT}) have special importance. They are arrays which replaces runtime computation with simple array indexing which is time-saver in processing terms, as retrieving a value from memory is often faster than I/O operation \cite{wiki}. \acrshort{LUT} are in charge of storing the truth table of the any boolean function. It is a parameter sometimes used to estimate the \textit{capacity} of an \acrshort{FPGA}.
			%
			%Regarding the programming, the user normally provides a certain design using some Hardware Description Language, as mentioned. A netlist is generated by a procedure called \textbf{synthesis}, which will be fitted to the \acrshort{FPGA} architecture using a process called \textbf{place-and-route}. As a result of this process, the user will obtain a performance report exposed via timing analysis or simulation. Finally, the generated binary file is used to program the \acrshort{FPGA}.
			%
		%\acrshort{FPGA} are particularly suitable for this Co-processing module because of their versatility; the variety of implementations possible is countless; for example, one of the most featured applications is \textbf{embedding a processor} inside an \acrshort{FPGA}. It has many advantages, e.g.: specific peripherals can be chosen based on the application, mitigates obsolescence, reduces costs and allows impressive customization. For instance, there is a complete flexibility to select any combination of peripherals and controllers which can be directly connected to the processor's bus. This point allows meeting non-standard requirements, for example a \acrshort{COTS} processor with ten \acrshort{UART} may be impossible to find, but it is easy to implement in an \acrshort{FPGA} embedded processor \cite{fpgaproc}. An embedded processor implementation example can be found on \cite{fpgacalvo}. Another possible implementations using \acrshort{FPGA} are generating digital interfaces such as \acrshort{I2C} or generating \acrshort{PWM}.
	%
%\subsubsection{Communications}
%
%For obvious reasons, Communications subsystem or, more specifically, Telecommand~\&~Telemetry~(\acrshort{TCC}) are vital in CubeSats. Communications must be understood not only as the ones between the \glsname{cubesat} and the \glsname{ground} but also inside the \glsname{cubesat} itself. One of the simplest ways to communications is between \textbf{wired} and \textbf{wireless}. Next, a brief analysis of the useful technologies for CubeSats is introduced.
%
%\paragraph{Wired}
%
%Normally, wired connections are necessary even in CubeSats intended to fly to be able to program data, perform test and so forth. In this project, because of the simulation parcel, wired communications will be specially important.
%
%\subparagraph{Serial}
%
%This in an historical wired interface. Serial communication is based on sequentially sending data one bit at a time. Serial ports are typically identified as such which comply with RS-232 standard. One of the main drawbacks of serial communications is their slowness.
%
%Although it is considered to be deprecated, it is still usual in many electronics systems, as debugging port or just as a contingency way of communication. In fact, these could its usefulness in a \glsname{cubesat}. Throughout years, serial ports have led to new derived technologies, such as \acrshort{USB}. \autoref{serial} depicts DE-9 connectors, to be used with RS-232 standard.
%
			%\begin{figure} [H] 				
				%\centering
				%\includegraphics[page=1,trim={0cm 0cm 0cm 0cm},clip=true,width=65mm]{figurastfm/Chapter3/Imagenes/serial.jpeg}
				%\caption{DE-9 connector \cite{wiki}} \label{serial}
				%\vspace{-1cm}
			%\end{figure}
%
%
%\subparagraph{Universal Serial Bus}
%
%One of the most extended interfaces in industry. Derived from serial ports, they were initially intended to standardize the connection of peripherals, and finally have largely replaced its preceding interfaces.
%
%\acrshort{USB} has plenty advantages, improved ease of use as it is self-configurable and hot pluggable. It is also much more faster than serial ports, reaching 5 Gb/s in the 3.0 version, one of the latest \cite{wiki}. They are also extensively used as power ports. \autoref{usbfig} shows some of the standard connectors pinout.
%
			%\begin{figure} [H] 				
				%\centering
				%\includegraphics[page=1,trim={0cm 0cm 0cm 0cm},clip=true,width=80mm]{figurastfm/Chapter3/Imagenes/usb.png}
				%\caption{USB standard connectors pinout \cite{wiki}} \label{usbfig}
				%%\vspace{-0.5cm}
			%\end{figure}
%
%Because of their versatility, \acrshort{USB} interfaces may have numerous applications on CubeSats, from programming to power ports or to downloading data.
%
%\subparagraph{Ethernet}
%
%Standardized first in 1983 as \acrshort{IEEE}802.3 is by far the most used wired technologies in \acrshort{LAN}. Systems communicating over this interface divide streams of data into \textbf{frames} which contains source and destination addresses and allows error-checking. 
%
%In CubeSats, Ethernet interfaces may be useful to allow communications using high-level protocols such as \acrshort{SSH}, which ease communications and configuration of the system.
%
			%\begin{figure} [H] 				
				%\centering
				%\includegraphics[page=1,trim={0cm 0cm 0cm 0cm},clip=true,width=70mm]{figurastfm/Chapter3/Imagenes/ethernet.png}
				%\caption{Ethernet standard connector \cite{ethernet}}
				%\vspace{-0.2cm}
			%\end{figure}
%
%
%\paragraph{Wireless}
%
%Wireless communications are vital in real CubeSats, for obvious reasons. In orbit CubeSats communicates with \glsname{ground} using radio-links and dedicated antennas. It is normally a whole subsystem within Communications system. It goes beyond the scope of this project analyzing that kind of implementations; extended radio-link analysis and calculation can be found on \cite{tfg}. In this section, it will be briefly introduced local area wireless communications, which can be used as base for future developments focused on real long-distance links.
%
%\subparagraph{Radio-frequency} \label{radiofr}
%
%Communications over the air are the only choice when dealing with space communications. There is variety of frequencies used in radio-links, depending on distance, power and numerous additional factors. According to the frequency used, the link will have certain particularities which shall be taken into account. \autoref{tfgradio} schematically shows an example of radio path between a \glsname{ground} based on Granada and \acrshort{ISS}; it illustrates some of the aspects to be considered in a long-distance radio-frequency communication such as elevation angle or effective distance, see \autoref{tfgradio}.
%
			%\begin{figure} [H] 				
				%\centering
				%\includegraphics[page=1,trim={0cm 0cm 0cm 0cm},clip=true,width=130mm]{figurastfm/Chapter3/PDF/radio_l.pdf}
				%\caption{Radio-path basic elements \cite{tfg}} \label{tfgradio}
			%\end{figure}
%
%
%Although this Project will not tackle this kind of communications, short-distance radio link may be considered useful for the Simulation Platform branch of the project. Contrary to the links designed for larger distances, short-distance communications over the air exhibits way lower complexity and cost due to the minor requirements in relation with power, losses constraints, among others.
%
%
%\subparagraph{IEEE802.11}
%
%Commonly known as Wi-Fi (a trademark from the Wi-Fi Alliance), it is based on \acrshort{IEEE}802.11 standard and it is the most extended wireless communication protocol for local areas. Once again, this makes it specially suitable for this Simulation Platform because of its ease of use and extended use. On the other hand, it cannot be used for space radio-links.
%
%Since its origins, back in 1997, the standard has going through a significant number of reviews. Nowadays the most usual versions are \acrshort{IEEE}802.11g/n/ac reaching data rates up to 1300 Mbit/s in the 5 GHz band \cite{wiki}.
%
%As a result of the wide use of \acrshort{IEEE}802.11, the majority of the systems related with it are highly standardized (\acrshort{IC}, antennas, amplification subsystems, etc.) which reduces complexity and eases integration.
%
%\subsubsection{Payload. Sensors.}
%
%The \glsname{payload} is generally known as the amount of cargo capacity of an aircraft, including fuel and people. However, it may also refer to the equipment specifically intended to perform a certain mission while in orbit, for instance a camera or a \glsname{startracker}.
%
%In CubeSats field, a common \glsname{payload} example are \textbf{sensors}: they are typically a crucial part of any mission, not only as part of the \glsname{payload} but also a necessary part for the correct functioning of the whole system. Normally they are some kind of electronic device which takes some measurement or perform a certain action depending on the inputs it receives. Some examples of sensors have already been treated before, see \ref{imus}. Another possible sensors are: barometer, thermometer, magnetometer, lightmeter, tachometer... the list is countless. Many of them will be analyzed and implemented in \autoref{chap:chapter4}. \autoref{sensores} shows some examples of them.
%
%While this analysis considers sensors as part of the \acrshort{OBC}, given that eventually the data will be treated there, they can belong to a different subsystem such as \acrshort{ADCS}, see \ref{imus}. 
%
			%\begin{figure}[H]
			%\centering
			%\subfloat[\acrshort{IMU}]{\includegraphics[page=1,trim={0cm 0cm 0cm 0cm},clip=true,width=65mm]{figurastfm/Chapter3/Imagenes/MPU.png}}
			 %\quad
			%\subfloat[Barometer and thermometer]{\includegraphics[page=1,trim={0cm 0cm 0cm 0cm},clip=true,width=65mm]{figurastfm/Chapter3/PDF/BMP.pdf}}
			%\caption{Sensor examples \cite{hobbyking}} \label{sensores}
			%\vspace{-0.55cm}
%\end{figure}	
%
%
%\subsubsection{Flight Software. On-Board Data Handling.}
%
%In charge of all the exposed, there must be a \textbf{governor}, a coordinator which controls the hardware and deals with the different stages and situations of a mission, this is the \textbf{Flight Software} (FSW). As in a significant number of the subsections analyzed through this Master's Thesis, Flight Software is such a large field that it allows a complete project focused on it. Therefore, hereby only a basic classification is introduced; this will be an important issue at the design stage and also for future improvements of the system.
%
%As far as this project concerns, considered flight software can be subdivided into the following.
%
%\paragraph{Non-real-time Operative System}
%
%Non-real-time Operative Systems is basically a general purpose \acrshort{OS} to be used with personal computers, servers, etc. The main difference with \acrshort{RTOS} is \textbf{determinism}. Non-\acrshort{RTOS} are not deterministic as tasks will not run at a certain time and for a certain time, there are no guarantees for critical tasks, exhibits high latency because of using unpredictable virtual memory, as well as \glsname{jitter}.
%
%The previous characteristics translates into the use of \textbf{non-preemptive schedulers}. Examples of these kind of \acrshort{OS} are widely used Windows\,\textsuperscript{\textregistered}, Debian, etc.
%
%\paragraph{Real-time Operative System}
%
%In contrast with the latter, \acrshort{RTOS} are completely deterministic, this is, how and when a task will run given whatever conditions defined for it to do so, it is \textbf{guaranteed}. They are intended to process data as it comes in, without buffer delays. There are a huge amount of systems nowadays which must use \acrshort{RTOS}: cars, \textbf{spacecrafts}, avionics, critical systems...
%
%
%In sum, \acrshort{RTOS} must be able to compute a task in a \textbf{limited or predictable amount of time} i.e., is time-bounded. However, this behaviour has nothing to do with processing speed but with a known deadline (a second, an hour or a month) and reduced \glsname{jitter}.
%
%From a technical perspective, this implies not to use virtual memory, strict scheduling (preemptive) and avoid non-deterministic elements. \acrshort{RTOS} are usually much smaller than general purpose \acrshort{OS} in order to ease maintenance and find sources of delay. One of the most used \acrshort{RTOS} is \textbf{VxWorks} \cite{vxworks}; Linux non-preemptive kernel can also be \textbf{patched} to allow real-time behaviour.
%
%
			%\begin{figure} [H] 				
				%\centering
				%\includegraphics[page=1,trim={0cm 0cm 0cm 0cm},clip=true,width=85mm]{figurastfm/Chapter3/Imagenes/vxworks.png}
				%\vspace{0.3cm}
				%\caption{VxWorks logo \cite{vxworks}} 
				%%\vspace{-2cm}
			%\end{figure}
%
%\subsubsection{Formal Requirements Definition} \label{obcreqs}
%
%
%\begin{table} [H]
%\centering
%
%\begin{tabularx}{\linewidth}{lX}
%
%\multicolumn{1}{c}{\textbf{Ref.}}                      & \multicolumn{1}{c}{\textbf{Formal Requirements}}                    \tabularnewline \specialrule{1.1pt}{1pt}{1pt}
%OBC.FoR.1                                              & The \acrshort{OBC} shall have an \acrshort{I2C} interface which allows communication with the rest of subsystems and components. \tabularnewline \midrule
%OBC.FoR.2                                              & The \acrshort{OBC} shall have an \acrshort{SPI} interface which allows communication with the rest of subsystems and components. \tabularnewline \midrule
%OBC.FoR.3                                            & The \acrshort{OBC} shall have a programmable \acrshort{FPGA} of at least 2500 \acrshort{LUT}. \tabularnewline \midrule
%OBC.FoR.4                                                   & The \acrshort{OBC} shall have one RJ-45 connector. \tabularnewline \midrule
%OBC.FoR.5                                                   & The \acrshort{OBC} shall have at least two \acrshort{USB} Micro-B device-ports and capability to deal with four.  \tabularnewline \midrule
%OBC.FoR.6                                                   & The \acrshort{OBC} shall be programmable with an external computer using an \acrshort{USB} Micro-B connector. \tabularnewline \midrule
%OBC.FoR.7                                                   & The \acrshort{OBC} shall have a \acrshort{LED} system to check the correct functioning of the \acrshort{EPS}  \tabularnewline \midrule
%OBC.FoR.8                                                   & The \acrshort{OBC} shall have at least two \acrshort{USB} device-ports and capability to deal with four.  \tabularnewline \midrule
%OBC.FoR.9                                                   & The \acrshort{OBC} shall have a barometer with a resolution of at least 0.1 hPa.  \tabularnewline \midrule
%OBC.FoR.10                                                   & The \acrshort{OBC} shall have a thermometer with a resolution of at least 0.1\textdegree~C.  \tabularnewline \midrule
%OBC.FoR.11                                                   & The \acrshort{OBC} shall have at least 8 \acrshort{ADC} channels. \tabularnewline \midrule
%OBC.FoR.12                                                   & The \acrshort{OBC} shall have 6 analog sun sensors, one per each face of the \glsname{cubesat}. \tabularnewline \midrule
%OBC.FoR.13                                                   & The \acrshort{OBC} shall have a Real Time Clock (\acrshort{RTC}) which provides the whole system with time information. \tabularnewline \midrule
%OBC.FoR.14                                                   & The \acrshort{OBC} shall have an \acrshort{SD} card holder, to store telemetry data.  \tabularnewline \midrule
%OBC.FoR.15                                                   & The \acrshort{OBC} shall have an integrated \acrshort{IEEE}802.11g device which allows wireless communications. \tabularnewline \midrule
%OBC.FoR.16                                                   & The \acrshort{OBC} shall implement a linux-based \acrshort{RTOS}. \tabularnewline \midrule
%
%\end{tabularx}
%\caption{\acrshort{OBC} - Formal Requirements}
%\vspace{-0.5cm}
%%\label{forgstation}
%
%\end{table}
%
%
%
%
%
%
%
%
%

% Capítulo 4: Prototipo LM11
\chapter{System Design} \label{chap:chapter4}

After the different stages of analysis and requirements definition performed in former chapters, this forth chapter will address the \textbf{system design}. It will be structured as \autoref{chap:chapter3}, so the flow design coincides as much as possible with the analysis flow. 

Firstly, the mechanical platform is designed and characterized using professional \acrshort{CAD} tools such as \glsname{solid}. Also it is manufactured \textit{in-house} using aluminum milling machine. Regarding the \glsname{ground}, it is developed using an extended framework in space field, even used by \acrshort{NASA}. Finally, some of the most important subsystems of \glsname{cubesat} are developed, concluding with a complete \glsname{engmodel}, expandable and which can be used as a solid base to keep working on participating in the \textit{Fly~your~Satellite!} \cite{flyyour} program.

\section{Inertial 2D Orbit Simulator (I2DOS)} \label{i2dossec}

Among the different platforms analyzed in \autoref{mechplat}, this project will include an \textbf{inertial} one, which has been called Inertial 2D Orbit Simulator (\acrshort{I2DOS}). Particularly, the designed platform will be \textbf{rotational} along Z axis. \acrshort{I2DOS} will be composed of two parts, on the one hand, the rotational platform itself, and on the other hand, the base which will support the platform and will allow including external elements such as irradiance sources, according to the \textbf{formal requirements}.

\subsection{Inertial Platform}

\subsubsection{Design and mechanical characterization}


The inertial platform is designed using \glsname{solid} and sequentially improved through different redesigns. It is fully designed from scratch, with a physical background behind, reasoned later. It must have some kind of support which hosts the \glsname{cubesat} during simulation. Given the circular shape of the platform itself, that support will also be circular. Besides, it must have adequate attachment for a 1U \glsname{cubesat}. \autoref{superior} depicts a high quality render of that piece, including attachments.

\vspace{-2cm}
\begin{figure}[H]
			\centering
			\includegraphics[page=1,trim={0cm 0cm 0cm 0cm},clip=true,width=165mm]{figurastfm/Chapter4/Imagenes/superior.png}
			\vspace{-1.5cm}
			\caption{1U \glsname{cubesat} support} \label{superior} 
						\vspace{-1cm}
		\end{figure}


Top view and bottom view are depicted in Figure\autoref{topplat} and Figure\autoref{botplat} while another 3D render of the final design is shown in \autoref{inertiaplat3d}.


			\begin{figure}[H]
			\centering
			\subfloat[Top view \label{topplat}]{\includegraphics[page=1,trim={0cm 0cm 0cm 0cm},clip=true,width=75mm]{figurastfm/Chapter4/Imagenes/vista_superior.png}}
			 \quad
			\subfloat[Bottom view \label{botplat}]{\includegraphics[page=1,trim={0cm 0cm 0cm 0cm},clip=true,width=75mm]{figurastfm/Chapter4/Imagenes/vista_inferior.png}}
			\caption{Inertial platform} 
						\vspace{-2cm}

\end{figure}


\newpage
\vspace*{\fill}
\begin{figure}[H]
			\centering
			\includegraphics[page=1,trim={0cm 0cm 0cm 0cm},clip=true,width=165mm]{figurastfm/Chapter4/Imagenes/inertial.png}
%			\vspace{-2cm}
			\caption{Inertial platform 3D render} \label{inertiaplat3d} 
		\end{figure}
\vspace*{\fill}
\newpage

Regarding the \textbf{mechanical characterization}, simplifying, the \textit{functioning} is based on the inertia moment stored on the bigger ring at the bottom area ($I_1$) plus the one stored on the \glsname{cubesat} support ($I_2$). The first is considered a thick-walled cylindrical tube with open ends of inner radius $r_1$, outer radius $r_2$, length $h$, mass $m$ and a density $\rho$ whose moment of inertia is given by the first term in \autoref{inertiaring}; the second can be approximated by a solid cylinder of radius $r$, height $h$ and mass $m$ and its moment of inertia is given by the second term in \autoref{inertiaring} \cite{wiki}. The inertia moment stored at the rods is considered negligible as well as the components related to X and Y axis.

\begin{equation} \label{inertiaring}
I_{\text{z}}=I_{\text{z}1}+I_{\text{z}2}=\frac{\pi\rho h}{2}(r^{4}_{2}-r^{4}_{1})+\frac{1}{2}mr^{2}
\end{equation}

It has been considered a density $\rho$ of \SI{1.24}{g/cm^3} corresponding to \acrshort{PLA} material and a mass $m$ given by \glsname{solid} of 453 g.  Substituting with the rest of design parameters, the approximated total moment of inertia stored when rotating around Z axis is \SI{94.6}{kg/cm^2}. \glsname{solid} also calculates the inertia moment exactly, featuring a value of  \SI{108}{kg/cm^2}. It makes a difference of about 13 \% which is coherent with the approximations taken and verifies the result of \autoref{inertiaring}.

\acrshort{I2DOS}, and particularly this inertial platform, simulates the moment of inertia that the \glsname{cubesat} will face when deployed, in \glsname{tumbling}. It will have to use different detumbling mechanisms to counter that uncontrolled movement.

			\vspace{-0.5cm}


\subsubsection{Manufacturing}

The inertial platform is 3D-printed using \acrshort{PLA} and aluminum rods as the ones depicted in \autoref{inertiaplat3d}. The final result is shown in \autoref{fotoinertial}.

\vspace{-0.3cm}
\begin{figure}[H]
			\centering
			\includegraphics[page=1,trim={0cm 0cm 0cm 0cm},clip=true,width=69mm]{figurastfm/chapter4/imagenes/i2dos_b.png}
			\caption{Manufactured Inertial platform} \label{fotoinertial} 
			\vspace{-2cm}
\end{figure}

\subsection{Base}

\subsubsection{Design}

In order for the inertial platform to rotate, it is necessary a base which supports it. It is also designed using \glsname{solid} according to the requirements (see section \ref{frmech}); it must allow including external components such as sun simulators. The high quality 3D render of the base is depicted in \autoref{basemech}.

\begin{figure}[H]
			\centering
			\includegraphics[page=1,trim={0cm 0cm 0cm 0cm},clip=true,width=120mm]{figurastfm/chapter4/imagenes/base.png}
			\caption{Base 3D render} \label{basemech} 
			%\vspace{-2cm}
\end{figure}

The center hole holds the inertial base while the external ones allows adding external components; the attachments of the inertial platform are reused. It is designed to be wooden manufactured and the additional components can be attached using 4 cm aluminum rods. \autoref{basecosas} depicts the base with the different rods.

\begin{figure}[H]
			\centering
			\includegraphics[page=1,trim={0cm 0cm 0cm 0cm},clip=true,width=120mm]{figurastfm/chapter4/imagenes/base_cosas.png}
			\caption{3D render of the base with the aluminum rods} \label{basecosas} 
\end{figure}
\subsubsection{Manufacturing}

The base was manufactured at one of the \textit{makerspace} of the University of Granada in which the base could be cut. Video \autoref{videojulio} shows part of that process (Adobe Reader needed).


\begin{videoFloat}[H]
\centering
\includemovie[text={\includegraphics[page=1,trim={0cm 0cm 0cm 0cm},clip=true,width=90mm]{figurastfm/Chapter4/Imagenes/foto_video.jpg}}]{10cm}{6cm}{FigurasTFM/Chapter4/Video/makerspace.wmv}
\caption{Wooden base manufacturing (double click)} \label{videojulio}
\vspace{-0.5cm}
\end{videoFloat}

Figure\autoref{base1} and Figure\autoref{base2} depict the final result, with an external irradiance source attached, analyzed next in \autoref{irrasources}. 


			\begin{figure}[H]
			\centering
			\subfloat[Front \label{base1}]{\includegraphics[page=1,trim={0cm 0cm 0cm 0cm},clip=true,width=80mm]{figurastfm/Chapter4/Imagenes/frontbase.jpg}}
			 \quad
			\subfloat[Detail \label{base2}]{\includegraphics[page=1,trim={0cm 0cm 0cm 0cm},clip=true,width=80mm]{figurastfm/Chapter4/Imagenes/detailbase.jpg}}
			\caption{Inertial 2D Orbit Simulator (I2DOS)} \label{i2dosfinal}
			\vspace{-0.5cm}
\end{figure}

This completes the Inertial 2D Orbit Simulator (\acrshort{I2DOS}).

\subsection{Irradiance sources characterization} \label{irrasources}

In order to have a simulation platform as realistic as possible, it is needed an irradiance source which emulates the Sun. In this subsection, different lightning options are compared, following the analysis performed in section \ref{irradsources}. As studied before, it is desirable a spectrum as similar as possible to the Sun's, so the spectral response of the solar cells is such that the output power is maximum.

\subsubsection{Xenon Arc Sun Simulator} \label{xenonar}

Firstly, the Xenon Arc Sun Simulator of the \glsname{GranaSAT} laboratory is characterized following the procedures detailed in \autoref{chap:chapter3}. Although it is not usable in the proper mechanical platform designed in \autoref{i2dossec}, it is useful as a reference and can be eventually used in order to perform real missions simulations. This simulator features six different power modes, in increasing order. Figure\autoref{canon1} and Figure\autoref{canon2} show the simulator to be characterized.

\begin{figure}[H]
			\centering
			\subfloat[Side view \label{canon1}]{\includegraphics[page=1,trim={0cm 0cm 0cm 0cm},clip=true,width=85mm]{figurastfm/Chapter4/Imagenes/canon1.png}}
			 \\
			\subfloat[Front detail \label{canon2}]{\includegraphics[page=1,trim={0cm 0cm 0cm 0cm},clip=true,width=50mm]{figurastfm/Chapter4/Imagenes/canon2.png}} 
			\caption{Xenon Arc Sun Simulator at \glsname{GranaSAT} laboratory}
			%\vspace{-2cm}
\end{figure}
\newpage
\paragraph{Spectral response} \label{spectral}

To characterize the spectral response of the simulator, we had the collaboration of the  Department of Optics from the University of Granada. To perform the measurement, it was used an spectrometer THORLABS CCS200/M \cite{spectrometro} as the one depicted in \autoref{spetrometer}.

\begin{figure}[H]
			\centering
			\includegraphics[page=1,trim={0cm 0cm 0cm 0cm},clip=true,width=80mm]{figurastfm/chapter4/imagenes/spectrometer.jpg}
			\caption{Spectrometer THORLABS CCS200/M \cite{spectrometro}} \label{spetrometer} 
						\vspace{-1cm}

\end{figure}


The Sun simulator was measured with the different selectable power levels, getting the plot shown in \autoref{canonspectral}. 

\begin{figure}[H]
			\centering
			\includegraphics[page=1,trim={0cm 0cm 0cm 0cm},clip=true,width=165mm]{figurastfm/chapter4/PDF/Spectrum_sun_Sim.pdf}
			\caption{Spectral response of the different power levels} \label{canonspectral}
			%\vspace{-2cm}
\end{figure}

\autoref{canonspectral} illustrates a relative intensity plot, as the used spectrometer cannot measure absolute values, which will be determined later. It is easy to see that the spectrum remains the same regardless of the chosen level, as expected, featuring an spectral response in the range between 400 and 700 nm. It is worth mentioning that power does not increase uniformly with levels, for instance, while the increase is greater between levels 3 and 4, the difference between levels 5 and 6 is almost imperceptible.  

Recalling that the \glsname{cubesat} will count with solar panels as power source, it is interesting to see how the spectral response of the simulator matches the Sun's. It is depicted in \autoref{suncanon}, using \acrshort{AM}1.5 again.


\begin{figure}[H]
			\centering
			\includegraphics[page=1,trim={0cm 0cm 0cm 0cm},clip=true,width=160mm]{figurastfm/chapter4/PDF/Spectrum_comparison.pdf}
			\caption{Spectrum comparison} \label{suncanon}
			%\vspace{-2cm}
\end{figure}

The simulator exhibits a narrower wavelength range in comparison with Sun's, which features a wide spectral response between 300 nm up to 1000 nm; however, the area with a greater intensity is similar to the simulator's one, between 450 and 600 nm. It must be recalled from \autoref{negrous} that a standard Si solar cell gets excited in the wavelength range 400-1100 nm, so it is reasonable to expect a considerably worse performance when this Xenon Arc Sun Simulator is used.  

\paragraph{Irradiance}

After characterizing the xenon arc simulator in terms of spectral response, this section addresses the irradiance power and its variation with distance. In order to perform that measurement, it is used the Silicon-base pyranometer \textbf{Apogee~SP-110-SS}~\cite{apogee}. Figure\autoref{pyranomapog} depicts it while Figure\autoref{apogeespec} shows its spectral response.

			\begin{figure}[H]
			\centering
			\subfloat[Detail \label{pyranomapog}]{\includegraphics[page=1,trim={0cm 0cm 0cm 0cm},clip=true,width=25mm]{figurastfm/Chapter4/Imagenes/APOGEE.jpg}}
			 \quad
			\subfloat[Spectral response \label{apogeespec}]{\includegraphics[page=1,trim={0cm 0cm 0cm 0cm},clip=true,width=110mm]{figurastfm/Chapter4/PDF/pyranomspec.pdf}}
			\caption{Apogee SP-110-SS \cite{apogee}}
			\vspace{-0.75cm}
\end{figure}

 Using the pyranometer, several measurements are taken at different distances and plotted in \autoref{irradistance}.

\begin{figure}[H]
			\centering
			\includegraphics[page=1,trim={0cm 0cm 0cm 0cm},clip=true,width=140mm]{figurastfm/chapter4/PDF/irradiance_Distance.pdf}
			\caption{Irradiance decay with distance} \label{irradistance}
			\vspace{-2cm}
\end{figure}

As \autoref{irradistance} shows, the irradiance decays with distance approximately in a slightly exponential or power tendency. Marked in green, it is the point in which irradiance matches the amount received at the Earth's surface (\SI{1120}{W/m^2}), approximately at a distance of 140 cm. This shall be the point to perform simulations when emulating the Sun.

\paragraph{Lightning distribution}

As performed in \autoref{chap:chapter3} with the measurements from \cite{catalan}, it is analyzed the lightning distribution of this Sun simulator, in order to detect off-centered beam or any other inconsistencies. This is a complex procedure which may be performed with different instruments. In this case, the previously mentioned pyranometer is used along with a designed \textbf{test template} to be lighted up with the xenon arc simulator. \autoref{plantillacanon} shows that template, designed with AutoCAD.

\begin{figure}[H]
			\centering
			\includegraphics[page=1,trim={0cm 0cm 0cm 0cm},clip=true,width=100mm]{figurastfm/chapter4/PDF/template.pdf}
			\caption{Lightning distribution test template} \label{plantillacanon}
			\vspace{-0.7cm}
\end{figure}

The procedure consists in performing measurements with the pyranometer at the different angles depicted in the test template, until completing the whole circumference. Besides, these measurements must be taken at several distances, marked in red. When completed, it is possible to plot the lightning distribution, just as performed in \autoref{catorigin} and \autoref{catinterp}. The results are shown in \autoref{distrib1} and \autoref{distrib2}.

	\afterpage{
\begin{landscape}
%\vspace*{\fill}
\begin{figure}
			\centering
			\includegraphics[page=1,trim={0cm 0cm 0cm 0cm},clip=true,width=175mm]{figurastfm/Chapter4/PDF/Lightning_distrib_pcolor_jet.pdf}
			\caption{Lightning distribution (raw data)} \label{distrib1}
\end{figure}
%\vfill
\end{landscape}
}
	
	
\afterpage{
\begin{landscape}
%\vspace*{\fill}
\begin{figure}
			\centering
			\includegraphics[page=1,trim={0cm 0cm 0cm 0cm},clip=true,width=175mm]{figurastfm/Chapter4/PDF/Lightning_distrib_contourf_jet.pdf}
			\caption{Interpolated Lightning distribution} \label{distrib2}
\end{figure}
%\vfill
\end{landscape}
}

As in \autoref{chap:chapter3}, the measurements are plotted twice, one with the \textit{raw} data in \autoref{distrib1} and another interpolating the data, allowing a more understandable graphic in \autoref{distrib2}. In order to ease the comparison, they are both plotted together in \autoref{comparison}. The ones above correspond to the simulator at \glsname{GranaSAT} laboratory and the ones below to the one used in \cite{catalan}.

\begin{figure}
			\centering
			\includegraphics[page=1,trim={0cm 0cm 0cm 0cm},clip=true,width=165mm]{figurastfm/Chapter4/PDF/Lightning_distrib_comparison.pdf}
			\caption{Lightning distribution comparison} \label{comparison}
\end{figure}

The first conclusion which can be extracted in comparison with the simulator used in \cite{catalan} is a noticeably less powerful xenon bulb in my case, even at a shorter distance (1.8~m vs. 2.7~m) the maximum irradiance is 33 \% lower than his. By inspection at \autoref{irradistance}, the equivalent irradiance at 2.8 m is about \SI{500}{W/m^2}, once again one third lower. Therefore, although the power of the xenon bulb used is unknown, assuming similar losses, it can be estimated around \textbf{650 W}. Another possibility is that the bulb has reached its end-of-life and power gets lower. 

As stated in \ref{irradsources}, the closer the simulated is used, the greater the beam gets off-centered; indeed, as seen in \autoref{comparison}, the beam of the simulator in our case is clearly off-centered, even more so than the other, at a distance of about 5 cm from the center, within the range between 180\textdegree~and~290\textdegree. Therefore, there are two possibilities in order to use this Sun simulator:
\newpage
\begin{itemize} [noitemsep,topsep=0pt]
\item Use it at greater distances, so the beam gets centered, sacrificing power, which will get below the \SI{1100}{W/m^2}. Allows characterizing the solar panel in the center of the test template (see \autoref{plantillacanon}) but it is not useful for real mission simulations.	\\
\item Use it at closer distances, with an off-centered beam as depicted in previous plots, but with realistic irradiance. Allows emulating Sun's irradiance, needed in order to check functioning in a real mission scenario, but needs to place the solar panel in the area where the beam focus, stated before.
\end{itemize}

The choice will depend on the purpose of the simulation.

\subsubsection{LED Sun Simulator} \label{ledsim}
%\vspace{-0.5cm}

\begin{wrapfigure}{r}{0.28\textwidth} 
	\centering
	\includegraphics[width=0.20\textwidth]{figurastfm/chapter4/imagenes/led.jpg}
	\vspace{0.1cm}
	\caption{\acrshort{LED} proposed} \label{ledchino}
\end{wrapfigure}


Because of its size, the xenon arc simulator from \ref{xenonar} cannot be used along with the inertial platform designed. Therefore, in this section it will be characterized another technology previously analyzed in \autoref{chap:chapter3}, \acrshort{LED}. According to the requirements, it has to be usable with the platform, so it must be small-sized. It is proposed using a low-cost thermally-bonded \acrshort{LED} lamp, such as the one shown in \autoref{ledchino}. \\ \\


%\begin{figure} [H]
			%\centering
			%\includegraphics[page=1,trim={0cm 0cm 0cm 0cm},clip=true,width=40mm]{figurastfm/Chapter4/Imagenes/led.jpg}
			%\caption{\acrshort{LED} proposed} \label{ledchino}
			%\vspace{-0.5cm}
%\end{figure}
\vspace{-1cm}
It is a 100 W lamp to be supplied between 20 V and 32 V, with a maximum consumption of 3 A. It may reach high temperatures so it needs cooling; besides, as it is desired a lightning beam as concentrated as possible, it raises the need for a \textbf{collimator}. To satisfy these needs, it is used a kit designed for this kind of lamps, which includes both, a collimator and a 12 V fan. It is shown in \autoref{collimator}. When it is completely assembled, the result is the one which could be seen in \autoref{i2dosfinal}.

%\vspace{-0.5cm}


			\begin{figure}[H]
			\centering
			\subfloat[Collimator and body]{\includegraphics[page=1,trim={0cm 0cm 0cm 0cm},clip=true,width=50mm]{figurastfm/Chapter4/PDF/cuerpoled.pdf}}
			\quad
			\subfloat[Copper base to cool the \acrshort{LED}]{\includegraphics[page=1,trim={0cm 0cm 0cm 0cm},clip=true,width=50mm]{figurastfm/Chapter4/Imagenes/topled.jpg}}
			 \quad
			\subfloat[Fan]{\includegraphics[page=1,trim={0cm 0cm 0cm 0cm},clip=true,width=50mm]{figurastfm/Chapter4/Imagenes/cooler.jpeg}}
			\caption{Kit used for the \acrshort{LED} lamp} \label{collimator}
			\vspace{-1cm}
\end{figure}


\paragraph{Spectral response}

The spectral response of the \acrshort{LED} is measured using the same procedure detailed in \ref{spectral}. The results are plotted in \autoref{specralled}, along with the Sun's.

\begin{figure} [H]
			\centering
			\includegraphics[page=1,trim={0cm 0cm 0cm 0cm},clip=true,width=165mm]{figurastfm/Chapter4/PDF/Spectrum_comparison_led.pdf}
			\caption{Spectrum comparison} \label{specralled}
\end{figure}

In this case, the \acrshort{LED} spectrum covers the range 400 nm to 650 nm, with a pronounced peak at 500 nm in which almost no intensity is irradiated. Therefore, the range of a typical solar panel is slightly covered and in order to get at least an acceptable efficiency, a really high power would be needed.  

\paragraph{Irradiance and consumption}

Irradiance is measured again using the same pyranometer. As this lamp is intended to be used with the inertial platform at an approximately fixed distance, it will be characterized at that single point, \textbf{20 cm far}. \autoref{irradiancemeas} shows part of the assembly needed to perform the measurement. 

\begin{figure} [H]
			\centering
			\includegraphics[page=1,trim={0cm 0cm 0cm 0cm},clip=true,width=100mm]{figurastfm/Chapter4/Imagenes/irrad_meas.jpg}
			\caption{\acrshort{LED} irradiance measurement} \label{irradiancemeas}
\end{figure}

The measurement needs to control simultaneously two devices: on the one hand, the power source, which must sweep between 20 V and 32 V and on the other hand, the oscilloscope to check for signal integrity first and perform the measurements. Specifically, the power source used is a \textbf{Siglent SPD3303X} while the oscilloscope is an \textbf{Agilent MXO-X-4104A}, both are shown in \autoref{equipos}.

			\begin{figure}[H]
			\centering
			\subfloat[Siglent SPD3303X \cite{Siglent}]{\includegraphics[page=1,trim={0cm 0cm 0cm 0cm},clip=true,width=50mm]{figurastfm/Chapter4/Imagenes/SIGLENT.jpg}}
			\quad
			\subfloat[Agilent MXO-X-4104A \cite{agilent}]{\includegraphics[page=1,trim={0cm 0cm 0cm 0cm},clip=true,width=55mm]{figurastfm/Chapter4/Imagenes/AGILENT.png}}
			\caption{Measurement devices} \label{equipos}
			%\vspace{-1cm}
\end{figure}


Firstly, the irradiance is measured along with the \acrshort{LED} current consumption.  In order to get it done synchronically, the following Python code is developed.


%Python code highlighting
\lstset{  backgroundcolor=\color{white},  commentstyle=\color{codegreen}, keywordstyle=\color{magenta},   numberstyle=\tiny\color{codegray}, stringstyle=\color{codepurple},  basicstyle=\ttfamily\footnotesize,  breakatwhitespace=false,           breaklines=true,                   captionpos=b,                      keepspaces=true,                   numbers=left,                      numbersep=5pt,                    showspaces=false,                  showstringspaces=false,  showtabs=false,                    tabsize=2,  stepnumber=3, numberstyle=\tiny, basicstyle=\linespread{0.4}}

\begin{lstlisting}[language=Python, caption=Polling based measurements script]
import visa
from visa import constants
import vxi11
import csv
import pandas as pd
import time
import math
import os
import numpy as np

# GPIB INIT
# visa.log_to_screen()
SG =  vxi11.Instrument("192.168.1.119")
OSC = visa.ResourceManager('@py').get_instrument('TCPIP0::192.168.1.121::inst0::INSTR')
OSC.timeout=2500000

# IDENTIFYING
print("SG found: " + SG.ask("*IDN?").strip())
print("OSC found: " + OSC.query('*IDN?').strip())

#OSC Setting-up

OSC.write(':CHANnel1:DISPlay ON')
OSC.write(':DISPlay:SIDebar MEASurements')
OSC.write(':MEASure:VPP CHANnel1')

#SG Setting-up

SG.ask("CH1:VOLT 12")
SG.ask("CH2:VOLT 20")

SG.ask("CH1:CURRent 0.3")
SG.ask("CH2:CURRent 3.2")

SG.ask("OUTPut CH1,ON")
SG.ask("OUTPut CH2,ON")

volt_Sweep=np.arange(20,32.5,0.5)
current_Sweep=[]
measured_Osc=[]

#Allow time for the measurement to stabilize at the final values

for V in volt_Sweep:
	
	SG.ask("CH2:VOLT %f" %V)
	
	if V==31 or V==31.5 or V==32:
	
		time.sleep(60)
		
	else:
	
		time.sleep(1)	

	current_Sweep.append(SG.ask("MEASure:CURRent? CH2"))	
	measured_Osc.append(OSC.query(':MEASure:VRMS? CHANnel1'))


voltage_Data=pd.DataFrame(volt_Sweep)
current_Data=pd.DataFrame(current_Sweep)
measured_Voltage_Data=pd.DataFrame(measured_Osc)

pd.concat([voltage_Data,current_Data,measured_Voltage_Data],axis=1).to_csv("data5.csv")

print("Data written to CSV")

SG.ask("OUTPut CH1,OFF")
SG.ask("OUTPut CH2,OFF")

\end{lstlisting}

The results are plotted in \autoref{irracurr}. The tendency is clearly lineal, as expected; the higher the current, the greater the irradiance, until reaching the maximum current allowed by the \acrshort{LED}. The highest irradiance accomplished by the simulator at that distance (20 cm) is \SI{235}{W/m^2}. It is about five times lower than the one received at Earth's surface, therefore, it can be expected a poor performance from the solar panels with this irradiance. Of course it is also way lower than the value produced by the xenon arc lamp at that distance. However, that comparison is not realistic; while this \acrshort{LED} is about 80 W (see \autoref{ivpower}), the xenon arc simulator has an estimated power of about 600 W and is not intended to work at such a short distance. \textbf{Their simulation purposes are different.}

\begin{figure} [H]
			\centering
			\includegraphics[page=1,trim={0cm 0cm 0cm 0cm},clip=true,width=135mm]{figurastfm/Chapter4/PDF/IRRAD_CURR.pdf}
			\caption{Irradiance with \acrshort{LED} current consumption} \label{irracurr}
						\vspace{-0.7cm}
\end{figure}

 However, the irradiance is not linear with \textbf{voltage supply}, because of the exponential relationship between voltage and current in a \acrshort{LED}; indeed, neither the functioning range given by the manufacturer nor the power drawn are completely correct. This can be concluded from the plot in \autoref{ivpower}, which relates I-V curve of the \acrshort{LED} with the power drawn. 

\begin{figure} [H]
			\centering
			\includegraphics[page=1,trim={0cm 0cm 0cm 0cm},clip=true,width=135mm]{figurastfm/Chapter4/PDF/ivpower.pdf}
			\caption{I-V curve and power drawn} \label{ivpower}
			\vspace{-2cm}
\end{figure}

That plot is important because it allows checking the real requirements of the lamp and therefore, of the \acrshort{LED} driver. It can be concluded that the \acrshort{LED} does not start lightning until supplied with 26 V, and the maximum power drawn is about 80 W, unlike the range stated by the manufacturer 20-32 V and 100 W power, respectively. The power losses were obviously expected, but an under consumption of 20 \% is worth noting. Although the expected maximum is not reached, the lamp heats considerably, which makes it necessary to install the fan mentioned before. Even when the fan is on, the temperature borders on 100 \textdegree C, as shown in \autoref{temp}.


\begin{figure} [H]
			\centering
			\includegraphics[page=1,trim={0cm 0cm 0cm 0cm},clip=true,width=40mm]{figurastfm/Chapter4/Imagenes/temp.jpg}
			\caption{Measured temperature on the \acrshort{LED}} \label{temp}
						\vspace{-0.7cm}
\end{figure}
In sum, in order to supply the \acrshort{LED} Sun simulator, it is necessary a \textbf{driver} with an output voltage about 32 V and 100 W of power capability. \autoref{driverchino} shows the chosen one.

\begin{figure} [H]
			\centering
			\includegraphics[page=1,trim={0cm 0cm 0cm 0cm},clip=true,width=100mm]{figurastfm/Chapter4/Imagenes/led_driver.jpeg}
			\caption{Low cost 100 W \acrshort{LED} driver} \label{driverchino}
						%\vspace{-0.7cm}
\end{figure}


\subsubsection{Comparative}

The irradiance sources characterization finalizes comparing the spectrums of the proposed technologies with the Sun's. \autoref{compall} plots this comparison with all the data normalized to the Sun's spectrum irradiance, in order to ease the spectral visualization.


\begin{figure} [H]
			\centering
			\includegraphics[page=1,trim={0cm 0cm 0cm 0cm},clip=true,width=145mm]{figurastfm/Chapter4/PDF/Spectrum_comparison_all.pdf}
			\caption{Spectral comparison of the different irradiance sources} \label{compall}
						%\vspace{-0.7cm}
\end{figure}

Xenon arc and \acrshort{LED} spectrums are actually pretty similar, except for the 500 nm area in the latter. However, from \autoref{compall} and previous plots it is easy to conclude that \textbf{none of them matches the Sun's}, neither spectrum nor power. That is not necessarily a problem, though; the \acrshort{LED} is intended to be used along with \acrshort{I2DOS} so its final purposes are not strictly real simulation. When that is required however, it will be needed to use a different xenon arc simulator. Even if spectral matching is enough, low power and off-centered beam makes it unreliable for real missions simulations.


\section{Ground Station}

Regarding the \glsname{ground}, this Master's Thesis will address the design and launch of the central controller. It will be based on the environment \textbf{Ball~Aerospace~COSMOS} and will count with different user interfaces and a basic telemetry database. It is intended to function as the main point of control and communications with the \glsname{cubesat}.

The designed \glsname{ground} software will be expandable so it can be used as it would in a real mission.

\subsection{Ball Aerospace COSMOS} \label{cosmosgs}

COSMOS is a complete environment designed by Ball Aerospace which is composed of some sample applications and allows defining new interfaces and functionalities, covering the \textbf{full lifecycle} of the mission, as the same interface can be used from test to operation. It is written in \textbf{Ruby} and provides with a command and control system to interact with a variety of embedded systems such as a \glsname{cubesat} but also electronics equipment such as a radio transmitter.  It is also compatible with international standards analyzed before like \acrshort{CCSDS} or \acrshort{XTCE} and is \textbf{open source}.


Using COSMOS allows adding a new abstraction layer to develop the tools needed by the mission focusing on the required functionalities, instead of rewriting existing tools. This makes it a really powerful and productive tool, which has been widely used by international institutions including \acrshort{NASA}, particularly the Goddard Space Flight Center. Implementing it in \glsname{GranaSAT} laboratory gives future students the valuable opportunity to work and train themselves using real cutting-edge technology, which satisfies the academic perspective of this Master's Thesis.

\vspace{-0.5cm}
\subsubsection{Design and configuration}

Once installed, COSMOS design and configuration is based on plain text files which can be either modified directly or using a dedicated tool included with COSMOS, as it is proposed. \autoref{configeditor} displays that configuration editor.


\begin{figure} [H]
			\centering
			\includegraphics[page=1,trim={0cm 0cm 0cm 0cm},clip=true,width=155mm]{figurastfm/Chapter4/Imagenes/configeditor.png}
			\caption{Configuration editor} \label{configeditor}
			\vspace{-1cm}
\end{figure}

Next, some of the most important concepts about COSMOS are briefly introduced and configured. \newpage

\begin{itemize} [noitemsep,topsep=0pt]
	\item \textbf{Target} 
\end{itemize}

	Device or spacecraft intended to establish a communication. As for this project, there will be only one target, the designed \glsname{cubesat}. Firstly, they have to be \textbf{declared}; this can be done in the \textbf{system.txt} file, under the following path.

	
	\begin{figure} [H]
\centering
\begin{minipage}{5cm}
\dirtree{%
.1 /GranaSAT-GroundStation.
.2 +config.
.3 system.
.4 system.txt.
}
\end{minipage}
\end{figure}



Besides, it is necessary to create a folder with the name of the target and replicate the following folder structure:

	\begin{figure} [H]
\centering
\begin{minipage}{5cm}
\dirtree{%
.1 /GranaSAT-GroundStation.
.2 +config.
.3 +targets.
.4 GRANASAT-I.
.5 +cmd-tlm.
.6 target.txt.
}
\end{minipage}
\end{figure}

The file \textit{target.txt} must contain the name of the database with the telemetry and telecommands information.


\begin{itemize} [noitemsep,topsep=0pt]
	\item \textbf{Interface}
\end{itemize}

In COSMOS, interfaces refer to the protocol used to communicate with a certain target, i.e., TCP/IP, serial or any other defined. Particularly, TCP/IP is already defined and is very convenient to use in a local environment such as the one used. They are defined in the file \textit{cmd-tlm-server.txt} under the following path.


	\begin{figure} [H]
\centering
\begin{minipage}{5cm}
\dirtree{%
.1 /GranaSAT-GroundStation.
.2 +config.
.3 +tools.
.4 +cmd-tlm-server.
.5 cmd-tlm-server.txt.
}
\end{minipage}
\end{figure}

\autoref{configinterf} illustrates the information needed in that file. Further information on the different options can be found on the official website \cite{cosmos}.

\begin{figure} [H]
			\centering
			\includegraphics[page=1,trim={0cm 0cm 0cm 0cm},clip=true,width=165mm]{figurastfm/Chapter4/Imagenes/interfaces.png}
			\caption{Interface configuration screenshot} \label{configinterf}
			%\vspace{-0.7cm}
\end{figure}

\subsubsection{Telecommand \& Telemetry database}

The whole framework revolves around the database with the telemetry and telecommand information. Once again, it can be defined using plain text; however, COSMOS is also compatible with \acrshort{XTCE} format, with an XML syntax. Once the implementation is decided, the data structure can be standardized using \acrshort{CCSDS}. Out of simplicity, in this Master's Thesis the database will be implemented using \acrshort{XTCE} and \acrshort{CCSDS}, but this one only partially. Expanding the database using completely that format is proposed as a future line of work, specially when working on a real mission.

There are countless XML editors, any of which can be used to create the needed database. It is proposed \textbf{Altova XMLSpy} because of the variety of different representations it allows. \autoref{codigoxml} shows an example of telemetry packets defined using \acrshort{XTCE}.

%xml code highlighting
\definecolor{gray}{rgb}{0.4,0.4,0.4}
\definecolor{darkblue}{rgb}{0.0,0.0,0.6}
\definecolor{cyan}{rgb}{0.0,0.6,0.6}
%%%%%%%%%%%%%%%%%
\lstdefinelanguage{XML}
{
  morestring=[b]",
  morestring=[s]{>}{<},
  morecomment=[s]{<?}{?>},
  stringstyle=\color{black},
  identifierstyle=\color{darkblue},
  keywordstyle=\color{cyan},
  morekeywords={xtce,encoding,UnitSSet}% list your attributes here
}
\lstset{  backgroundcolor=\color{white},  commentstyle=\color{codegreen}, keywordstyle=\color{magenta},   numberstyle=\tiny\color{codegray}, stringstyle=\color{codepurple},  basicstyle=\fontsize{4}\ttfamily\footnotesize,  breakatwhitespace=false,           breaklines=true,                   captionpos=b,                      keepspaces=true,                   numbers=left,                      numbersep=5pt,                    showspaces=false,                  showstringspaces=false,  showtabs=false,                    tabsize=2,  stepnumber=3, numberstyle=\tiny, basicstyle=\linespread{0.4}}

\begin{lstlisting}[language=XML, caption=Example of telemetry implemented using XTCE, label={codigoxml}]
<xtce:FloatParameterType name="OBC_Bosch_Temperature" shortDescription="OBC_Bosch_Temperature" signed="true">
				<xtce:UnitSet>
					<xtce:Unit description="deg">C</xtce:Unit>
				</xtce:UnitSet>
				<xtce:FloatDataEncoding sizeInBits="32" encoding="signed" bitOrder="LITTLE_ENDIAN">
					</xtce:FloatDataEncoding>
			</xtce:FloatParameterType>
			<xtce:FloatParameterType name="ESP_Data_Rate" shortDescription="ESP_Data_Rate" signed="true">
				<xtce:UnitSet>
					<xtce:Unit description="Mb/s">Mb/s</xtce:Unit>
				</xtce:UnitSet>
				<xtce:FloatDataEncoding sizeInBits="32" encoding="signed" bitOrder="LITTLE_ENDIAN">
					</xtce:FloatDataEncoding>
			</xtce:FloatParameterType>
			<xtce:FloatParameterType name="ESP_Signal_Level" shortDescription="ESP_Signal_Level" signed="true">
				<xtce:UnitSet>
					<xtce:Unit description="dBm">dBm</xtce:Unit>
				</xtce:UnitSet>
				<xtce:FloatDataEncoding sizeInBits="32" encoding="signed" bitOrder="LITTLE_ENDIAN">
					</xtce:FloatDataEncoding>

\end{lstlisting}

\autoref{schemaxml} shows part of the implemented database with one of the visual representations provided by Altova.

\begin{figure} [H]
			\centering
			\includegraphics[page=1,trim={0cm 0cm 0cm 0cm},clip=true,width=165mm]{figurastfm/Chapter4/Imagenes/xtce.png}
			\caption{Implemented telemetry and telecommand database using \acrshort{XTCE} format} \label{schemaxml}
			%\vspace{-0.7cm}
\end{figure}



\subsubsection{Interface}

Once the configuration parameters are set, COSMOS allows defining different user interfaces following the same procedure. For example, \autoref{maingui} displays the main interface of the developed \glsname{ground} software.

It shows the different tools developed in this Master's Thesis. It is divided into three differentiated areas: the first one, which initiates the controller, the second one 'Commanding and Scripting´, with the tools necessary to send commands to the \glsname{cubesat} and the last one 'Telemetry´, with tools to visualize and post-process the telemetry. Some of them are briefly described next.

\begin{figure} [h]
			\centering
			\includegraphics[page=1,trim={0cm 0cm 0cm 0cm},clip=true,width=85mm]{figurastfm/Chapter4/Imagenes/gui.png}
			\caption{Main interface of the \glsname{GranaSAT} \glsname{ground} controller} \label{maingui}
			\vspace{0.5cm}
\end{figure}

\begin{itemize} [noitemsep]
	\item \textbf{Command and Telemetry Server:} Used to connect to the different targets. It shows information regarding the number of packages received, errors detected and so forth.
\end{itemize}


\begin{itemize} [noitemsep,topsep=0pt]
	\item \textbf{Replay:} Used to \textit{repeat} a past mission. By loading the log of a certain session, it allows repeating the processing performed in real time, plot the information again with a temporal reference, etc. \autoref{replay} shows the interface of this tool.
\end{itemize}

\begin{figure} [H]
			\centering
			\includegraphics[page=1,trim={0cm 0cm 0cm 0cm},clip=true,width=85mm]{figurastfm/Chapter4/Imagenes/replay.png}
			\caption{Replay tool} \label{replay}
			%\vspace{-0.7cm}
\end{figure}


\begin{itemize} [noitemsep,topsep=0pt]
	\item \textbf{Packet viewer:} Used to visualize the telemetry data as received, either in real time or using the \textbf{Replay} tool. \autoref{viewer} shows an example of received packet in this tool.
\end{itemize}

\begin{figure} [H]
			\centering
			\includegraphics[page=1,trim={0cm 0cm 0cm 0cm},clip=true,width=85mm]{figurastfm/Chapter4/Imagenes/packetview.png}
			\caption{Packet viewer} \label{viewer}
			%\vspace{-0.7cm}
\end{figure}

Telemetry Grapher is not described because it will be widely seen in \autoref{chap:chapter5} to verify the functioning of the system.

\subsection{Real time 3D viewer}

Besides the COSMOS environment developed and explained in the latter section, it has been developed a web-based real time 3D viewer. Once again, it is intended to be in a local environment; in this case the reason is obvious: the communication delay existing in a real in orbit communication makes it impossible a real time display. However, it is an attractive tool for students learning about CubeSats with this platform. 

The diagram of \autoref{3dview} depicts the functioning and technologies used in the viewer.

\begin{figure} [H]
			\centering
			\includegraphics[page=1,trim={0cm 0cm 0cm 0cm},clip=true,width=165mm]{figurastfm/Chapter4/PDF/3dviewer.pdf}
			\caption{3D viewer architecture} \label{3dview}
			%\vspace{-0.7cm}
\end{figure}

The proper \glsname{cubesat} hosts in the \acrshort{OBC}, along with the rest of components (\acrshort{OBDH}, Flight Software...), a \textbf{Node.js} \cite{nodejs} server. When the attitude-related sensors (accelerometer, gyroscope and magnetometer) have data available, they pass it directly to the web server. Once the server has the attitude information, it communicates with the user interface using secure \glsname{websockets}. The user interface is designed as a \glsname{SPAG}~(\acrshort{SPA}) so it is fully downloaded just once. On the other hand, the server is listening at port 8000, waiting for connections.

The designed SPA uses three technologies to plot a real time 3D view of the \glsname{cubesat}. The one in the bottom is \textbf{WebGL} \cite{webgl}, graphic libraries optimized for web applications, in charge of rendering the 3D visualization. Right above it, \textbf{D3.js} \cite{d3} which adds an abstraction layer over WebGL, easing the process by managing \acrshort{STL} models, controlling light, cameras, background, etc; in sum, it is a 3D graphic library. Finally, on the top of the stack,  the library \textbf{C3.js} \cite{c3js} is used to generate real time 2D plots.

The website can be seen in \autoref{plotter3d}.

\begin{figure} [H]
			\centering
			\includegraphics[page=1,trim={0cm 0cm 0cm 0cm},clip=true,width=165mm]{figurastfm/Chapter4/PDF/3dplotter_.pdf}
			\caption{3D viewer web application} \label{plotter3d}
			%\vspace{-0.7cm}
\end{figure}

\section{Simulation CubeSat}

The design phase finalizes with the simulation \glsname{cubesat}, which will lay the foundations to complete a mission-ready device in the future. This section addresses all the aspects around \glsname{cubesat} design, from the mechanical design of the structure to the electronics, which makes it \textbf{the most complex part of the whole platform}.

\subsection{Mechanical structure}

As described in \autoref{cubesatdefin}, one of the most extended \glsname{cubesat} standards was defined by the California Politechnical State University \cite{calpoly}, depicted in \autoref{cubesatstd}. Therefore, the design performed in this Master's Thesis will stick as much as possible to it.

\subsubsection{Design}

Once again, the mechanical structure is designed using \glsname{solid}. \autoref{drawings} shows the drawings of the design. On the other hand, \autoref{cubesatrender} shows another high quality 3D render.


\begin{figure} [H]
			\centering
			\includegraphics[page=1,trim={0cm 0cm 0cm 0cm},clip=true,width=135mm]{figurastfm/Chapter4/Imagenes/csat3d.png}
			\caption{Mechanical structure 3D render} \label{cubesatrender}
			%\vspace{-1cm}
\end{figure}

Besides the outer edge design, the non-standard inner is important because it will constrain the \acrshort{PCB} edge.

\begin{landscape}
\begin{figure} [H]
			\centering
			\includegraphics[page=1,trim={0cm 0cm 0cm 0cm},clip=true,width=220mm]{figurastfm/Chapter4/PDF/soliddrw.pdf}
			\caption{Mechanical structure design} \label{drawings}
			%\vspace{-0.7cm}
\end{figure}
\end{landscape}

\subsubsection{Manufacturing}

The structure is made of aluminum and manufactured using a milling machine (see \autoref{milling}) in our own laboratory. 

\begin{figure} [H]
			\centering
			\includegraphics[page=1,trim={0cm 0cm 0cm 0cm},clip=true,width=80mm]{figurastfm/Chapter4/Imagenes/milling.jpg}
			\caption{Aluminum milling machine} \label{milling}
			\vspace{-0.5cm}
\end{figure}

After an iterative process, an optimized structure is accomplished, displayed in \autoref{manufactu}. Some testing solar panels \acrshort{PCB} have been attached in order to check the correctness of the design.

			\begin{figure}[H]
			\centering
			\subfloat[1]{\includegraphics[page=1,trim={0cm 0cm 0cm 0cm},clip=true,width=83mm]{figurastfm/Chapter4/Imagenes/manuf1.jpg}}
			\quad
			\subfloat[2]{\includegraphics[page=1,trim={0cm 0cm 0cm 0cm},clip=true,width=75mm]{figurastfm/Chapter4/Imagenes/manuf2.jpg}}
			\caption{Mechanical structure manufactured} \label{manufactu}
			\vspace{-2cm}
\end{figure}


\subsection{On-board computer (OBC)} \label{obcdesig}

As stated in \autoref{chap:chapter3}, the \acrshort{OBC} is probably the most complex part of the \glsname{cubesat}. It is in charge of managing and synchronizing the rest of the system \textbf{in-time} and acting consequently. This section is divided into different subsections according to the main functioning modules of the \acrshort{OBC}. 

\subsubsection{Central Processing Unit}

The Central Processing Unit of a \glsname{cubesat}, typically abbreviated as \acrshort{CPU}, can indeed be understood in the same manner of a standard computer. It is some kind of circuitry, typically a microprocessor, which performs certain instructions consisting in arithmetic, I/O operations, logic, etc.

This task can be addressed by a variety of architectures. One of the most common options is using a \textbf{microcontroller}, for example the MSP430 by Texas Instruments \cite{msp}. However, in this project the design will go one step beyond using a so-called \linebreak micro-computer or more specifically an embedded \textbf{Single-Board Computer} (\acrshort{SBC}). They are a complete computer built on a single \acrshort{PCB}, including not only microprocessor, but also memory, RAM, or I/O capabilities. Unlike microcontrollers, \acrshort{SBC} allows storage memory and can run an operating system completely. On the other hand, they can be more difficult to implement.  

\acrshort{SBC} can be found in a variety of formats; for instance, \autoref{sbcformats} shows a couple of examples in \textit{standalone} format.

			\begin{figure}[H]
			\centering
			\subfloat[Raspberry Pi 4 B+]{\includegraphics[page=1,trim={0cm 0cm 0cm 0cm},clip=true,width=83mm]{figurastfm/Chapter4/Imagenes/rpi.jpg}}
			\quad
			\subfloat[Orange Pi]{\includegraphics[page=1,trim={0cm 0cm 0cm 0cm},clip=true,width=75mm]{figurastfm/Chapter4/Imagenes/orangepi.jpg}}
			\caption{Different \acrshort{SBC} available on the market} \label{sbcformats}
			%\vspace{-2cm}
\end{figure}



\textit{Standalone} \acrshort{SBC} just need a loaded operating system and a power supply. Although that format is convenient for hobby or \textit{domestic} projects, as for a \glsname{cubesat} project is too heavy and furthermore really difficult to fit into the structure. For these reason, this design will make use of \textbf{industrial} \acrshort{SBC}, characterized by a more flexible form  factor with standard DDR2 SODIMM connector. It  increases substantially the complexity of the design but also allows a much more professional result and gets the project closer to the ticket-to-fly in the \textit{Fly Your Satellite!} program, described in \ref{fys}.

\paragraph{Comparative}

In order to choose the \acrshort{SBC} which fits best the requirements needed, it is performed a comparative between different options available in the market. The comparative is also needed because of the cost of this part, which is usually the highest of the whole \glsname{cubesat}.  Normally, this is also the component with the greatest power consumption, therefore an extensive analysis is performed on the option chosen.

\subparagraph{Apalis TK1}

The Apalis TK1 by Toradex \cite{apalistk1} is based on the NVIDIA Tegra K1 \acrshort{SoC}. It counts with a Cortex A15 quad-core CPU up to 2.1 GHz and a powerful GPI Geforce Kepler, also from NVIDIA, as well as 2 GB RAM.

Additionally, the Apalis TK1 features a low power ARM Cortex M4 up to 100 MHz which extends the \acrshort{ADC}, \acrshort{GPIO} and several other interfaces. It fits into a standard 200 pins SODIMM connection. \autoref{tk1} shows this \acrshort{SBC}.


\begin{figure} [H]
			\centering
			\includegraphics[page=1,trim={0cm 0cm 0cm 0cm},clip=true,width=80mm]{figurastfm/Chapter4/Imagenes/tk1.jpg}
			\caption{Apalis TK1} \label{tk1}
			%\vspace{-1cm}
\end{figure}


\subparagraph{Raspberry Pi 3 Compute Module}

The second option is one of the most known in the \acrshort{SBC} field: the Raspberry Pi 3 Compute Module. It is the industrial equivalent to the standard Raspberry Pi displayed in \autoref{sbcformats}; therefore, it is based on the same \acrshort{SoC}, the Broadcom BCM2837 up to 1.2 GHz and features 1 GB RAM, all integrated on a 67.6 mm x 31 mm board with a SODIMM connector. As for storage capabilities, it allows both 4 GB integrated flash memory and \acrshort{SD} card interface.

One of the greatest advantages of this solution is the amount of drivers already available when Linux is used as \acrshort{OS}, because of the broad usage of this board. \autoref{rpicm} shows the \acrshort{PCB}.


\begin{figure} [H]
			\centering
			\includegraphics[page=1,trim={0cm 0cm 0cm 0cm},clip=true,width=80mm]{figurastfm/Chapter4/Imagenes/rpicm.jpg}
			\caption{Raspberry Pi 3 Compute Module \cite{cm3}} \label{rpicm}
			%\vspace{-1cm}
\end{figure}


\subparagraph{Colibri VF50}


This solution is provided by Toradex \cite{colibrivf50} and is based on the NXP Freescale Vybrid \acrshort{SoC}, featuring a microprocessor Cortex A5 at 400 MHz which delivers cost effective processing and graphic performance. It counts with 128 MB RAM and 128 MB of storage.

Its major advantage is its low power consumption and value for money, however its storage capabilities are not enough for this design. \autoref{colibri} shows it.

\begin{figure} [H]
			\centering
			\includegraphics[page=1,trim={0cm 0cm 0cm 0cm},clip=true,width=80mm]{figurastfm/Chapter4/Imagenes/COLIBRI.jpg}
			\caption{Colibri VF50 \cite{colibrivf50}} \label{colibri}
			%\vspace{-1cm}
\end{figure}

\subparagraph{CL-SOM-iMX8X}

Based on the i.MX8X processor family by NXP, features a quad-core ARM Cortex-A35 up to 1.2 GHz and an integrated GPU. It also counts with 4 GB RAM and WiFi 802.11ac capabilities. Another interesting point is its wide temperature range, -40 \textdegree C to 85 \textdegree C.

Regarding connectivity, it counts with 96 \acrshort{GPIO}, 4 \acrshort{UART} and is \acrshort{USB} 3.0 ready. It is shown in \autoref{cls}.

\begin{figure} [H]
			\centering
			\includegraphics[page=1,trim={0cm 0cm 0cm 0cm},clip=true,width=80mm]{figurastfm/Chapter4/Imagenes/cls.jpg}
			\caption{CL-SOM-iMX8X} \label{cls}
			%\vspace{-1cm}
\end{figure}


\autoref{tablacm} sums up the main features of each solution.
% Please add the following required packages to your document preamble:
% \usepackage{booktabs}
\begin{table}[H]
\centering
\begin{tabular}{@{}cccccc@{}}
\toprule
\textbf{SBC}      & \textbf{System-On-Chip} & \multicolumn{1}{l}{\textbf{RAM (GB)}} & \multicolumn{1}{l}{\textbf{Flash (GB)}} & \multicolumn{1}{l}{\textbf{GPIO}} & \multicolumn{1}{l}{\textbf{Price (€)}} \\ \midrule
Apalis TK1        & NVIDIA Tegra K1         & 2                                     & 16                                      & 87                                & 211.25                                 \\ \midrule
Raspberry Pi 3 CM & Broadcom BCM2837        & 1                                     & 4                                       & 48                                & 30.30                                  \\ \midrule
Colibri VF50      & NXP Freescale Vybrid    & 0.128                                 & 0.128                                   & 103                               & 40.40                                  \\ \midrule
CL-SOM-iMX8X      & ARM Cortex A-35         & 4                                     & 4                                       & 96                                & 75                                     \\ \bottomrule
\end{tabular}
\caption{SBC comparative}
\label{tablacm}
\end{table}

The first option, Apalis TK1 has a powerful GPU which is not likely to be used for the \glsname{cubesat} purposes; therefore, taking into account its high cost, it is not a viable solution. On the hand, the Colibri VF50 is not powerful enough, although it counts with the highest number of \acrshort{GPIO} it lacks RAM and storage; the Raspberry Pi 3 CM exhibits better features at a lower price. Among the two possibilities left, besides a better value for money, the Raspberry solution, as mentioned, counts with one of the greatest communities because of its extended usage. For all these reasons, it is the option chosen.

\paragraph{Raspberry Pi 3 CM consumption analysis} \label{rpipower}

As described before, the \acrshort{CPU} is one of the most power consuming components of the \glsname{cubesat}. Therefore, it is necessary to perform a consumption analysis in different situations which allows characterizing the device completely, so the \acrshort{EPS} can be adequately designed, increasing reliability. 

To perform this analysis, different performance tests will be executed using the Raspberry Pi 3 CM. To ease these tests, the operating system \textbf{DietPi} \cite{diet} is used; it is a lightweight system which already includes different stress tests. By executing those tests, the device will increase its consumption until different ranges, including maximum before powering itself off. 

Regarding the data gathering, it is used the DC Power Analyzer Keysight N6705B, shown in \autoref{dcana}, by establishin a \acrshort{TCP} connection over a \acrshort{LAN} network, using \glsname{matlab}.

\begin{figure} [H]
			\centering
			\includegraphics[page=1,trim={0cm 0cm 0cm 0cm},clip=true,width=90mm]{figurastfm/Chapter4/Imagenes/keysight.png}
			\caption{Keysight N6705B \cite{agilent}} \label{dcana}
			%\vspace{-1cm}
\end{figure}



The structure of the tests will be the following: as the CM allows being supplied with two voltages (3.3 V and 5 V), both must be tested in order to decide. On the other hand, DietPi allows choosing the number of cores of the microprocessor used (1 to 4) as well as the CPU \textbf{governor} (\textit{ondemand, conservative, powersave or perfomance}). With all these parameters, the CM will be analyzed under three situations: at the boot (power peaks are common), in stationary state and when dealing with high CPU loads (maximum consumption).

Of course, this schema has produced a large number of plots; out of simplicity, only a selection of the most representative ones are plotted. Particularly, it is compared the power consumption for each voltage supply in every three situations. Additionally, they are compared twice: with the \textit{Powersave} governor and 1 core enabled (minimum consumption) and with \textit{Perfomance} governor and 4 core enabled (maximum consumption). This setting-up will allow determining the whole consumption range, from the minimum to the maximum.


\subparagraph{\textit{Powersave} governor - 1 core}



\begin{figure} [H]
			\centering
			\includegraphics[page=1,trim={0cm 0cm 0cm 0cm},clip=true,width=135mm]{figurastfm/Chapter4/PDF/powersboot.pdf}
			\caption{Power Consumption at \textbf{Boot}} %\label{dcana}
			\vspace{-1cm}
\end{figure}

\begin{figure} [H]
			\centering
			\includegraphics[page=1,trim={0cm 0cm 0cm 0cm},clip=true,width=135mm]{figurastfm/Chapter4/PDF/powersstat.pdf}
			\caption{Power Consumption at \textbf{Stationary State}} % \label{dcana}
			\vspace{-2cm}
\end{figure}


\begin{figure} [H]
			\centering
			\includegraphics[page=1,trim={0cm 0cm 0cm 0cm},clip=true,width=135mm]{figurastfm/Chapter4/PDF/powershigh.pdf}
			\caption{Power Consumption at \textbf{High load}} % \label{dcana}
			\vspace{-0.5cm}
\end{figure}

\subparagraph{\textit{Performance} governor - 4 cores}


\begin{figure} [H]
			\centering
			\includegraphics[page=1,trim={0cm 0cm 0cm 0cm},clip=true,width=135mm]{figurastfm/Chapter4/PDF/performboot.pdf}
			\caption{Power Consumption at \textbf{Boot}} %\label{dcana}
			\vspace{-2cm}
\end{figure}

\begin{figure} [H]
			\centering
			\includegraphics[page=1,trim={0cm 0cm 0cm 0cm},clip=true,width=135mm]{figurastfm/Chapter4/PDF/performstat.pdf}
			\caption{Power Consumption at \textbf{Stationary State}} % \label{dcana}
			%\vspace{-1cm}
\end{figure}


\begin{figure} [H]
			\centering
			\includegraphics[page=1,trim={0cm 0cm 0cm 0cm},clip=true,width=135mm]{figurastfm/Chapter4/PDF/performhigh.pdf}
			\caption{Power Consumption at \textbf{High load}} % \label{dcana}
			\vspace{-2cm}
\end{figure}

From the previous plots, some conclusions can be extracted:

\begin{itemize} [topsep=0pt]
	
	\item At \textbf{boot}, the \textbf{minimum} power consumption peak (with \textit{Powersave} governor and 1~core enabled) is \textbf{2.25 W, with 5 V supply}.
	\item At \textbf{boot}, the \textbf{maximum} power consumption peak (with \textit{Performance} governor and 4 core enabled) is \textbf{2.5 W}, regardless of the voltage supply; 11.1 \% higher than the minimum.
	\item At \textbf{stationary state}, power consumption is higher in the \textit{Performance} configuration, regardless of the voltage supply (0.49 W vs 0.61 W with 3.3 V supply and 0.62 W vs 0.78 W with 5 V supply).
	\item When dealing with a \textbf{high work load}, the power consumption is again higher with 5 V supply (1.5 W and 3.5 W with \textit{Powersave} and \textit{Performance} configuration, respectively) than with 3.3 V supply (1.2 W and 3.25 W with \textit{Powersave} and \textit{Performance} configuration, respectively). This case features the greatest increase in power consumption between the voltages supply: \textbf{170.83 \%} and \textbf{133.33 \%} increase with 3.3 V and 5 V respectively.
	\item Considering equal configuration, \textbf{power consumption with 3.3 V supply is always lower} than the equivalent with 5 V supply.
	
\end{itemize}

Therefore, in order to save the most energy possible (given that the system will be supplied by batteries) 3.3 V supply is chosen. All this information will be used in \ref{battscharac} to estimate the operating time of the \acrshort{OBC} depending on the battery used.


\subsubsection{Communications}

As described in previous sections, the Communications subsystem is included within the \acrshort{OBC}. As for this part of the \glsname{cubesat}, only the simulation purposes are considered. In a real mission, COMMS area will count with different radio-frequency systems, antennas and the rest of elements mentioned in section \ref{radiofr}. Once  again, that task can suppose a complete project by itself, so in this section, the focus will be on communications system which allows simulation-related tasks such as programming the device, downloading data, connecting to the Internet and so on.
\paragraph{Wired}

Wired connections will be the easiest way of communication when working with the platform. Although Wireless communications are intented to be used under operation, some tasks need to be performed over a wire, for example, the first programming.

\subparagraph{USB \& Ethernet}

\acrshort{USB} is used to first program the device, perform low-level tasks and \textbf{charge} the batteries. This is important considering the poor performance expected from the solar panels, as described in \autoref{ledsim}; although that charging procedure can be used to test the \acrshort{EPS}, when the simulations are related to another subsystems there must be some quick charge method, in this case provided through \acrshort	{USB}.  Several connections will be disposed, allowing not only the mentioned tasks but also connecting external devices for prototyping, such as cameras. 

Both \acrshort{USB} and Ethernet connections are complex to design as they are \textbf{differential transmission lines}. Differential signaling is a technique to transmit information using complementary signals (positive and negative) using different wires, in contrast to the more common \textbf{single-ended signaling}. \autoref{diffsignal} shows a couple of differential lines with a representation of each signal when traveling across each of the lines.

\begin{figure} [H]
			\centering
			\includegraphics[page=1,trim={0cm 0cm 0cm 0cm},clip=true,width=105mm]{figurastfm/Chapter4/PDF/diff_Signaling.pdf}
			\caption{A differential pair with a signal propagating on each line \cite{bogatin}} \label{diffsignal}
			%\vspace{-2cm}
\end{figure}


In an impedance-matched circuit, external \acrshort{EMI} affects both wires equally. Since the receiver only detects the difference between the conductors, this technique resists electromagnetic noise better than a single-ended conductor, as well as improving \acrshort{SNR} and minimizing crosstalk \cite{wiki}. \autoref{diffsch} depicts the advantages of this method.


\begin{figure} [H]
			\centering
			\includegraphics[page=1,trim={0cm 0cm 0cm 0cm},clip=true,width=115mm]{figurastfm/Chapter4/Imagenes/Diff.png}
			\caption{Differential signaling functioning \cite{wiki}} \label{diffsch}
			%\vspace{-2cm}
\end{figure}

In order to design a \acrshort{USB} line, it is necessary a \textbf{Host Controller Interface} \acrshort{HCI} or simply a controller. It is an integrated circuit which allows the computer to communicate with the \acrshort{USB} device. Indeed, Ethernet connections need another too. That is the reason for this common section: it will be used a common controller for both, particularly the \textbf{Microchip LAN9514} \cite{lan9514}. It is a dual \acrshort{USB} 2.0 and 10/100~Ethernet (full-duplex) Controller which features four \acrshort{USB} downstream ports and one upstream. \autoref{laninten} depicts the internal block diagram of the chip.


\begin{figure} [H]
			\centering
			\includegraphics[page=1,trim={0cm 0cm 0cm 0cm},clip=true,width=125mm]{figurastfm/Chapter4/PDF/lan9514_.pdf}
			\caption{LAN9514 internal block diagram \cite{lan9514}} \label{laninten}
			%\vspace{-2cm}
\end{figure}

This common usage has drawbacks though. First, the Raspberry Pi 3 Compute Module has a single connection point (which is actually composed of two differential pins) for both, Ethernet and \acrshort{USB} connections. However, if the controller is directly connected to the \acrshort{SBC}, the Compute Module cannot be turned into \textbf{slave mode}, which means that \textbf{it cannot be programmed}. In this context, \textit{programmed} means mounting the memory storage of the Compute Module as \acrshort{USB} slave and then \textbf{flash} it with an operating system image. Therefore, it is necessary a \textbf{\acrshort{USB} switch} that allows both modes depending on the connector used and the configuration selected. For this purpose, it is chosen the integrated circuit \textbf{ON Semiconductor FSUSB42} \cite{onsemi}; it is a bi-directional two \acrshort{USB} 2.0 ports switch in charge of turning the Compute Module into slave mode when a \acrshort{USB} is connected to the programming port. \autoref{swusb} depicts the functioning of this circuitry. 

\begin{figure} [H]
			\centering
			\includegraphics[page=1,trim={0cm 0cm 0cm 0cm},clip=true,width=165mm]{figurastfm/Chapter4/PDF/swusb.pdf}
			\caption{LAN9514 and \acrshort{USB} switch mechanism} \label{swusb}
			%\vspace{-2cm}
\end{figure}
Ethernet connections are designed for a differential impedance of \SI{100}{\text{ \ohm}}. Differential impedance depends on a variety of factors, featuring: track width, differential tracks separation, substrate used, substrate height and frequency. This kind of design is a complex issue which usually needs to assume different tradeoffs. There are a variety of methods to compute differential impedance; in short, every couple of differential impedance exhibits two different impedances: on the one hand, the differential one and, on the other hand, the single-ended corresponding to each of the lines (i.e., the \textbf{characteristic impedance of the line}). The equations under differential impedance are experimental, as there is no exact solution. Some of the most typically used are shown in \autoref{diffeqs}, where $\epsilon_\text{r}$ is the \glsname{permittivity}, $h$ is the substrate's height and $w$ and $t$ the width and thickness of the line, respectively. It has been used the values given by the manufacturer which will be in charge to produce the \acrshort{PCB}.

\begin{equation}
\begin{split}
Z_{\text{0}} &= \frac{87}{\sqrt{\epsilon_{\text{r}}+1.41}} \ln \frac{5.98h}{0.8w+t} \\
Z_{\text{diff}} &= 2\times Z_{\text{0}} (1-0.48\mathrm{e}^{-0.96\frac{d}{h}})\\
\end{split}\label{diffeqs}
\end{equation}


To illustrate the complexity under this operation, the differential impedance corresponding to an Ethernet line is graphically calculated, using the equations above. Particularly, \autoref{difftrace} depicts the differential impedance depending on the trace width, for three different space lengths between the lines. The first noticeable thing is that differential impedance decreases as trace width increases. The horizontal discontinuous lines depicts the target impedance for both, differential and single-ended, in the Ethernet case \textbf{100 $\Omega$} and \textbf{50 $\Omega$}, respectively. 
 
\begin{figure} [H]
			\centering
			\includegraphics[page=1,trim={0cm 0cm 0cm 0cm},clip=true,width=145mm]{figurastfm/Chapter4/PDF/diff_trace.pdf}
			\caption{Graphical calculation of differential impedance} \label{difftrace}
			\vspace{-2cm}
\end{figure}

The differential impedance varies with the spacing between lines, while the single-ended is obviously independent of this parameter. By plotting vertical lines (in discontinuous black) at the point in which each series matches its target is easy to see that \textbf{it is not possible to simultaneously have 100 $\Omega$ as differential impedance and 50~$\Omega$ as single-ended}. A certain \textbf{tradeoff} has to be assumed. In this case, the closest option is using traces of \textbf{10.8 mils} width with a separation of \textbf{14 mils} (purple series). With that design, single-ended impedance would be about 55 $\Omega$ while differential would be about 100 $\Omega$. 

While the latter example helps to understand the complexity under this kind of designs, the calculation is usually performed using dedicated and more accurate tools such as \textbf{Saturn PCB Design Toolkit} \cite{saturn}. \autoref{saturnpic} shows the calculation using that tool. It allows checking that the result accomplished by analytical calculation before is really close to the one estimated by Saturn, \textbf{105 $\Omega$}/\textbf{59 $\Omega$}, which falls into a 10 \% tolerance, enough for the design to work properly.


\begin{figure} [H]
			\centering
			\includegraphics[page=1,trim={0cm 0cm 0cm 0cm},clip=true,width=165mm]{figurastfm/Chapter4/Imagenes/saturn.png}
			\caption{Calculation of differential impedance using Saturn PCB Design Toolkit} \label{saturnpic}
			\vspace{-2cm}
\end{figure}

In order for the signals to reach the end of the line at the same time, it is important for the lines to be \textbf{length matched}. It is accomplished by tweaking the lines as shown in \autoref{ethernetlines}.

\begin{figure} [H]
			\centering
			\includegraphics[page=1,trim={0cm 0cm 0cm 0cm},clip=true,width=120mm]{figurastfm/Chapter4/PDF/eth_pair.pdf}
			\caption{Lines tweaking to achieve length matching} \label{ethernetlines}
			%\vspace{-2cm}
\end{figure}

Just like Ethernet lines, \acrshort{USB} designs must be based on differential signaling, particularly with a 90 $\Omega$ differential impedance and 45 $\Omega$ single-ended, which is accomplished using the same procedure.

Besides the embedded Ethernet design described above, in the \acrshort{OBC} design it is included a ENC28J60 \cite{enc28j60break} Ethernet breakout which provides with direct Ethernet connection by \acrshort{SPI}, shown in Figure\autoref{enc}. Its obvious drawback is its size, as seen in Figure \autoref{encobc}, with the breakout on, the \acrshort{PCB} cannot be fitted into the \glsname{cubesat} structure; also the data rate accomplished with it is significantly inferior. However, in case the embedded design fails, this could be used as \textbf{contingency option}.

			\begin{figure}[H]
			\centering
			\subfloat[Breakout\label{enc}]{\includegraphics[page=1,trim={0cm 0cm 0cm 0cm},clip=true,width=40mm]{figurastfm/Chapter4/Imagenes/enc.jpg}}
			\quad
			\subfloat[Placed in the \acrshort{OBC}, impeding its placement in the \glsname{cubesat}\label{encobc}]{\includegraphics[page=1,trim={0cm 0cm 0cm 0cm},clip=true,width=90mm]{figurastfm/Chapter4/PDF/enc.pdf}}
			\caption{ENC28J60} \label{sbcformats}
			%\vspace{-2cm}
\end{figure}

\paragraph{Wireless}

The formal requirement \textbf{OBC.FoR.15} described in \autoref{forobc}, states the need for a wireless communication system, provided by an integrated device fully compatible with \acrshort{IEEE} 802.11g standard. After designing the wired communications subsystem, this section addresses the wireless ones, specially important in order to allow a realistic simulation using the \acrshort{I2DOS}: wired communications allow neither a free rotation nor a convenient telemetry sending. Therefore, it will be crucial to accomplish a reliable communication over the air.

\subparagraph{WiFi}
Another communications related requirement was wireless communications. It will be provided with the typical technology family \textbf{WiFi}, based on the \acrshort{IEEE}~802.11~standards.

In order to include this feature, it is used the \textbf{ESP8266} microcontroller by Espressif~\cite{espressif}. It is a highly integrated Wi-Fi programmable \acrshort{SoC} solution, really power-efficient with an output power of up to 17.5 dBm. It will be used as slave to the Raspberry~Pi~3~Compute~Module with which will be completely integrated as a network interface. Figure\autoref{espfoto} shows the module while Figure\autoref{espblocks} depicts the functional block diagram.




\begin{figure}[H]
			\centering
			\subfloat[External aspect\label{espfoto}]{\includegraphics[page=1,trim={0cm 0cm 0cm 0cm},clip=true,width=55mm]{figurastfm/Chapter4/Imagenes/esp.jpg}}
			\\
			\subfloat[Functional block diagram\label{espblocks}]{\includegraphics[page=1,trim={0cm 0cm 0cm 0cm},clip=true,width=165mm]{figurastfm/Chapter4/PDF/esp.pdf}}
			\caption{ESP8266 \cite{espressif}} 
			%\vspace{-2cm}

\end{figure}

\subsubsection{Payload. Sensors}

As defined in previous sections, the payload is the \textbf{useful equipment} carried by the \glsname{cubesat}, able to provide with valuable information. Besides dedicated payloads such as cameras, also sensors and equipment necessary for the correct functioning of the \glsname{cubesat} itself  can be considered as a payload, from \glsname{startracker} to magnetometers or similar. This Master's Thesis provides with the basis for a simulation \glsname{cubesat} which includes the majority of the sensors needed to successfully perform a mission, allowing adding more elements as payload in the future. In this section, the design of the different sensors implemented is outlined. Although most of them are used by the \acrshort{ADCS} in its functioning, they can be considered part of the \acrshort{OBC} and will be placed on its board.


\paragraph{Inertial Measurement Unit}


\acrshort{IMU} have been previously defined in this project. It is in charge of measuring different parameters such as gravitational acceleration, angular rate and orientation of the device, by using accelerometers, gyroscopes and magnetometers. This design will use the \textbf{MPU9250} by Invensense \cite{mpu9250}.


%\begin{wrapfigure}{r}{0.18\textwidth} 
	%\centering
	%\includegraphics[width=0.13\textwidth]{figurastfm/chapter4/imagenes/mpu.png}
	%\vspace{-0.1cm}
	%\caption{MPU9250} %\label{MPU9250PIC}
%\end{wrapfigure}

\begin{figure} [H]
			\centering
			\includegraphics[page=1,trim={0cm 0cm 0cm 0cm},clip=true,width=40mm]{figurastfm/Chapter4/imagenes/mpu.png}
			\caption{MPU9250} \label{mpupack}
			\vspace{-0.5cm}
\end{figure}


MPU9250  is a multi-chip module consisting of two dies integrated into a single package, displayed in \autoref{mpupack}; while the first hosts the 3-Axis gyroscope and the 3-Axis accelerometer, the second houses the 3-Axis magnetometer, resulting in a 9-axis MotionTracking
device that combines a 3-axis gyroscope, 3-axis accelerometer, 3-axis magnetometer in a small 3x3x1 mm package, challenging to solder. It allows communicating using both \acrshort{I2C} and \acrshort{SPI} and features nine 16-bit \acrshort{ADC} to digitize the output of each sensor.

In order to gather the data from the sensors, it is implemented a Python library in which the different registers addressed are assigned a variable name, as well as the default \acrshort{I2C} slave addresses. Also, the different methods to read from the sensors are coded. For example, \autoref{coderead} shows the method to read the accelerometer data, using pseudocode.

\begin{lstlisting}[language=Python, caption=Method to read from the accelerometer sensors, label=coderead]
    def readAccelerometer(self):
 
        x = round(data[0]*resolution, 3)
        y = round(data[1]*resolution, 3)
        z = round(data[2]*resolution, 3)

        return {x,y,z}
\end{lstlisting}

A similar procedure is followed for the rest of the sensors, which eases accessing to the data from the main script.


\paragraph{Barometer \& Thermometer}

Specially useful for local simulation purposes, implementing a sensor easy to understand such as this gets students closer to the device and increases its interest. In this case, the integrated circuit used is the tiny Bosch BMP280 \cite{bosch}, shown in \autoref{bmp}.

\begin{figure} [H]
			\centering
			\includegraphics[page=1,trim={0cm 0cm 0cm 0cm},clip=true,width=55mm]{figurastfm/Chapter4/Imagenes/BMP.jpg}
			\caption{Bosch BMP280 \cite{bosch}} \label{bmp}
			\vspace{-0.5cm}
\end{figure}



Bosch BMP280 exhibits high accuracy and linearity, as well as an almost negligible power consumption (with an average current of \SI{2.74}{\micro A}. Besides, it allows a wide range of voltage supply, between 1.71 V and 3.6 V. Its technical parameters comply with the requirements defined in \autoref{chap:chapter3}.

Again, it is developed a Python library to  read from the integrated sensors, thermometer and barometer. This chip is also able to estimate the altitude above sea level from the pressure detected. Also, it is able to be programmed with different over-sampling settings for both barometer and temperature sensor; each oversampling step reduces noise and increases the output resolution by one bit. \autoref{bmppcb} shows the BMP280 placed in the \acrshort{OBC} \acrshort{PCB}.

\begin{figure} [H]
			\centering
			\includegraphics[page=1,trim={0cm 0cm 0cm 0cm},clip=true,width=100mm]{figurastfm/Chapter4/PDF/bosch.pdf}
			\caption{Bosch BMP280 placed in the \acrshort{OBC} \acrshort{PCB}} \label{bmppcb}
			\vspace{-2cm}
\end{figure}


\paragraph{Sun sensors. ADC}

In order to determine the position of the real or simulated Sun, it is necessary to include the so-called \textbf{Sun sensors}. As for every component implied in a space-oriented design, there are qualified Sun sensors, intended to be used in a real mission, \autoref{ssensor} shows an example from \cite{csatshop}.


\begin{figure} [H]
			\centering
			\includegraphics[page=1,trim={0cm 0cm 0cm 0cm},clip=true,width=70mm]{figurastfm/Chapter4/Imagenes/ssensor.png}
			\caption{NCSS-SA05 Sun Sensor} \label{ssensor}
				%\vspace{-2cm}
\end{figure}

As usual, the \textit{problem} with space qualified components is the cost: the NCSS-SA05 costs around 3300 €. That value makes it unfeasible for this project so, instead, the Sun sensors used will be based on \acrshort{LDR}. 

\acrshort{LDR} will make it possible, depending on the intensity measured, to determine with enough accuracy the position of the light. Given that \acrshort{LDR} are ultimately resistances, it is necessary to arrange some circuitry to determine the Sun's position. The easiest way to accomplish that point is designing a voltage divider between the \acrshort{LDR} itself and another resistor. When supplied, a certain voltage falls in the measurement point, which is gathered using an \acrshort{ADC}. \autoref{ldrcircuit} depicts the circuit described.


\begin{figure} [H]
			\centering
			\includegraphics[page=1,trim={0cm 0cm 0cm 0cm},clip=true,width=70mm]{figurastfm/Chapter4/PDF/ldr_circuit.pdf}
			\caption{Sun sensor based on \acrshort{LDR} circuit} \label{ldrcircuit}
			\vspace{-2cm}
\end{figure}

As seen in \autoref{ldrcircuit}, it is needed another resistor to complete the voltage divider. Given that a certain \acrshort{LDR} will exhibit a fixed resistance range, the \textbf{excursion~output~range} of the voltage divider (i.e., the difference between the minimum and the maximum voltage measured at the point $ADC\_READ$) will depend on the resistor magnitude chosen. Therefore, it is necessary to choose that resistor in order to \textbf{optimize} the excursion output range and get the most resolution possible.

In the \autoref{ldrcircuit} above, \acrshort{LDR} is denoted as $R^*$ and $R$ is another resistance completing the voltage divider. Depending on the light falling upon the \acrshort{LDR}, $R^*$ varies, so does the voltage $V$. That voltage is directly connected to the \acrshort{ADC} so it is converted into a discrete value. With this circuit, $V$ will increase when light level is high (facing Sun) whereas it will decrease when light reduces (eclipse). 

Ideally, $U$ voltage should go through the range $[0 \text{ V},U_0]$ to maximize distinguishable values. Considering that the only degree of freedom is given by $R$ value, it must be the point of the optimization. On the other hand, extreme resistance values of the \acrshort{LDR} must be known so the whole resistance of the voltage divider can be calculated. Therefore, for a certain resistance $R^* \in [R^*_{\text{min}}, R^*_{\text{max}}]$:

\begin{align*}
U(R)=U_0 \frac{R}{R^* + R}
\end{align*}

Since the aim is maximizing the voltage margin for $U$, it will be denoted as $\Delta U(R)$, which is the electric potential differential across $R^*$:

% <![CDATA[
\begin{align*}
\Delta U(R) &= U_{\text{max}}(R) - U_{\text{min}}(R) \\
&= U_0 \Bigg[\frac{R}{R^*_{\text{min}} + R} - \frac{R}{R^*_{\text{max}} + R} \Bigg] \\
&= U_0 \Bigg[\frac{R(R^*_{\text{max}} - R^*_{\text{min}})}{(R^*_{\text{min}} + R)(R^*_{\text{max}} + R)} \Bigg] \\
\end{align*}


In order to maximize the function $\Delta U(R)$, the derivative $\frac{\partial}{\partial R} \Delta U $ is calculated:

\begin{align*}
\centering
\frac{\partial}{\partial R} \Delta U &= U_0\Bigg[\frac{(R^*_{\text{min}} - R^*_{\text{max}})\big[R^2-R^*_{\text{min}}R^*_{\text{max}}\big]}{(R^*_{\text{min}} + R)^2(R^*_{\text{max}} + R)^2}\Bigg]
\end{align*} %]]>

Then, it is solved $\frac{\partial}{\partial R} \Delta U = 0$ to find extreme point or \textbf{extrema}:

\begin{align*}
\centering
U_0\Bigg[\frac{(R^*_{\text{min}} - R^*_{\text{max}})\big[R^2-R^*_{\text{min}}R^*_{\text{max}}\big]}{(R^*_{\text{min}} + R)^2(R^*_{\text{max}} + R)^2}\Bigg] =  0 \rightarrow 
(R^*_{\text{min}} - R^*_{\text{max}})\big[R^2-R^*_{\text{min}}R^*_{\text{max}}\big] =  0
\end{align*} %]]>

Assuming $R^*_{\text{max}}>R^*_{\text{min}}$, solving the last equation yields to:

\begin{align*}
R &= \sqrt{R^*_{\text{min}}R^*_{\text{max}}} \\
\end{align*} %]]>


Therefore, the optimum resistance in order to get the greatest excursion in the output voltage is given by the square root of the product of the  \acrshort{LDR} extreme values, which means that the optimum resistor depends on the chosen \acrshort{LDR}. In order to make that decision, \autoref{ldrcomp} shows the behavior of the function $\Delta U(R)$ depending on the extreme resistance values of different photoresistor. Because of availability reasons, the \acrshort{LDR} corresponding to the \textbf{blue} series is chosen which, with an associated optimum resistance $ R = $ \SI{1}{k\ohm}, reaches around 80 \% of the possible excursion output.
		
		\begin{figure} [H]
			\centering
			\includegraphics[page=1,trim={0cm 0cm 0cm 0cm},clip=true,width=140mm]{FigurasTFM/Chapter4/PDF/ldr_comp.pdf}
			\caption{$\Delta{U}$ ratio for different LDR extreme values}      		\label{ldrcomp}
			\vspace{-1cm}
\end{figure}

Once it has been accomplished an optimum excursion output, that value has to be read using an \acrshort{ADC}. As for this design, the one chosen is the \textbf{Texas~Instruments~ADS1015}~\cite{ads15}. It is a low-power \acrshort{I2C} compatible \acrshort{ADC} with a maximum resolution of 12 bits and up to four measurements; therefore, considering that there is one Sun Sensor per face of the \glsname{cubesat}, two \acrshort{ADC} are needed. They must be assigned two different \acrshort{I2C} addresses in order to allow communication with both.

Following the same procedure, a simple \acrshort{I2C} communication library has been developed, allowing data polling. This \acrshort{ADC} also features the \textit{ALERT} function, which uses a digital comparator that can issue an alert on the corresponding pin when conversion data exceeds the limits set in  threshold registers. 

\paragraph{Real-Time Clock}

A Real-Time Clock or \acrshort{RTC} is an integrated circuit which keeps track of the time. They must be distinguished from the typical hardware clocks, which are simple signals not related to time in typical units.

\acrshort{RTC} are based on crystal oscillators, typically at a frequency of \textbf{32.768 kHz}, because of its convenience with binary counter circuits. In this case, it is used in conjunction with the \acrshort{CPU}, so it is able to be aware of the current time; this is not needed when it is connected to the Internet, as it gets time by that via, but in any other case the system will start in \textbf{Unix Time}, the first of January, 1970. 

As for this design, the \acrshort{RTC} chosen is the PCF8523 by NXP \cite{rtcnxp}, along with a CR2032 battery which ensures it keeps tracking time when the central system is unpowered. \autoref{rtcpic} shows the placement of the ADS1015 with the crystal oscillator in the \acrshort{PCB}.

		\begin{figure} [H]
			\centering
			\includegraphics[page=1,trim={0cm 0cm 0cm 0cm},clip=true,width=100mm]{FigurasTFM/Chapter4/PDF/rtc.pdf}
			\caption{\acrshort{RTC} and crystal oscillator in the \acrshort{OBC} \acrshort{PCB}}      		\label{rtcpic}
			\vspace{-2cm}
\end{figure}


\newpage

\subsubsection{Flight Software. On-Board Data Handling}

As stated in section \ref{flight}, Flight Software is in charge of controlling the whole system and dealing with all the possible situations during a mission. In this section, it is addressed the software under the whole \glsname{cubesat} as well as implemented a basic \acrshort{OBDH}.


\vspace{-0.5cm}
\paragraph{Operating System} \label{rtosdes}

\acrshort{OS} were divided into Real-time and Non-real-time. The solution taken in this project is somewhere in the middle of both: it is used an initially Non-\acrshort{RTOS} such as Linux, which is patched to be turned into a \acrshort{RTOS}. This way, it is possible to maintain the great drivers compatibility of Linux but adding the deterministic nature of the \acrshort{RTOS}.

In order to perform this change, it is necessary to \textbf{recompile the Linux kernel}. By default, Linux kernel is non-preemptive. To turn it into a preemptive scheduling, it is applied the \textbf{PREEMPT\_RT patch} \cite{linuxrtos}, which provides with faster response times and removes unbounded latencies. 

The first task is downloading the desired version of the kernel. Then, the corresponding version of the patch is also downloaded and unpacked. It can be patched following the instructions in \cite{linuxpat}. After patching, the kernel is rebuilt and configured as \textbf{preemptive}.


\paragraph{Telemetry System}

As for this first prototype, it is designed a simple telemetry sending system using \textbf{Python}, particularly the \textbf{socket} library. It connects to a certain IP address and port with the instructions shown in \autoref{socket}. 


\begin{lstlisting}[language=Python, caption=Connection to server using the \textbf{socket} library, label=socket]
  
	s = socket.socket(socket.AF_INET, socket.SOCK_STREAM)
	host = '192.168.1.118'
	port = 8445
	s.connect((host, port)) 	
	
\end{lstlisting}

Particularly, the tuple IP:Port must match the existing at the other end of the communication, in which \textbf{the COSMOS \glsname{ground} is listening}, as seen in \autoref{cosmosgs}. Once the connection is established, the flow of communication with telemetry and telecommands can start. %As for the downlink (to the \glsname{ground}), the information is packed and sent down using the established connection.

\subsubsection{PCB design}

This section includes the different schematics and information related to the \acrshort{PCB} designed for the \acrshort{OBC}. In order to allow its production, this design must comply with the constraints given by the manufacturer; these depends on the number of layers of the \acrshort{PCB}, while the \acrshort{OBC} is a 4-layers board, the \acrshort{ADCS} \acrshort{PCB} (addressed in section \ref{adcsdes}) is a 2-layers board and the constraints \textbf{change}. Particularly, \autoref{conschino} depicts the different design requirements expected by the manufacturer: the first one for \textbf{2-layers} and the second one for\textbf{ 4-layers}.

			\begin{figure}[H]
			\centering
			\subfloat[2-layers]{\includegraphics[page=1,trim={0cm 0cm 0cm 0cm},clip=true,width=165mm]{figurastfm/Chapter4/PDF/pcbcons.pdf}}
			\\
			\subfloat[4-layers]{\includegraphics[page=1,trim={0cm 0cm 0cm 0cm},clip=true,width=165mm]{figurastfm/Chapter4/PDF/pcbcons4layer.pdf}}
			\caption{\acrshort{PCB} Manufacturing constraints \cite{pablotesis}} \label{conschino}
			\vspace{-1cm}
\end{figure}

\newpage

\autoref{final} shows a 3D render of the top and bottom layer in the final design, while \autoref{persp} shows a perspective view from different angles. 			
%\vspace{1cm}





			\begin{figure}[H]
			\centering
			\subfloat[Top]{\includegraphics[page=1,trim={0cm 0cm 0cm 0cm},clip=true,width=80mm]{figurastfm/Chapter4/PDF/obc_3d.pdf}}
			\quad
			\subfloat[Bottom]{\includegraphics[page=2,trim={0cm 0cm 0cm 0cm},clip=true,width=80mm]{figurastfm/Chapter4/PDF/obc_3d.pdf}}
			\caption{\acrshort{OBC} \acrshort{PCB} final design} \label{final}
			%\vspace{-0.5cm}
\end{figure}

\vspace{1cm}


		\begin{figure}[H]
			\centering
			\subfloat[]{\includegraphics[page=1,trim={0cm 0cm 0cm 0cm},clip=true,width=80mm]{figurastfm/Chapter4/PDF/3d_persp.pdf}}
			\quad
			\subfloat[]{\includegraphics[page=1,trim={0cm 0cm 0cm 0cm},clip=true,width=80mm]{figurastfm/Chapter4/PDF/3d_persp_3.pdf}}
			\caption{\acrshort{OBC} \acrshort{PCB} 3D perspective} \label{persp}
			\vspace{-2cm}
\end{figure}
%
%\includepdf[pages=-,link=true,landscape,linkname=schematics]{FigurasTFM/Chapter4/PDF/sch_OBC.PDF}
%\includepdf[pages=-,link=true,landscape,linkname=schematics]{FigurasTFM/Chapter4/PDF/OBC_PCB.PDF}
\includepdf[pages=-,link=true,landscape,linkname=schematics]{FigurasTFM/Chapter4/PDF/obc_3dview.PDF}


\newpage

\subsubsection{Manufacturing and soldering} \label{obcmanuf}

Once the \acrshort{PCB} are designed, they are sent to a professional manufacturer. The result is shown in \autoref{obcpcb}.

\begin{figure}[H]
			\centering
			\subfloat[Top]{\includegraphics[page=1,trim={0cm 0cm 0cm 0cm},clip=true,width=80mm]{figurastfm/Chapter4/Imagenes/obc_top.png}}
			\quad
			\subfloat[Bottom]{\includegraphics[page=1,trim={0cm 0cm 0cm 0cm},clip=true,width=80mm]{figurastfm/Chapter4/Imagenes/obc_bot.png}}
			\caption{\acrshort{OBC} \acrshort{PCB} manufactured} \label{obcpcb}
			%\vspace{-2cm}
\end{figure}

Because of certain components packages (for instance, the LAN9514 or the Compute Module), this \acrshort{PCB} cannot be soldered by hand. It is needed a technique called \textbf{oven reflow}. It is a procedure in which the first step consists in covering the \acrshort{PCB} with \textbf{solder paste} using a \glsname{stencil} and then the components are placed onto the board. \autoref{obcsolder} shows different moments of this process.

Once the \acrshort{PCB} is fully assembled and all the components are in place, it undergoes the reflow process, which in this case is performed with a domestic oven which had been previously characterized with a trial-and-error approach. Video \autoref{reflow} shows the \acrshort{PCB} undergoing the process (video available with Adobe Acrobat Reader).


\begin{figure}[H]
			\centering
			\subfloat[Solder paste transfer using a \glsname{stencil}]{\includegraphics[page=1,trim={0cm 0cm 0cm 0cm},clip=true,width=81.5mm]{figurastfm/Chapter4/Imagenes/stencil.jpg}} 
			\quad
			\subfloat[QFN package placement]{\includegraphics[page=1,trim={0cm 0cm 0cm 0cm},clip=true,width=76mm]{figurastfm/Chapter4/Imagenes/detail_place.jpg}}
			%\\
			%\subfloat[\acrshort{PCB} partially assembled]{\includegraphics[page=1,trim={0cm 0cm 0cm 0cm},clip=true,width=75mm]{figurastfm/Chapter4/Imagenes/detail_place2.jpg}}
			\caption{Soldering process} \label{obcsolder}
			\vspace{2cm}
\end{figure}

\begin{videoFloat}[H]
	\centering
\includemovie[text={\includegraphics[page=1,trim={0cm 0cm 0cm 0cm},clip=true,width=90mm]{figurastfm/Chapter4/Imagenes/inoven.jpg}}]{10cm}{6cm}{FigurasTFM/Chapter4/Video/makerspace.wmv}
\vspace{2.8cm}
\caption{Oven reflow (video available with Adobe Acrobat Reader)} \label{reflow}
%\vspace{-2cm}
\end{videoFloat}

\newpage

\autoref{topfinalobc} and \autoref{botfinalobc} show the final product after completing soldering.


		\begin{figure} [H]
			\raggedright
			\begin{minipage}{10cm}
 			\includegraphics[page=1,trim={0cm 0cm 0cm 0cm},clip=true,width=100mm]{figurastfm/Chapter4/Imagenes/final_obc.png}
			\centering
			\caption{Top layer of the \acrshort{OBC}}     \label{topfinalobc}
			  \end{minipage}
			\vspace{-1cm}
\end{figure}


		\begin{figure} [H]
		\raggedleft
			\begin{minipage}{10cm}
			\includegraphics[page=1,trim={0cm 0cm 0cm 0cm},clip=true,width=100mm]{figurastfm/Chapter4/Imagenes/final_obc_bot.png}
			\caption{Bottom layer of the \acrshort{OBC}}     \label{botfinalobc}
						  \end{minipage}
			\vspace{-2cm}

\end{figure}


\newpage
\subsection{Attitude Determination and Control System (ADCS)}

This section addresses the design of the \acrshort{ADCS}, in charge of keeping the \glsname{cubesat} correctly oriented and managing external torques properly. At this point, some clarifications must be made with respect to the \acrshort{ADCS} analysis in \autoref{chap:chapter3}. For instance, this subsection will include the \textbf{Co-processing Programmable Core} initially analyzed in the \acrshort{OBC} section; this is due to the need of an independent processing unit which frees the \acrshort{CPU} of this tasks and allows having a real \acrshort{ADCS} managing the issues regarding this matter. On the  other hand, the \textbf{tachometer} is the only sensor which is included in this section, as it is placed in the \acrshort{ADCS} board and is explicitly necessary for its operation. The \textbf{actuators} included in the design are also considered here.

\subsubsection{Co-processing Programmable Core}

As described in section \ref{copropo}, this architectures are usually implemented using \acrshort{FPGA}. Using an \acrshort{FPGA} provides with \textbf{time-critical} capabilities which goes in the same line of the \acrshort{RTOS} introduced in \autoref{rtosdes}: while with the operating system that capability is provided inherently by software, an \acrshort{FPGA} allows expanding that concept to the hardware parcel. Time accuracy and management is a vital issue in space field, which justifies the different designs addressed in this project on said line. 

Therefore, this Master's Thesis will make use of an \acrshort{FPGA}, particularly the ICE40UP5K by Lattice \cite{up5k}. It counts with 5000 \acrshort{LUT} and is provided in a variety of packages, in this case, \textbf{QFN} is chosen. Typically,  \acrshort{FPGA} are programmed from an external memory in a \textbf{master-slave} configuration in which the \acrshort{FPGA} acts as master. In this project, however, given the presence of a central processing unit (the Raspberry Pi 3 Compute Module), it is designed the other way around: it acts as master while the \acrshort{FPGA} does as \textbf{slave}. This way, each time the \acrshort{SBC} boots, it sends (as master) a \textbf{bitstream} through an \acrshort{SPI} channel which contains the programming code. Additionally, the same channel is used to transfer the information during the \acrshort{FPGA} operation, detailed next.

\paragraph{PWM generation}

In this project, the \acrshort{FPGA} main functionality will be generating different \acrshort{PWM} signals to control the actuators, given that the \acrshort{CPU} cannot produce the amount needed; however, it could be used for a variety of purposes, which is the idea under the Co-processing Programmable Core. As described before, the same \acrshort{SPI} channel is used to configure each of the \acrshort{PWM}. \autoref{fpgadiag} depicts the basis of this operation.

		\begin{figure} [H]
			\centering
			\includegraphics[page=1,trim={0cm 0cm 0cm 0cm},clip=true,width=165mm]{FigurasTFM/Chapter4/PDF/fpga.pdf}
			\caption{\acrshort{FPGA} functioning as \acrshort{PWM} generator }      		\label{fpgadiag}
			%\vspace{-2cm}
\end{figure}

The \textbf{bitstream} is composed of two bytes; while the first of them chooses the particular \acrshort{PWM} to be set up, the second one states the configuration of it, in terms of \textbf{frequency} and \textbf{duty cycle}. The \acrshort{PWM} are generated by using a 25 MHz  oscillator, which is divided as needed to accomplish the desired frequency. Finally, the generated \acrshort{PWM} outputs in \textbf{one of the pins} of the \acrshort{GPIO} from the banks. As \textbf{Hardware Description Language} it has been used \textbf{Verilog}.


\subsubsection{Actuator: reaction wheel}

Actuators are used to change or control the attitude of the \glsname{cubesat}. As for this prototype, it has been implemented a \textbf{reaction wheel}. However, the \acrshort{ADCS} is capable of managing up to three \glsname{magnetorquers} which can be added in the future, using also the different \acrshort{PWM} signals available. 

As reaction wheel, it has been reused the one available at \glsname{GranaSAT} laboratory, shown in \autoref{3dreact}. In order to propel the wheel, it is necessary a DC motor. As usual, it is not possible to supply a DC motor directly; in this case, for example, it can provide with just a few mA, so an \textbf{H-bridge} is needed. It is a device composed of four switches that control the flow of large currents to a load, which is precisely its usefulness: allows, with a low-power control signal (in this case, the \acrshort{PWM}) to drive high loads, such as a motor. It is composed of a simple circuit, depicted in \autoref{hbridge}.

		\begin{figure} [H]
			\centering
			\includegraphics[page=1,trim={0cm 0cm 0cm 0cm},clip=true,width=65mm]{FigurasTFM/Chapter4/PDF/H_bridge.pdf}
			\caption{H-bridge basic circuit \cite{wiki}}      		\label{hbridge}
			%\vspace{-2cm}
\end{figure}

In sum, the \acrshort{PWM} generated by the \acrshort{FPGA} attack the H-bridge (control signals) and it supplies the load with the needed current (larger than the one the \acrshort{FPGA} can provide). The H-bridge could attack the motor directly and supply it with all the current demanded, however, this implies two drawbacks: on the one hand, the motor would be driven at the maximum speed continuously, which reduces its control dramatically; on the other hand, it would be propelled in the same fixed direction. An H-bridge allows controlling both parameters.

As for this design, several H-bridges are used, listed next:

\begin{itemize} [noitemsep,topsep=0pt]
\item\textbf{ ST L298N}
\end{itemize}

The L298N by ST \cite{l298n} is an integrated high voltage, high current Dual Full-Bridge Driver. It allows an operating supply voltage of up to 46 V @ 4 A. Additionally, it features two \textit{enable} inputs that allow the L298 to be enabled or disabled independently of the input signals. Because of its high power capabilities, it is used to control the motor driving the reaction wheel. It will be used a single unit, which provides with \textbf{four outputs}; as the motor only needs a couple, the other two are used for the X-axis magnetorquer.


		\begin{figure} [H]
			\centering
			\includegraphics[page=1,trim={0cm 0cm 0cm 0cm},clip=true,width=50mm]{FigurasTFM/Chapter4/Imagenes/l298.jpg}
			\caption{ST L298N \cite{l298n}}  %    		\label{hbridge}
			\vspace{-1cm}
\end{figure}
\newpage



\begin{itemize} [noitemsep,topsep=0pt]
\item \textbf{ROHM BD6211F-E2}
\end{itemize}

As stated before, the \acrshort{ADCS} will count with the capability to use \glsname{magnetorquers} in the future. They must be supplied with an H-bridge, because of the same reasons of a motor, essentially. However, the current needed for \glsname{magnetorquers} is much lower, so a smaller and cheaper bridge is used, particularly the BD6211F-E2 by ROHM \cite{rohm}. It is able to supply with up to 1 A with 5.5 V. 

		\begin{figure} [H]
			\centering
			\includegraphics[page=1,trim={0cm 0cm 0cm 0cm},clip=true,width=45mm]{FigurasTFM/Chapter4/Imagenes/ROHM.jpg}
			\caption{ROHM BD6211F-E2 \cite{rohm}}  %    		\label{hbridge}
			%\vspace{-2cm}
\end{figure}

Each BD6211F-E2 exhibits two outputs; considering that the X-axis magnetorquer is supplied by the L928N, each of the axis left needs another couple of outputs so they are necessary two BD6211F-E2. 



\subsubsection{Sensor: Tachometer}

In order to adequately control the reaction wheel, it is implemented a \textbf{tachometer}, able to measure the rotation speed of the wheel. To design the tachometer, it is used a \textbf{Reflective Optical Sensor}. These devices produce an infrared light onto a certain surface and the reflection is measured; when there is no reflection the sensor will show a \textbf{0} at the output, while a certain reflection will produce a \textbf{1}. Particularly, the sensor used is the CNY70 by Vishay \cite{cny}, shown in Figure\autoref{cny70}.


\begin{figure}[H]
			\centering
			\subfloat[Package \label{cny70}]{\includegraphics[page=1,trim={0cm 0cm 0cm 0cm},clip=true,width=35mm]{figurastfm/Chapter4/Imagenes/cny70.jpg}} 
			\quad
			\subfloat[Top view]{\includegraphics[page=1,trim={0cm 0cm 0cm 0cm},clip=true,width=45mm]{figurastfm/Chapter4/PDF/cny70.pdf}}
			\caption{Vishay CNY70 \cite{cny}} %\label{obcsolder}
			%\vspace{-2cm}
\end{figure}

The maximum distance at which the CNY70 is able to measure reflection depends on the \textbf{forward current} of the transistor, which in turn is subject to the emitter diode current. Therefore, the resistors needed to bias the device have to be optimized. The maximum measurement distance is accomplished when the forward current is about \textbf{20 mA}, so regarding the \textbf{diode}, the following equations can be considered:


\begin{equation}
\centering
I_{\text{forward}}=\frac{V_{\text{cc}}-V_{\text{forward}}}{R_{\text{diode}}}
\end{equation} %]]>


Considering a supply of 3.3 V and taking into account that, according to the datasheet \cite{cny}, the typical forward voltage is about\textbf{ 1.25 V}, it yields to:


\begin{equation} 
\centering
20 \text{ mA}=\frac{5 \text{ V}- 1.25 \text{ V}}{R_{\text{D}}} \rightarrow R_{\text{D}} = \SI{102.5}{\ohm}
\end{equation} %]]>

As for the phototransistor, for a typical collector current of 0.1 mA, at a distance of 0.3 mm, the maximum $V_{\text{CE}}$ under saturation is 0.3 V. So, substituting in \autoref{eqtrans}:



\begin{align} \label{eqtrans}
\centering
I_{\text{C}}=\frac{V_{\text{cc}}-V_{\text{CE}}}{R_{\text{collector}}} \rightarrow R_{\text{collector}} = \SI{30}{k\ohm}
\end{align} %]]>

In order to test the accuracy of the system, the CNY70 is used with the configuration determined to measure the rotation rate of a reaction wheel. At the same time, it is measured using a professional tachometer. The results are plot together at \autoref{tachom}.

		\begin{figure} [H]
			\centering
			\includegraphics[page=1,trim={0cm 0cm 0cm 0cm},clip=true,width=137mm]{FigurasTFM/Chapter4/PDF/rpm.pdf}
			\caption{Rotation rate measured with a professional tachometer and with the CNY70 circuit designed} 	\label{tachom}
			\vspace{-2cm}
\end{figure}

As expected, the rotation rate is linear with the voltage supply. On the other hand, the error made is completely negligible, with an almost exact coincidence between both measurements. This validates the design of the tachometer based on the CNY70.


\subsubsection{Control law}

The control law is the software implementation governing the hardware above. It deals with the variety of situations along a mission and acts consequently. As starting point for this prototype, it is implemented a control law based on a \acrshort{PID}, described in section \ref{pid3}. Particularly, it has been used the library \textbf{simple\_pid} in Python~\cite{simplepid}. \autoref{pidop} shows an example of the \acrshort{PID} counteracting the movement of the  \glsname{cubesat}.

\begin{figure}[H]
			\centering
			\subfloat[Long operation]{\includegraphics[page=1,trim={0cm 0cm 0cm 0cm},clip=true,width=147mm]{figurastfm/Chapter4/PDF/pid3.pdf}} 
			\quad
	\subfloat[Detail]{\includegraphics[page=1,trim={0cm 0cm 0cm 0cm},clip=true,width=147mm]{figurastfm/Chapter4/PDF/pid2.pdf}} 
			\caption{\acrshort{PID} counteracting Z-axis rotation} \label{pidop}
			\vspace{-2cm}
\end{figure}


\newpage
\subsubsection{PCB design} \label{adcsdes}
\autoref{finaladcs} shows a 3D render of the top and bottom layer in the final design, while \autoref{perspadcs} shows a perspective view from different angles. 			
%\vspace{1cm}





			\begin{figure}[H]
			\centering
			\subfloat[Top]{\includegraphics[page=1,trim={0cm 0cm 0cm 0cm},clip=true,width=80mm]{figurastfm/Chapter4/PDF/adcs_top.pdf}}
			\quad
			\subfloat[Bottom]{\includegraphics[page=1,trim={0cm 0cm 0cm 0cm},clip=true,width=80mm]{figurastfm/Chapter4/PDF/adcs_bot.pdf}}
			\caption{\acrshort{ADCS} \acrshort{PCB} final design} \label{finaladcs}
			%\vspace{-0.5cm}
\end{figure}

\vspace{1cm}


		\begin{figure}[H]
			\centering
			\subfloat[]{\includegraphics[page=1,trim={0cm 0cm 0cm 0cm},clip=true,width=80mm]{figurastfm/Chapter4/PDF/adcs_persp.pdf}}
			\quad
			\subfloat[]{\includegraphics[page=1,trim={0cm 0cm 0cm 0cm},clip=true,width=80mm]{figurastfm/Chapter4/PDF/adcs_persp2.pdf}}
			\caption{\acrshort{ADCS} \acrshort{PCB} 3D perspective} \label{perspadcs}
			\vspace{-2cm}
\end{figure}

%\includepdf[pages=-,link=true,landscape,linkname=schematics]{FigurasTFM/Chapter4/PDF/ADCS.PDF}
%\includepdf[pages=-,link=true,landscape,linkname=schematics]{FigurasTFM/Chapter4/PDF/ADCS_PCB.PDF}
\includepdf[pages=-,link=true,landscape,linkname=schematics]{FigurasTFM/Chapter4/PDF/adcs_3d.PDF}

\subsubsection{Manufacturing and soldering}

In this case, the \acrshort{PCB} is composed of 2 layers. The design is sent to a professional manufacturer too, resulting in the \acrshort{PCB} shown in \autoref{adcspcb}.

\begin{figure}[H]
			\centering
			\subfloat[Top]{\includegraphics[page=1,trim={0cm 0cm 0cm 0cm},clip=true,width=80mm]{figurastfm/Chapter4/Imagenes/adcs_top.png}}
			\quad
			\subfloat[Bottom]{\includegraphics[page=1,trim={0cm 0cm 0cm 0cm},clip=true,width=80mm]{figurastfm/Chapter4/Imagenes/adcs_bot.png}}
			\caption{\acrshort{ADCS} \acrshort{PCB} manufactured} \label{adcspcb}
			\vspace{-0.5cm}
\end{figure}

As for the soldering process, the procedure is the same as the one described in section \ref{obcmanuf}; the result is shown in \autoref{adcspcbsolder}.

		\begin{figure} [H]
			\centering
			\includegraphics[page=1,trim={0cm 0cm 0cm 0cm},clip=true,width=115mm]{FigurasTFM/Chapter4/Imagenes/adcs_final.png}
			\caption{\acrshort{ADCS} \acrshort{PCB} after assembling and soldering} 	\label{adcspcbsolder}
			\vspace{-2cm}
\end{figure}



\newpage
\subsection{Electrical Power System (EPS)}

Regarding the \acrshort{EPS}, this Master's Thesis addresses some of the crucial choices to be made when facing the design of this kind of system. While the electronics design is left for future projects, this one will lay the foundations for the photovoltaic subsystem of the \textbf{GranaSAT-I}, by analyzing and characterizing different solar panels in order to determine the optimum one according to the rest of the \glsname{cubesat}. On the other hand, a set of batteries are characterized at different charge/discharge rates, solving part of the decisions regarding the energy storage system. 


\subsubsection{Solar panels}

Solar cells were profusely analyzed in \autoref{chap:chapter3} so in this section, the \textbf{testbench} designed for its characterization is presented first and then different solar panels are characterized.

\paragraph{Measurements testbench}

Characterizing solar panels is a complex matter. As described in previous sections, generally speaking, solar cells behavior is described by its I-V curve. However, obtaining said curve is not as simple as it could seem: if the output of a solar panel is measured when it is under light, a certain voltage is obtained, the $V_{\text{OC}}$, which is useless as it cannot deliver any power at that point (see section \ref{solarcellsexpl}). So in order to get the complete curve it is necessary to sweep a voltage range when the solar panel is \textbf{loaded}, this is, emulating a real operation environment. It is easy to understand under the basis stated in previous explanations: if the solar panel is forced to work at a certain voltage, against a certain load, it will deliver a given power; if that voltage varies, so does the rest, getting the complete I-V curve. Particularly, the voltage sweep must cover the whole expected functioning range. 

Simulating all these variables requires an adequate \textbf{testbench} composed of different equipment, depicted in \autoref{testbench}. The testbench has been designed in a way which allows measuring in both, \textbf{outdoor} (solar light) and \textbf{indoor} (simulators) environments. Firstly, the solar panel is connected to the \textbf{load}, which is composed of two devices: the \textbf{HP 6063B} \cite{hpdata1} and the \textbf{HP E3631A} \cite {hpdata2}. While the first is an electronic load, used to perform the \textbf{voltage sweep}, the second is a simple DC power source used to impose a certain voltage offset to the HP 6063B output, needed for correct functioning. Both, acting as load, forces the solar panel output a certain voltage (\textbf{$V_{\text{sweep}}$}); as it is exposed to an irradiance source (whether actual Sun light or simulated), the solar panel provides with the \textbf{current} corresponding to the impedance seen from the panel, which varies with the voltage sweep. As the sweep is performed, both the voltage imposed and the current provided is measured using the multimeter SIGLENT~SDM3065X \cite{siglentmult}. Additionally, the testbench counts with a pyranometer to be placed along with the solar panel, so the irradiance present at the test can be taken into account; it is also connected to the multimeter, and the voltage measured is converted into an irradiance measurement.

The whole equipment is controlled using a \acrshort{matlab} library developed progressively throughout the years at the \glsname{GranaSAT} laboratory, which also allows downloading the data. Finally, the I-V curve can be plotted.

		\begin{figure} [H]
			\centering
			\includegraphics[page=1,trim={0cm 0cm 0cm 0cm},clip=true,width=165mm]{FigurasTFM/Chapter4/PDF/testbench.pdf}
			\caption{Solar panels characterization testbench} 	\label{testbench}
			%\vspace{-0.5cm}
\end{figure}


\paragraph{Characterizations}

In this section, different solar panels are characterized using the testbench just described. Firstly,  in order to assure the correct functioning of the characterization bench, it is performed a test under Sun's illumination, on a high quality solar panel with a known IV curve. It is a 5 W polycristalline solar panel with an open-circuit voltage of 22 V \cite{panelrs}, shown in \autoref{rs}.

The irradiance measured outdoors is \SI{1000}{W/m^2} and it slightly varied for all the characterizations with solar lightning performed.


		\begin{figure} [H]
			\centering
			\includegraphics[page=1,trim={0cm 0cm 0cm 0cm},clip=true,width=70mm]{FigurasTFM/Chapter4/PDF/panel_rs.pdf}
			\caption{5 W solar panel \cite{panelrs}} 	\label{rs}
			\vspace{-0.5cm}
\end{figure}

\autoref{rs1} shows the IV curve of the panel, along with the power delivered by in those conditions.


\begin{figure}[H]
			\centering
			\includegraphics[page=1,trim={0cm 0cm 0cm 0cm},clip=true,width=145mm]{figurastfm/Chapter4/PDF/plot_panel_rs.pdf}
			\caption{Characterization of the RS solar panel} \label{rs1}
			\vspace{-0.5cm}
\end{figure}

The panel behaves exceptionally well, with an output even \textbf{above 5 W} and an \textbf{MPP~=~[17~V,~0.31~A]}. This \textbf{validates the functioning of the testbench}, so in the next pages different small-sized solar panels are also characterized.

The solar panels candidates to be used in the \glsname{cubesat} will be characterized facing both Sun's and \acrshort{LED} lightning. \autoref{montaje} shows the assembly for the measurement using the \acrshort{LED} simulator.

\begin{figure}[H]
			\centering
			\subfloat[]{\includegraphics[page=1,trim={0cm 0cm 0cm 0cm},clip=true,width=80mm]{figurastfm/Chapter4/Imagenes/ledpanel1.jpg}}
			\quad
			\subfloat[]{\includegraphics[page=1,trim={0cm 0cm 0cm 0cm},clip=true,width=80mm]{figurastfm/Chapter4/Imagenes/ledpanel2.jpg}}
			\caption{Assembly to characterize solar panels with \acrshort{LED} lightning} \label{montaje}
			%\vspace{-0.5cm}
\end{figure}

The results are divided into the three solar panels characterized, each one with two plots depending on the illumination. For each panel, it is indicated the theoretical performance according to the manufacturer. \autoref{fotospanel} shows the panels under test.
\begin{landscape}


\vspace*{\fill}
\begin{figure}[H]
			\centering
			\subfloat[1: 5 V @ 1.3 W]{\includegraphics[page=1,trim={0cm 0cm 0cm 0cm},clip=true,width=65mm]{figurastfm/Chapter4/Imagenes/panel1.png}}
			\quad
			\subfloat[2: 12 V @ 1.5 W]{\includegraphics[page=1,trim={0cm 0cm 0cm 0cm},clip=true,width=65mm]{figurastfm/Chapter4/Imagenes/panel2.png}}
			\quad
				\subfloat[3: 5.5 V @ 1 W]{\includegraphics[page=1,trim={0cm 0cm 0cm 0cm},clip=true,width=65mm]{figurastfm/Chapter4/Imagenes/panel3.png}}			
			\caption{Solar panels characterized} \label{fotospanel}
			%\vspace{-0.5cm}
\end{figure}
\vspace*{\fill}


\end{landscape}


\begin{itemize} [noitemsep,topsep=0pt]
\item \textbf{Solar panel 1: 5 V @ 1.3 W}
\end{itemize}


\begin{figure}[H]
			\centering
			\subfloat[Solar lightning - MPP = (11 V, 0.097 A)]{\includegraphics[page=1,trim={0cm 0cm 0cm 0cm},clip=true,width=135mm]{figurastfm/Chapter4/PDF/panel2_sol.pdf}}
			\\
			\subfloat[\acrshort{LED} lightning - MPP = (11.5 V, 0.0197 A)]{\includegraphics[page=1,trim={0cm 0cm 0cm 0cm},clip=true,width=135mm]{figurastfm/Chapter4/PDF/panel2_led.pdf}}
			\caption{Solar panel 1. Characterization results} %\label{adcspcb}
			\vspace{-2cm}
\end{figure}

\newpage


\begin{itemize} [noitemsep,topsep=0pt]
\item \textbf{Solar panel 2: 12 V @ 1.5 W}
\end{itemize}

\begin{figure}[H]
			\centering
			\subfloat[Solar lightning - MPP = (5 V, 0.165 A)]{\includegraphics[page=1,trim={0cm 0cm 0cm 0cm},clip=true,width=135mm]{figurastfm/Chapter4/PDF/panel3_sol.pdf}}
			\\
			\subfloat[\acrshort{LED} lightning - MPP = (5 V, 0.035 A)]{\includegraphics[page=1,trim={0cm 0cm 0cm 0cm},clip=true,width=135mm]{figurastfm/Chapter4/PDF/panel3_led.pdf}}
			\caption{Solar panel 2. Characterization results} %\label{adcspcb}
			\vspace{-2cm}
\end{figure}



\newpage

\begin{itemize} [noitemsep,topsep=0pt]
\item \textbf{Solar panel 3: 5.5 V @ 1 W}
\end{itemize}

\begin{figure}[H]
			\centering
			\subfloat[Solar lightning - MPP = (5 V, 0.1 A)]{\includegraphics[page=1,trim={0cm 0cm 0cm 0cm},clip=true,width=135mm]{figurastfm/Chapter4/PDF/panel1_sol.pdf}}
			\\
			\subfloat[\acrshort{LED} lightning - MPP = (5.3 V, 0.024 A)]{\includegraphics[page=1,trim={0cm 0cm 0cm 0cm},clip=true,width=135mm]{figurastfm/Chapter4/PDF/panel1_led.pdf}}
			\caption{Solar panel 3. Characterization results} %\label{adcspcb}
			\vspace{-2cm}
\end{figure}


\newpage

From the previous plots, several parameters can be obtained; they are arranged in \autoref{tablepanel} along with the figures of merit of each panel. Definitions and formulas for the parameters exposed can be found in section \ref{solarcellsexpl}.


\begin{table}[]
\centering
\begin{tabular}{cccccccc}
\hline
\textbf{Panel}              & \textbf{Lightning} & \textbf{$V_{\text{MP}}$ (V)} & \textbf{$I_{\text{MP}}$ (A)} & \textbf{$P_{\text{MP}}$ (W)} & \multicolumn{1}{l}{\textbf{Fill Factor}} & \textbf{Incident power (W)} & \textbf{$\eta$} \\ \hline
\multicolumn{1}{c|}{1} &                    &                              &                              &                              & \multicolumn{1}{l}{}                     &                             &                 \\
\multicolumn{1}{c|}{}  & Sun                & 11                           & 0.097                        & 1.067                        & 0.70                                     & 8.8                         & 12 \%           \\
\multicolumn{1}{c|}{}  & \acrshort{LED}     & 11.5                         & 0.0197                       & 0.23                         & 0.82                                     & 2.07                        & 11 \%           \\ \hline
\multicolumn{1}{c|}{2} &                    &                              &                              &                              &                                          &                             &                 \\
\multicolumn{1}{c|}{}  & Sun                & 5                            & 0.165                        & 0.825                        & 0.75                                     & 9.78                        & 8 \%            \\
\multicolumn{1}{c|}{}  & \acrshort{LED}     & 5                            & 0.035                        & 0.175                        & 0.64                                     & 2.30                        & 8 \%            \\ \hline
\multicolumn{1}{c|}{3} &                    &                              &                              &                              &                                          &                             &                 \\
\multicolumn{1}{c|}{}  & Sun                & 5                            & 0.1                          & 0.5                          & 0.81                                     & 9.03                        & 6 \%            \\
\multicolumn{1}{c|}{}  & \acrshort{LED}     & 5.3                          & 0.024                        & 0.13                         & 0.80                                     & 2.12                        & 6 \%            \\ \hline
\end{tabular}
\caption{Results of the solar panels characterized}
\label{tablepanel}
\end{table}

Additionally, several conclusions can be extracted from the previous characterization, summarized next.



\begin{itemize} [topsep=0pt]
	
	\item In order to get realistic results, the incident power has been calculated for each panel, i.e., the power received by the solar panel is not the one irradiated by the source (\SI{1000}{W/m^2} for the Sun or \SI{235}{W/m^2} for the \acrshort{LED}), but the \textbf{proportional part} considering the panel's size. Indeed, the \textbf{packing density} also has to be taken into account, with an estimated value of \textbf{0.8}.
	\item Although, at first sight, the panels seem to exhibit a significantly worse performance when the \acrshort{LED} simulator is used, actually \textbf{the efficiency is almost the same}. Particularly, panels 2 and 3 features the same \textbf{power conversion efficiency} regardless of the lightning. 
	\item On the other hand, the power delivered is about a fifth when the \acrshort{LED} simulator is lightning, but that is logical, considering that the incident power is also about \textbf{20~\%} of the one provided by the Sun.
	\item The $V_{\text{MP}}$ is also similar regardless of the irradiance source used, which means that, for an equivalent load, \textbf{the spectrum of both power sources is reasonably similar too}. The decrease in power is then a consequence of the decrease in the \textbf{output current}, $I_{\text{MP}}$, directly related to the incident power. 
	\item In the same way, the \textbf{Fill Factor} slightly differs with the kind of lightning, indeed, the panel 3 features almost an identical factor. Furthermore, \textbf{in all cases the factor is pretty acceptable}, with a minimum value of 0.64 when the panel 2 is lighted with the \acrshort{LED} simulator. With a mean value of about 0.75, this implies that about that percentage of the power in shortcircuit is provided as \textbf{real power}. 
	\item Regarding the \textbf{power conversion efficency} or $\eta$ is really poor in all cases; this issue was expectable considering that these are low-cost solar panels. However, while the efficiency exhibited by panels 2 and 3 is even lower than typical values, the efficiency associated to the panel 1 (12 \%) \textbf{can be deemed acceptable.}
	
	\item In sum, this testbench has allowed, on the one hand: characterizing the behavior of different solar panels in terms of \textbf{power and figures of merit}; on the other hand, \textbf{validating} the \acrshort{LED} simulator, as the results are coherent with the ones obtained using solar lightning, considering the \textbf{difference in incident power}.
	
\end{itemize}

The final choice will depend on the overall design of the \acrshort{EPS} as well as on the battery used.

\subsubsection{Batteries characterization} \label{battscharac}

As for the energy storage subsystem, because this Master's Thesis has addressed the power consumption of some elements, such as the \textbf{Central Processing Unit} (see section \ref{rpipower}), some Li-ion batteries are characterized, so it can be used as a reference for the future \acrshort{EPS}. The definitions of the most relevant \textbf{figures of merit} related to batteries, mentioned next, can be found in section \ref{energystr}.

Particularly, several batteries with different \textbf{capacities} (1100 and 1800 mAh) are tested using a wide range of values for the \textbf{dimensional abuse}. This test is performed using the battery charging centre \textbf{VOLCRAFT ALC-8500}, shown in \autoref{volcraft}.


\begin{figure}[H]
			\centering
			\includegraphics[page=1,trim={0cm 0cm 0cm 0cm},clip=true,width=80mm]{figurastfm/Chapter4/Imagenes/alc.jpg}
			\caption{VOLCRAFT ALC-8500} \label{volcraft}
			%\vspace{-0.5cm}
\end{figure}

Once again, the results allow both, characterizing the batteries and making a first approximation to the energy needs of the system and the battery \textbf{run-time} when supplying it.  The results are shown in the following pages.

\includepdf[pages=-,link=true,landscape,linkname=schematics]{FigurasTFM/Chapter4/PDF/completo_1100.PDF}

\includepdf[pages=-,link=true,linkname=tp,landscape,linkname=schematics]{FigurasTFM/Chapter4/PDF/charge_cap_1800.PDF}

Some conclusions can be extracted from the previous results:

\begin{itemize} [topsep=0pt]
	
	\item \textbf{Neither voltage drop nor increase is linear} with discharge or charge respectively. Therefore, neither battery duration can be expected to behave linearly, e.g., with a \textbf{discharge rate} C/10, the run-time of the 1100 mAh battery is about 10 hours, and with C/20, instead of 20 hours, is slightly 19 hours.
	\item As expected, the battery behaves as a \textbf{constant current source}, whose voltage varies. This can be seen clearly in the plot analyzing the current drop when discharging. However, when charging the behavior is different, although current is constant for a while, at the end of the charge, \textbf{it decreases almost exponentially}.
	\item Because current drop is constant, capacity drop is too, shown at the third plot on the left. It is interesting to point out another phenomena: \textbf{the batteries do not completely discharge}. For instance, see the capacity increase in the 1100 mAh battery; it clearly \textbf{surpasses its theoretical capacity}, reaching more than 1200 mAh. However, the capacity drop reaches the 1100 mAh stated by the manufacturer. This is actually a \textbf{protection mechanism} of the charging centre, which prevents the battery voltage from dropping below a certain value (about 3.1 V), as this would permanently \textbf{damage the battery.}
	\item It is worth noting that, when charging the 1800 mAh battery, the capacity slightly surpasses 1500 mAh, but it should reach almost 1900 mAh, according to the behavior seen before. Of course, this supposes that the capacity drop is about 1500 mAh too, instead of 1800 mAh. In this case, the manufacturer has provided a \textbf{fake total capacity.}
	\item Another interesting point is that batteries \textbf{cannot work at every dimensional abuse}; it depends on the specific battery. If a certain battery is forced to work at an inadequate rate, its behavior will be somewhat erratic. An example of this can be seen when the 1800 mAh battery is charged at a C/2 rate. In fact, the current increase with charge is not constant at any point, after reaching the maximum it immediately start falling exponentially. The voltage increase with this rate differs from the others too. Although not strange behavior is noticed when discharged at this rate,\textbf{ it is inadvisable to use the battery with this configuration}, to prevent any damage.
\end{itemize}


In order to get a rough approximation of the battery run-time when supplying the \glsname{cubesat}, it can be applied the \textbf{Pareto principle}: the 80 \% of the whole power consumption is produced by only a 20 \% of the components, in this case, the Raspberry~Pi~3~Compute~Module. Therefore, it is estimated the battery run-time assuming that it supplies only that component, using a discharge \textbf{dimensional abuse} coherent with the results gathered in section \ref{rpipower}. The results are summarized in \autoref{runtime1} and \autoref{runtime2}, for the 1800 mAh battery.

% Please add the following required packages to your document preamble:
% \usepackage[table,xcdraw]{xcolor}
% If you use beamer only pass "xcolor=table" option, i.e. \documentclass[xcolor=table]{beamer}
\begin{table}[H]
\centering
\begin{tabular}{ccccccccccc}
\multicolumn{11}{c}{\textbf{3.3 V supply - 1800 mAh battery run-time estimation - Stationary state (C/10, 5.51 Wh)}}                                                                                                                                                                                                                                                                         \\ \hline
\textbf{Cores}         & \textbf{Mode}                    & \multicolumn{3}{c}{\textbf{Current (A)}}                                                                  & \multicolumn{3}{c}{\textbf{Power (W)}}                                                                    & \multicolumn{3}{c}{\textbf{Time (h)}}                                                                    \\ \hline
\textbf{}              & \textbf{}                        & \textbf{Mean}             & \textbf{RMS}              & \textbf{Peak}                                     & \textbf{Mean}             & \textbf{RMS}              & \textbf{Peak}                                     & \textbf{Mean}             & \textbf{RMS}              & \textbf{Peak}                                    \\ \cline{3-11} 
\multicolumn{1}{c|}{1} & Powersave                        & 0.147                     & 0.147                     & \multicolumn{1}{c|}{{\color[HTML]{FE0000} 0.156}} & 0.484                     & 0.484                     & \multicolumn{1}{c|}{{\color[HTML]{FE0000} 0.514}} & 11.38                     & 11.38                     & {\color[HTML]{FE0000} 10.72}                     \\
\multicolumn{1}{c|}{}  & Ondemand                         & 0.147                     & 0.147                     & \multicolumn{1}{c|}{{\color[HTML]{FE0000} 0.155}} & 0.485                     & 0.485                     & \multicolumn{1}{c|}{{\color[HTML]{FE0000} 0.511}} & 11.36                     & 11.36                     & {\color[HTML]{FE0000} 10.77}                     \\
\multicolumn{1}{c|}{}  & Performance                      & 0.188                     & 0.188                     & \multicolumn{1}{c|}{{\color[HTML]{FE0000} 0.213}} & 0.620                     & 0.620                     & \multicolumn{1}{c|}{{\color[HTML]{FE0000} 0.702}} & 8.88                      & 8.89                      & {\color[HTML]{FE0000} 7.85}                      \\
\multicolumn{1}{l|}{}  & \multicolumn{1}{l}{Conservative} & \multicolumn{1}{l}{0.146} & \multicolumn{1}{l}{0.146} & \multicolumn{1}{l|}{{\color[HTML]{FE0000} 0.155}} & \multicolumn{1}{l}{0.482} & \multicolumn{1}{l}{0.482} & \multicolumn{1}{l|}{{\color[HTML]{FE0000} 0.510}} & \multicolumn{1}{l}{11.43} & \multicolumn{1}{l}{11.43} & \multicolumn{1}{l}{{\color[HTML]{FE0000} 10.80}} \\ \hline
\multicolumn{1}{l}{}   & \multicolumn{1}{l}{}             & \multicolumn{1}{l}{}      & \multicolumn{1}{l}{}      & \multicolumn{1}{l|}{{\color[HTML]{FE0000} }}      & \multicolumn{1}{l}{}      & \multicolumn{1}{l}{}      & \multicolumn{1}{l|}{{\color[HTML]{FE0000} }}      & \multicolumn{1}{l}{}      & \multicolumn{1}{l}{}      & \multicolumn{1}{l}{{\color[HTML]{FE0000} }}      \\
\multicolumn{1}{c|}{4} & Powersave                        & 0.145                     & 0.145                     & \multicolumn{1}{c|}{{\color[HTML]{FE0000} 0.149}} & 0.477                     & 0.477                     & \multicolumn{1}{c|}{{\color[HTML]{FE0000} 0.491}} & 11.54                     & 11.54                     & {\color[HTML]{FE0000} 11.23}                     \\
\multicolumn{1}{c|}{}  & Ondemand                         & 0.148                     & 0.147                     & \multicolumn{1}{c|}{{\color[HTML]{FE0000} 0.168}} & 0.487                     & 0.486                     & \multicolumn{1}{c|}{{\color[HTML]{FE0000} 0.554}} & 11.32                     & 11.34                     & {\color[HTML]{FE0000} 9.94}                      \\
\multicolumn{1}{c|}{}  & Performance                      & 0.186                     & 0.186                     & \multicolumn{1}{c|}{{\color[HTML]{FE0000} 0.197}} & 0.612                     & 0.612                     & \multicolumn{1}{c|}{{\color[HTML]{FE0000} 0.651}} & 9.00                      & 9.00                      & {\color[HTML]{FE0000} 8.46}                      \\
\multicolumn{1}{c|}{}  & Conservative                     & 0.146                     & 0.145                     & \multicolumn{1}{c|}{{\color[HTML]{FE0000} 0.150}} & 0.481                     & 0.479                     & \multicolumn{1}{c|}{{\color[HTML]{FE0000} 0.495}} & 11.46                     & 11.50                     & {\color[HTML]{FE0000} 11.13}                     \\ \hline
\end{tabular}
\caption{1800 mAh battery run-time estimation in stationary state, with a dimensional abuse C/10}
\label{runtime1}
\end{table}
\vspace{-0.8cm}


As seen in the battery analysis performed, the capacity can vary depending on the dimensional abuse used to charge it; in this case, the battery exhibits a stored energy of 5.51 Wh. Recalling the results in section \ref{rpipower}, the \acrshort{CPU} current drain in stationary state can be approximated by C/10 for this battery. With these considerations, the battery-runtime is estimated taking three values: mean, \acrshort{RMS} and Peak. The results are really satisfactory as, for the maximum performance configuration, the battery is estimated to supply enough power for more than 8 hours. \autoref{runtime2} summarizes the results of this test under a high workload.


% Please add the following required packages to your document preamble:
% \usepackage[table,xcdraw]{xcolor}
% If you use beamer only pass "xcolor=table" option, i.e. \documentclass[xcolor=table]{beamer}
\begin{table}[H]
\centering
\begin{tabular}{ccccccccccc}
\multicolumn{11}{c}{\textbf{3.3 V supply - 1800 mAh battery run-time estimation - High load (C/5, 5.25 Wh)}}                                                                                                                                                                                                                                                                              \\ \hline
\textbf{Cores}         & \textbf{Mode}                    & \multicolumn{3}{c}{\textbf{Current (A)}}                                                                  & \multicolumn{3}{c}{\textbf{Power (W)}}                                                                    & \multicolumn{3}{c}{\textbf{Time (h)}}                                                                 \\ \hline
\textbf{}              & \textbf{}                        & \textbf{Mean}             & \textbf{RMS}              & \textbf{Peak}                                     & \textbf{Mean}             & \textbf{RMS}              & \textbf{Peak}                                     & \textbf{Mean}            & \textbf{RMS}             & \textbf{Peak}                                   \\ \cline{3-11} 
\multicolumn{1}{c|}{1} & Powersave                        & 0.294                     & 0.296                     & \multicolumn{1}{c|}{{\color[HTML]{FE0000} 0.370}} & 0.971                     & 0.976                     & \multicolumn{1}{c|}{{\color[HTML]{FE0000} 1.221}} & 5.41                     & 5.38                     & {\color[HTML]{FE0000} 4.30}                     \\
\multicolumn{1}{c|}{}  & Ondemand                         & 0.461                     & 0.459                     & \multicolumn{1}{c|}{{\color[HTML]{FE0000} 0.568}} & 1.521                     & 1.515                     & \multicolumn{1}{c|}{{\color[HTML]{FE0000} 1.874}} & 3.45                     & 3.47                     & {\color[HTML]{FE0000} 2.80}                     \\
\multicolumn{1}{c|}{}  & Performance                      & 0.461                     & 0.465                     & \multicolumn{1}{c|}{{\color[HTML]{FE0000} 0.567}} & 1.523                     & 1.533                     & \multicolumn{1}{c|}{{\color[HTML]{FE0000} 1.872}} & 3.45                     & 3.42                     & {\color[HTML]{FE0000} 2.80}                     \\
\multicolumn{1}{l|}{}  & \multicolumn{1}{l}{Conservative} & \multicolumn{1}{l}{0.456} & \multicolumn{1}{l}{0.459} & \multicolumn{1}{l|}{{\color[HTML]{FE0000} 0.563}} & \multicolumn{1}{l}{1.504} & \multicolumn{1}{l}{1.515} & \multicolumn{1}{l|}{{\color[HTML]{FE0000} 1.858}} & \multicolumn{1}{l}{3.49} & \multicolumn{1}{l}{3.47} & \multicolumn{1}{l}{{\color[HTML]{FE0000} 2.83}} \\ \hline
\multicolumn{1}{l}{}   & \multicolumn{1}{l}{}             & \multicolumn{1}{l}{}      & \multicolumn{1}{l}{}      & \multicolumn{1}{l|}{{\color[HTML]{FE0000} }}      & \multicolumn{1}{l}{}      & \multicolumn{1}{l}{}      & \multicolumn{1}{l|}{{\color[HTML]{FE0000} }}      & \multicolumn{1}{l}{}     & \multicolumn{1}{l}{}     & \multicolumn{1}{l}{{\color[HTML]{FE0000} }}     \\
\multicolumn{1}{c|}{4} & Powersave                        & 0.524                     & 0.524                     & \multicolumn{1}{c|}{{\color[HTML]{FE0000} 0.598}} & 1.729                     & 1.731                     & \multicolumn{1}{c|}{{\color[HTML]{FE0000} 1.974}} & 3.04                     & 3.03                     & {\color[HTML]{FE0000} 2.66}                     \\
\multicolumn{1}{c|}{}  & Ondemand                         & 1.016                     & 0.999                     & \multicolumn{1}{c|}{{\color[HTML]{FE0000} 1.082}} & 3.353                     & 3.298                     & \multicolumn{1}{c|}{{\color[HTML]{FE0000} 3.569}} & 1.57                     & 1.59                     & {\color[HTML]{FE0000} 1.47}                     \\
\multicolumn{1}{c|}{}  & Performance                      & 0.986                     & 0.978                     & \multicolumn{1}{c|}{{\color[HTML]{FE0000} 1.042}} & 3.254                     & 3.229                     & \multicolumn{1}{c|}{{\color[HTML]{FE0000} 3.437}} & 1.61                     & 1.63                     & {\color[HTML]{FE0000} 1.53}                     \\
\multicolumn{1}{c|}{}  & Conservative                     & 0.979                     & 0.981                     & \multicolumn{1}{c|}{{\color[HTML]{FE0000} 1.055}} & 3.231                     & 3.238                     & \multicolumn{1}{c|}{{\color[HTML]{FE0000} 3.483}} & 1.63                     & 1.62                     & {\color[HTML]{FE0000} 1.51}                     \\ \hline
\end{tabular}
\caption{1800 mAh battery run-time estimation under high load stress, with a dimensional abuse C/5}
\label{runtime2}
\end{table}
\vspace{-0.8cm}

When the \acrshort{CPU} is dealing with a high load stress, the battery run-time falls drastically, as expected. However, it reaches an hour and a half, which is a great value taking into account that this workload will not be usual.
% Capítulo 5: Conclusiones y trabajo futuro
\chapter{Validation}\label{chap:chapter5}


This fifth chapter is used to validate the solutions proposed in this Master's Thesis and check the accomplishment of the goals of the project. This verification is performed using two resources: on the one hand, different photos of the developed products, which complement the ones introduced in each section of the \autoref{chap:chapter4}; on the other hand, the \glsname{ground} software, COSMOS, is used to verify the functioning of the different subsystems, but simulating several operations using the \acrshort{I2DOS} and the GranaSat-I \glsname{cubesat} prototype.


\section{Final product} 
%
%\begin{wrapfigure}{r}{0.48\textwidth} 
	%\centering
		%\includegraphics[angle=270,page=1,trim={0cm 0cm 0cm 0cm},clip=true,width=60mm]{FigurasTFM/Chapter5/Imagenes/final_pr1.jpg}	%
			%\caption{\acrshort{I2DOS} and GranaSat-I prototype} 	\label{final1}
			%\vspace{1cm}
%\end{wrapfigure}

In this section, different photos of the design simulation platform are shown. For example, \autoref{final1} shows the \acrshort{I2DOS} with the developed 1U \glsname{cubesat}, while using the \acrshort{LED} simulator.

The \glsname{cubesat} is powered with an external power bank, which complies with the the \acrshort{EPS} requirement \textbf{EPS.FoR.4}. The solar panel test board has been kept in order to check the behaviour of the platform when the \textbf{mass distribution} is different to the expected; it is able to deal with the imbalance correctly. \autoref{final2} shows a detail of the \glsname{cubesat} electronics.

	\begin{figure} [H]
			\centering
			\includegraphics[angle=270,page=1,trim={0cm 0cm 0cm 0cm},clip=true,width=70mm]{FigurasTFM/Chapter5/Imagenes/final_pr1.jpg}
			\caption{\acrshort{I2DOS} and GranaSat-I prototype} 	\label{final1}
			\vspace{-0.5cm}
\end{figure}


		\begin{figure} [H]
			\centering
			\includegraphics[angle=270,page=1,trim={0cm 0cm 0cm 0cm},clip=true,width=90mm]{FigurasTFM/Chapter5/Imagenes/finalpr2.jpg}
			\caption{GranaSat-I \glsname{cubesat} detail} 	\label{final2}
			\vspace{-2cm}
\end{figure}

Video \autoref{videofinal} shows a slow motion clip (Adobe Reader or another compatible PDF reader needed) of the platform moving.

\vspace{-0.3cm}

\begin{videoFloat}[H]
\centering
\includemovie[text={\includegraphics[page=1,trim={0cm 0cm 0cm 0cm},clip=true,width=115mm]{figurastfm/Chapter5/Imagenes/finalvideo.png}}]{10cm}{7cm}{FigurasTFM/Chapter5/Video/videofinal.wmv}
%\vspace{0.5cm}
\caption{\acrshort{I2DOS} platform in movement (Adobe Reader needed)} \label{videofinal}
\end{videoFloat}
\vspace{-0.5cm}

%\newpage




\section{OBC} 

As for \acrshort{OBC}, once verified electrical connectivity and supply this section validates the functioning of the sensors performing different test and checking the results with COSMOS.

\vspace{-0.3cm}

\subsection{Sensors} 

As validation for the \acrshort{OBC} sensors, two screenshots of the \glsname{ground} software are included. \autoref{obcfig1} illustrates the state of the sensors in a static position, as inferable from the static illumination received (see ADC\_Measurements plot); the previous oscillations imply facing the \textit{Sun} repeatedly when rotating. Also, pressure is stable so there is no variation in height and it confirms that the \glsname{cubesat} \textbf{is stopped}. The plot \textbf{OBC\_Main\_Temperatures} shows the temperatures measured by two different sensors: the Bosch BMP280 and the internal sensor of the \acrshort{CPU}. It helps to see the \textbf{gradient} of temperature across the \acrshort{OBC}.

On the other hand, \autoref{obcfig2} shows the result sent by the sensors when the \glsname{cubesat} has gone through a complete \textbf{start-and-stop} cycle, i.e., it has started tumbling and the \acrshort{ADCS} has counteracted. Therefore, there is appreciable variations in pressure (height) and in illumination (the Sun sensors have faced it multiple times and  finally has left in eclipse position). Also, because of the higher workload and the operation of the actuators the overall temperature of the system have increased (especially in the \acrshort{CPU}, with a gradient of 10 \textdegree C).

\subsection{Communications} 

The whole telemetry data for this validation testbench is being sent over the air, using the wireless communication system of the \glsname{cubesat}. The \acrshort{OBC} sends several parameters regarding this subsystem, and they are plotted in one of the section designed in COSMOS, see \autoref{commscosmo}. Particularly, this data corresponds to the same start-and-stop cycle mentioned before; it is specially interesting the \textbf{signal level}, which varies according to the position of the antenna with respect to the receiver.


Although this validation is being performed using wireless communications, the \textbf{Ethernet} interface has been also tested, with successful results. In fact, the wired connection is necessary to configure the system the first time. In conjunction with the \glsname{ground} management software, COSMOS, the communications subsystem is able to interact in a swarm of CubeSats, which allows the use of this platform in environments such as classrooms or training sessions. \autoref{figcosoms} depicts this possible utilization.

\vspace{0.5cm}

	\begin{figure} [H]
			\centering
			\includegraphics[page=1,trim={0cm 0cm 0cm 0cm},clip=true,width=165mm]{FigurasTFM/Chapter5/PDF/COMSCOSMOS.pdf}
			\caption{Communication subsystem allows being used in a \acrshort{LAN} with multiple CubeSats} 	\label{figcosoms}
			%\vspace{-0.5cm}
\end{figure}


\begin{landscape}
\vspace*{\fill}
	\begin{figure} [H]
			\centering
			\includegraphics[page=1,trim={0cm 0cm 0cm 0cm},clip=true,width=230mm]{FigurasTFM/Chapter5/Imagenes/OBC_1.png}
			\caption{\acrshort{OBC} measurements with static \glsname{cubesat}} 	\label{obcfig1}
			\vspace{-0.5cm}
\end{figure}
\vspace*{\fill}
\end{landscape}

\begin{landscape}
\vspace*{\fill}
	\begin{figure} [H]
			\centering
			\includegraphics[page=1,trim={0cm 0cm 0cm 0cm},clip=true,width=230mm]{FigurasTFM/Chapter5/Imagenes/OBC_2.png}
			\caption{\acrshort{OBC} measurements with static \glsname{cubesat}} 	\label{obcfig2}
			\vspace{-0.5cm}
\end{figure}
\vspace*{\fill}
\end{landscape}


\begin{landscape}
\vspace*{\fill}
	\begin{figure} [H]
			\centering
			\includegraphics[page=1,trim={0cm 0cm 0cm 0cm},clip=true,width=230mm]{FigurasTFM/Chapter5/Imagenes/comms.png}
			\caption{COMMS subsystem telemetry during rotation} 	\label{commscosmo}
			\vspace{-0.5cm}
\end{figure}
\vspace*{\fill}
\end{landscape}


\section{ADCS}

This section addresses the validation of the \acrshort{ADCS} subsystem. Firstly, the magnetometers are characterized using the \textbf{Helmholtz Cage} available at \glsname{GranaSAT} laboratory. Then, just as with the \acrshort{OBC}, the correct functioning of the attitude-related sensors is verified with COSMOS.

\subsection{Magnetometer characterization} 

\autoref{magneto1} shows the testbench used to perform this characterization, with the GranaSat-I in the center of the cage.

	\begin{figure} [H]
			\centering
			\includegraphics[page=1,trim={0cm 0cm 0cm 0cm},clip=true,width=80mm]{figurastfm/Chapter5/Imagenes/testb2.jpg}
			\caption{Magnetometer characterization testbench} \label{magneto1}
			%\vspace{-0.q5cm}
\end{figure}

It is performed a simple characterization test, to assure the correct functioning of the sensor. In order to get the \glsname{cubesat} ready for a real mission, it is necessary to go through the procedure described in section \ref{magnetotest}, standardized in \cite{ESAmagnetic}. \autoref{magnetictest} shows the results; because timing is not important in this test, X-axis is shown in samples.


\begin{landscape}
\vspace*{\fill}
	\begin{figure} [H]
			\centering
			\includegraphics[page=1,trim={0cm 0cm 0cm 0cm},clip=true,width=230mm]{FigurasTFM/Chapter5/PDF/magnetic.pdf}
			\caption{Magnetic field measurement test} 	\label{magnetictest}
			%\vspace{-0.q5cm}
\end{figure}
\vspace*{\fill}
\end{landscape}
\newpage

The procedure to characterize the magnetometer is composed of 4 cycles, clearly differentiated in \autoref{magnetictest} with colours, described next.

\begin{itemize} [topsep=0pt]

\item The first cycle, in \textcolor{orange}{orange}, is a measurement of the \textbf{Earth's magnetic field}, i.e., the Helmholtz Cage is not generating any field, but measuring the field of the Earth. The coincidence between both measurements is already a good sign in itself.

\item The second cycle, in \textcolor{purple}{purple}, depicts the \textbf{cancellation of the Earth's magnetic field}. Once again, the measurements taken by MPU-9250 magnetometers are pretty accurate, and it is able to detect a null field in the three axes.

\item The third cycle, in \textcolor{lime}{green}, shows the measurements of a certain magnetic field generated by the cage. Particularly, with a magnitude of \SI{35}{\micro T}, which is the module of the components measured [$x_{\text{field}}$, $y_{\text{field}}$, $z_{\text{field}}$] = [\SI{-10}{\micro T}, \SI{27}{\micro T}, \SI{-18}{\micro T}].

\item The fourth cycle, in \textcolor{red}{red}, measures Earth's magnetic field again.

\item The test of the magnetometers can clearly be considered a success, taking into account the high degree of accuracy exhibited. However, it is worth noting the amount of \textbf{noise} present in the measurements, specially in the Z-axis. This spatial component is particularly affected by this phenomena probably because it is in the same plane of the DC-motor and the rest of the electronic components, which produce a non-negligible amount of noise in its plane. It would be recommendable to implement some \textbf{filtering} to improve the results.

\end{itemize}

\vspace{-0.3cm}
\subsection{Sensors} 

Once again, to validate the functioning of the sensors, they are included the results received and plotted in COSMOS during operation, with special interest in the maneuver of \textbf{detumbling}; \autoref{adcs1} shows this operation artificially repeated. It can be seen that the system has to deal with different rotational uncontrolled movements, in different senses of rotation; it is easy to see this situation in the \textit{\textbf{Angular\_Speed}} plot. In all cases, it is able to counteract that movement and get the system \textbf{to a stable position}. 

Regarding the magnetic field, it also varies in the Z-axis (plane of \textbf{rotation}) and \linebreak X-axis (because the platform is not completely \textbf{perpendicular} to the Z-axis) but remains approximately \textbf{constant in the Y-axis}, as the \glsname{cubesat} is at a constant \textbf{elevation}. Finally, the \textbf{Accelerometer} measurements, as expected are approximately null in X and Y axis, and \textbf{1 g in the Z-axis} (1 g = \SI{9.8}{m/s^2}).

On the other hand, \autoref{adcs2} shows a maneuver inverting its sense of rotation. By looking at the \textit{Angular\_Speed} plot, the \glsname{cubesat} is rotating counterclockwise, and the \acrshort{ADCS} instead of stabilizing it, inverts its rotation; although this is not a realistic maneuver, it is useful to check the capabilities and functioning of the control system. The inversion is also easy to see in the \textit{Magnetic\_Field} plot.

\begin{landscape}
\vspace*{\fill}
	\begin{figure} [H]
			\centering
			\includegraphics[page=1,trim={0cm 0cm 0cm 0cm},clip=true,width=230mm]{FigurasTFM/Chapter5/Imagenes/adcs1.png}
			\caption{Detumbling operation} 	\label{adcs1}
			%\vspace{-0.q5cm}
\end{figure}
\vspace*{\fill}
\end{landscape}
\newpage

\begin{landscape}
\vspace*{\fill}
	\begin{figure} [H]
			\centering
			\includegraphics[page=1,trim={0cm 0cm 0cm 0cm},clip=true,width=230mm]{FigurasTFM/Chapter5/Imagenes/adcs2.png}
			\caption{Rotation sense inversion} 	\label{adcs2}
			%\vspace{-0.q5cm}
\end{figure}
\vspace*{\fill}
\end{landscape}
\newpage



% Capitulo 6: ------
\chapter{Conclusions and Future Lines}\label{chap:chapter6}

This document has shown the tough process under space-related projects by developing an innovative and functional \glsname{cubesat} Simulation Platform. This Master's~Thesis has been undertaken following several standards intensively used in space industry, using professional equipment and applying complex design techniques which go beyond the theory typically learnt during this tuition period.

This project provides a complete integrated environment in relation with CubeSats. As stated all along this document, two perspectives have been addressed: \textbf{mission-oriented} and \textbf{academic approach}. Particularly, the main contributions of this Master's Thesis are:

\begin{itemize} [topsep=0pt]

\item The \acrshort{I2DOS} Platform, which allows simulating a vacuum environment by using a low-friction rotating table, as well as exposing the \glsname{cubesat} to external lightning (simulating the Sun) or magnetic fields (simulating Earth's field). 

\item An expandable \glsname{ground} software, based on the Ball Aerospace COSMOS framework, previously used by \acrshort{NASA} or \textbf{Lockheed~Martin}; therefore, with this environment, the \glsname{GranaSAT} laboratory counts now with a professional mission management tool which can be used not only on the basis of a simulation, but also in a real mission.

\item A simulation \glsname{cubesat} prototype, which includes an aluminum mechanical structure, and counts with the main subsystems such as \acrshort{OBC}, \acrshort{ADCS}, local~COMMS as well as a basic \acrshort{OBDH}. It is ready to send telemetry information to COSMOS and allows defining new telemetry packages or telecommands. The \acrshort{ADCS} includes a \acrshort{PID}-based control law, which can be configured to detumble the device or follow a certain target. Additionally, this Master's Thesis includes an intensive characterization of different solar panels and batteries to be used with a \glsname{cubesat}, which will ease design decisions in the future \acrshort{EPS}.

\end{itemize}

Following the System Engineering methodology proposed by Wertz and Everett \cite{smad} has allowed minimizing risks and costs in the project and furthermore, accomplishing the goals defined.

This variety has supposed an increasing complexity which has given realism to this project, but also demanded an incredible effort. As for me, it has supposed an unprecedented challenge not only because of the inherent difficulty associated to this kind of project, but also because of its \textbf{extraordinarily wide scope}. Indeed, this project has made me learn things from a variety of fields of the engineering, which undoubtedly will be priceless in my professional future.

In a project of these characteristics, the end is never reached. Therefore, during the development of this Master's Thesis, some improvements and or future lines of work have naturally arisen, featuring the following ones:


\begin{itemize} [topsep=0pt]

\item Diminishing the \textbf{moment of inertia} of the inertial platform by reducing its diameter. This would ease the rotation control and detumbling of the \glsname{cubesat} and prevent the motor of the reaction wheel from saturation.

\item Defining \textbf{new telemetry packages}, including all the information provided by the sensors, and implementing automatic \textit{keep alives} or \glsname{housekeeping}.

\item Developing \textbf{new COSMOS interfaces}, which cover all the aspects of a mission and take full advantage of the capabilities of the GranaSAT-I, defining new telecommands and reports.

\item Keeping with the development of an adequate \acrshort{EPS}, based on the analysis and characterizations addressed in this project, which provides with power lines of 5 V and 3.3 V.

\item Implementing printed \glsname{magnetorquers}, which increases the \textbf{control capabilities} of the \glsname{cubesat}. The designed \acrshort{ADCS} already takes this possibility into account and includes terminals for three \glsname{magnetorquers}.

\item Defining new \acrshort{ADCS}-related functions for the co-processing programmable cores (\acrshort{FPGA}) which frees the \acrshort{OBC} of tasks and enables a \textbf{time-critical} management when needed.

\item Generating documentation regarding the academic orientation of the platform, with examples, case of use and guided procedures.

\end{itemize}

All in all, this Master's Thesis supposes a finalizing touch to my academic period but it does not end here; it is a live project which is now another part of \glsname{GranaSAT}, and I hope it is useful and brings success to its members henceforth.

Once again, the future is exciting.



%Begin ----  Para que funcione bien el TOC en PDF
%\cleardoublepage
%\begin{appendices}
%\clearpage % o \cleardoublepage
%\chapter{Presupuesto}
%\label{cap:budget}
%asdf

%\end{appendices}

\begin{appendices}
\chapter{Some Appendix}
The contents...
\end{appendices}
% Bibliografia.
%\clearpage
%{
%    \pagestyle{fancy}
%    \cleardoublepage
%}
%
%\fancypagestyle{biblio}
%{
%    \fancyhf{}
%    \fancyfoot[C]{Limitador de sonido para locales de música}
%    \fancyfoot[LE,RO]{\thepage}
%    \fancyhead[LE,RO]{\small{Bibliografía}}
%    \renewcommand{\headrulewidth}{0.5pt}
%    \renewcommand{\footrulewidth}{0.5pt}
%}

% ------------------------------------------------

% Incluir en el ToC
\addcontentsline{toc}{chapter}{Bibliografía}
% Definir estilo de la página
\pagestyle{biblio}

% Tipos de Bibliografía: https://www.overleaf.com/learn/latex/Bibtex_bibliography_styles
\bibliographystyle{abbrv}
% Fichoer con las referencias
\bibliography{./bibliografia/bibliografia}

% ------------------------------------------------

% Para que funcione bien el TOC en PDF
\cleardoublepage
\phantomsection
\label{bib}
% Apendice
\begin{appendices}


\chapter{Project Budget} \label{cap:budget}


\section{Materials and hardware}

In this section, project cost regarding materials and hardware sections will be detailed. Each one of the different hardware subsections is differentiated. Human resources area also included.

\subsection{Inertial 2D Orbit Simulator}

According to the System Design, \autoref{chap:chapter4}, the materials related with the \acrshort{I2DOS} are mainly plywood for the base and \acrshort{PLA} for the 3D printed inertial platform. The plywood board has a cost of \textbf{11.30 €}, while the amount necessary of \acrshort{PLA} costs \textbf{4.50 €}. That makes \textbf{15.80~€} without manufacturing costs, detailed in \autoref{humcosts}.


\subsection{Simulation CubeSat}

\subsubsection{Mechanical structure}

The mechanical structure of the \glsname{cubesat} is made of aluminum EN AW 5754, also sold in boards. Because of the extremely light weight of a \glsname{cubesat}, the attributable cost for the manufacturing of the GranaSAT-I is only \textbf{2.50 €}.

\subsubsection{On-board Computer}

The cost of the \acrshort{OBC} is breakdown into the electronic components and the manufacturing, as follows.

\paragraph{Bill of Materials}

% Please add the following required packages to your document preamble:
% \usepackage{booktabs}
\begin{table}[H]
\centering
\begin{tabular}{@{}cccc@{}}
\toprule
\textbf{Item}            & \textbf{Units}       & \textbf{Cost/u. (€)} & \textbf{Cost (€)}    \\ \midrule
LDR                      & 6                    & 0.15                 & 0.9                  \\
Bosch BMP280             & 1                    & 2.68                 & 2.68                 \\
MPU-9250                 & 1                    & 5.1                  & 5.1                  \\
CR2032 Battery           & 1                    & 0.35                 & 0.35                 \\
Electrolitic cap.        & 4                    & 0.1                  & 0.4                  \\
Tantalum cap             & 2                    & 0.8                  & 1.6                  \\
Ceramic cap.             & 36                   & 0.012                & 0.432                \\
LED SMD                  & 10                   & 0.05                 & 0.5                  \\
1N4148W                  & 1                    & 0.05                 & 0.05                 \\
ENC28J60 breakout        & 1                    & 3                    & 3                    \\
ESP8266                  & 1                    & 2                    & 2                    \\
Micro USB 2.0 B          & 3                    & 3.5                  & 10.5                 \\
J00-0045L                & 1                    & 4.5                  & 4.5                  \\
ESQ-126-39-G-D           & 2                    & 7                    & 14                   \\
Canned Oscillator 25 MHz & 1                    & 3.5                  & 3.5                  \\
Transistors 2N7002P      & 5                    & 0.3                  & 1.5                  \\
Resistors                & 51                   & 0.01                 & 0.51                 \\
Raspberry Pi 3 CM        & 1                    & 32                   & 32                   \\
PCF8523                  & 1                    & 1.13                 & 1.13                 \\
ADS1015                  & 2                    & 2.84                 & 5.68                 \\
LAN9514                  & 1                    & 5.13                 & 5.13                 \\
FSUSB42                  & 1                    & 0.5                  & 0.5                  \\
\multicolumn{1}{l}{}     & \multicolumn{1}{l}{} & \multicolumn{1}{l}{} & \multicolumn{1}{l}{} \\
\multicolumn{3}{c}{\textbf{Subtotal (before VAT):}}                    & \textbf{79.31}       \\
\multicolumn{3}{c}{\textbf{Total (VAT included):}}                     & \textbf{95.96}       \\ \bottomrule
\end{tabular}
\caption{\acrshort{OBC} \acrshort{PCB} components}
\label{obcomp}
\end{table}


\paragraph{Manufacturing}

The manufacturing of a 4-layers \acrshort{PCB} like the \acrshort{OBC}, including \glsname{stencil} has a cost of \textbf{43.40 €}.

\subsubsection{Attitude Determination and Control System}
\paragraph{Bill of Materials}
% Please add the following required packages to your document preamble:
% \usepackage{booktabs}
\begin{table}[H]
\centering
\begin{tabular}{@{}cccc@{}}
\toprule
\textbf{Item}            & \textbf{Units}       & \textbf{Cost/u. (€)} & \textbf{Cost (€)}    \\ \midrule
BD6211F                  & 2                    & 1.9                  & 3.8                  \\
L298N                    & 1                    & 5.2                  & 5.2                  \\
DC motor                 & 1                    & 7                    & 7                    \\
ICE40UP5K                & 1                    & 7.1                  & 7.1                  \\
1N4148W                  & 4                    & 0.05                 & 0.2                  \\
Ceramic cap.             & 16                   & 0.012                & 0.192                \\
CNY70                    & 1                    & 1.8                  & 1.8                  \\
ESQ-126-39-G-D           & 2                    & 7                    & 14                   \\
Canned Oscillator 25 MHz & 1                    & 3.5                  & 3.5                  \\
Transistors 2N7002P      & 8                    & 0.3                  & 2.4                  \\
Resistors                & 25                   & 0.01                 & 0.25                 \\
TLV7111225               & 1                    & 0.5                  & 0.5                  \\
\multicolumn{1}{l}{}     & \multicolumn{1}{l}{} & \multicolumn{1}{l}{} & \multicolumn{1}{l}{} \\
\multicolumn{3}{c}{\textbf{Subtotal (before VAT):}}                    & \textbf{37.97}       \\
\multicolumn{3}{c}{\textbf{Total (VAT included):}}                     & \textbf{45.94}       \\ \bottomrule
\end{tabular}
\caption{\acrshort{ADCS} \acrshort{PCB} components}
\label{adcscomp}
\end{table}

\paragraph{Manufacturing}

The manufacturing of a 2-layers \acrshort{PCB} like the \acrshort{ADCS}, including \glsname{stencil} has a cost of \textbf{18.15 €}.

\subsubsection{Electrical Power system}
\paragraph{Bill of Materials}
% Please add the following required packages to your document preamble:
% \usepackage{booktabs}
\begin{table}[H]
\centering
\begin{tabular}{@{}cccc@{}}
\toprule
\textbf{Item}                & \textbf{Units}       & \textbf{Cost/u. (€)} & \textbf{Cost (€)}    \\ \midrule
Li-ion battery (1100 mAh)    & 1                    & 8.60                 & 8.60                 \\
Li-ion battery (1800 mAh)    & 1                    & 11.11                & 11.11                \\
Solar panel 1 (5 V @ 1.3 W)  & 1                    & 3.44                 & 3.44                 \\
Solar panel 2 (12 V @ 1.5 W) & 1                    & 2.83                 & 2.83                 \\
Solar panel 3 (5.5 V @ 1 W)  & 1                    & 2.48                 & 2.48                 \\
\multicolumn{1}{l}{}         & \multicolumn{1}{l}{} & \multicolumn{1}{l}{} & \multicolumn{1}{l}{} \\
\multicolumn{3}{c}{\textbf{Subtotal (before VAT):}}                        & \textbf{23.52}       \\
\multicolumn{3}{c}{\textbf{Total (VAT included):}}                         & \textbf{28.46}       \\ \bottomrule
\end{tabular}
\caption{\acrshort{EPS} components}
\label{epscomp}
\end{table}



%
%\section{Software Cost}
%
%This Project has required the use of different software. The price of the different suites used is detailed in table \ref{softw}
%
%% Please add the following required packages to your document preamble:
%% \usepackage{booktabs}
%\begin{table}[H]
%\centering
%\begin{tabular}{@{}ccc@{}}
%\toprule
%\textbf{Software}   & \textbf{License Owner} & \textbf{Cost (€)} \\ \midrule
%Altium Designer 19  & UGR                    & Free              \\
%Solidworks Student Edition 2017     & Author                    & Free              \\
%MatLab          & UGR                    & Free              \\
%TeXnicCenter        & Author                 & Free              \\
%Miktex              & Author                 & Free              \\
%\textbf{}           & \textbf{TOTAL}         & \textbf{0 €}      \\ \bottomrule
%\end{tabular}
%\caption{Software Cost}
%\label{softw}
%\end{table}

\section{Human Resources Cost} \label{humcosts}

In the first place, the manufacturing of the base of the \acrshort{I2DOS} required the collaboration of the \textbf{makerspace} at the Architecture School of the University of Granada. The cost for an hour of machinery and technical assistance was \textbf{15 €}.

Additionally, the development of this Master's Thesis has required hiring two people. The first one is a \textbf{junior engineer} (10 €/h), hired as a full-time worker during twelve months. Secondly, as Project Supervisor a \textbf{senior engineer} is hired (50 €/h), computing 5 hours per week. Then, Human Resources amounts to \textbf{31200 €}, as detailed in table \ref{wages}.

% Please add the following required packages to your document preamble:
% \usepackage{booktabs}
% \usepackage[normalem]{ulem}
% \useunder{\uline}{\ul}{}
\begin{table}[H]
\centering

\begin{tabular}{@{}ccc@{}}
\toprule
{\ul \textbf{Post}} & {\ul \textbf{Time (Hours)}} & {\ul \textbf{Cost (€)}} \\ \midrule
Junior Engineer     & 1920                        & 19200                   \\
Senior Engineer     & 240                         & 12000                   \\
                    & \textbf{TOTAL}              & \textbf{31200 €}
\end{tabular}
\caption{Human Resources Cost}
\label{wages}
\end{table}

\end{appendices}
%\end{document}




%
%\chapter*{Acronyms}%
%\fancyhead{}
%\fancyhead[RO]{\small{\nouppercase{Acronyms}}}
%\fancyhead[LE]{\small{\nouppercase{Acronyms}}}%
%\printglossary[type=acronym,title=Acrónimos y siglas,toctitle=Acrónimos y siglas]
\printglossary[style=long3col,type=\acronymtype,title=Acrónimos,toctitle=Acrónimos]
\clearpage{\pagestyle{plain}\cleardoublepage}%



%Glosario
\clearpage{\pagestyle{fancy}\cleardoublepage}%
\fancyhead{}
\fancyhead[RO]{\small{\nouppercase{Glosario}}}
\fancyhead[LE]{\small{\nouppercase{Glosario}}}%
%Begin ----  Para que funcione bien el TOC en PDF
\cleardoublepage
\phantomsection \label{Glosar}
%END  ---- Para que funcione bien el TOC en PDF

\printglossary[type=main,title=Glosario,toctitle=Glosario]


%Listado de Símbolos
\clearpage{\pagestyle{fancy}\cleardoublepage}%
\fancyhead{}
\fancyhead[RO]{\small{\nouppercase{Symbols List}}}
\fancyhead[LE]{\small{\nouppercase{Symbols List}}}%
%Begin ----  Para que funcione bien el TOC en PDF
\cleardoublepage
\phantomsection \label{SymbolsList}
%END  ---- Para que funcione bien el TOC en PDF

\printglossary[style=long3col,type=symbols,title=Symbols List,toctitle=Symbols List]
%\addcontentsline{toc}{chapter}{Símbolos}

%Índice Alfabético de apariciones
\clearpage{\pagestyle{fancy}\cleardoublepage}
\fancyhead{}
\fancyhead[RO]{\small{\nouppercase{indice de apariciones}}}
\fancyhead[LE]{\small{\nouppercase{indice de apariciones}}}%



%


\index{ASIC}
\index{ASSP}



\index{lenguajes!Verilog-A}
\index{lenguajes!FORTRAN}
\index{lenguajes!lenguaje-C}
\index{lenguajes!MATLAB}

\index{Verilog-A!analog|ii}
\index{Verilog-A!begin|ii}
\index{Verilog-A!end|ii}
\index{Verilog-A!\$vt|ii}
\index{Verilog-A!\$temperature|ii}
\index{Verilog-A!\$simparam|ii}
\index{Verilog-A!\gls{contribucion}|ii}
\index{Verilog-A!\gls{descVar}|ii}
\index{Verilog-A!\gls{ddt}|ii}
\index{Verilog-A!\gls{ddx}|ii}
\index{Verilog-A!\gls{WhiteNoise}|ii}
\index{Verilog-A!\gls{FlicekrNoise}|ii}




\index{spice}
\index{netlist}
\index{bites!animal!gnats|ii}
\index{gnat!size of}
\index{gnat|ii}
\index{gnat!anatomy}
\index{gnus!good}
\index{gnus!bad}



%Begin ----  Para que funcione bien el TOC en PDF
\cleardoublepage
\phantomsection \label{IndiceApariciones}
%END  ---- Para que funcione bien el TOC en PDF

\printindex



\end{document}